% Transformation et conservation des fruits de saison — SVTEEHB Tle C — Cours signé OnBuch+
% Programme officiel MINESEC. Compiler avec : tectonic N10-conservation-fruits-de-saison.tex
\newcommand{\DOCMATIERE}{SVTEEHB}
\newcommand{\DOCNIVEAU}{Tle C}
\newcommand{\DOCMODULE}{Module IV : La Biotechnologie}
\newcommand{\DOCLECON}{N10}
\newcommand{\DOCTITRE}{Transformation et conservation des fruits de saison}
% OnBuch+ — préambule « cours signés » ENRICHI (compilé avec tectonic / XeLaTeX)
% Système visuel « Pop » d'OnBuch+ : fond crème (jamais blanc), bordures encre
% épaisses, ombres « dures » décalées, titres Archivo Black en capitales.
% Version MATHÉMATIQUES : graphes (pgfplots/tikz), boîtes pédagogiques
% (définition, propriété, méthode, attention, le savais-tu, exercice résolu,
% exercice, TP, bilan), page de garde et signature OnBuch+ ancrée en bas.
%
% Macros attendues dans main.tex AVANT \input{preamble} :
%   \DOCMATIERE  \DOCNIVEAU  \DOCMODULE  \DOCLECON (numéro)  \DOCTITRE
\documentclass[11pt]{article}

\usepackage{fontspec}
\usepackage[french]{babel}
\usepackage[a4paper,margin=2cm,top=3.4cm,bottom=2.8cm,headheight=1.4cm,footskip=1.3cm]{geometry}
\usepackage{amsmath,amssymb,amsfonts}
\usepackage{mathtools} % \xmapsto, \coloneqq... souvent utilisés par les modèles
\usepackage{cancel} % simplifications barrées (bilans de forces)
% \abs{z} (module, valeur absolue) et \norm{u} (norme), utilisés par les modèles
\providecommand{\abs}{}\let\abs\relax\DeclarePairedDelimiter{\abs}{\lvert}{\rvert}
\providecommand{\norm}{}\let\norm\relax\DeclarePairedDelimiter{\norm}{\lVert}{\rVert}
\usepackage[table]{xcolor} % [table] : fournit aussi \rowcolors (alternance de lignes)
\usepackage{pagecolor}
\usepackage{tikz}
\usetikzlibrary{arrows.meta,positioning,calc,shapes.geometric,shapes.misc,shapes.arrows,decorations.pathmorphing,decorations.pathreplacing,decorations.markings,patterns,shadows,mindmap,trees,fit,backgrounds,matrix,intersections,angles,quotes}
\usepackage{pgfplots}
\usepackage[version=4]{mhchem} % équations nucléaires \ce{^{235}_{92}U}
\usepackage[european,siunitx]{circuitikz} % schémas électriques
\pgfplotsset{compat=1.18}
\usepgfplotslibrary{fillbetween,groupplots} % aires entre courbes (calcul intégral)
\usepackage{enumitem}
\usepackage{fancyhdr}
% [most] charge toutes les bibliothèques tcolorbox (listings, minted, raster,
% documentation...) et gonfle inutilement la mémoire TeX sur un document long
% (~300 boîtes) : on ne charge que ce qui est réellement utilisé.
\usepackage[skins,breakable]{tcolorbox}
\usepackage{graphicx}
\usepackage{colortbl}
\usepackage{booktabs}
\usepackage{tabularx}
\usepackage{multirow}
\usepackage{diagbox} % case d'en-tête diagonale des tableaux à double entrée
% Colonnes extensibles centrée / gauche / droite (C, L, R), souvent utilisées par les modèles
\newcolumntype{C}{>{\centering\arraybackslash}X}
\newcolumntype{L}{>{\raggedright\arraybackslash}X}
\newcolumntype{R}{>{\raggedleft\arraybackslash}X}
\usepackage{array}
\usepackage{multicol}
\usepackage{siunitx}
% parse-numbers=false : accepte aussi \SI{2,5 \times 10^{-3}}{...}
\sisetup{locale=FR,output-decimal-marker={,},group-minimum-digits=4,parse-numbers=false}
\DeclareSIUnit\ohm{\ensuremath{\symup{\Omega}}} % U+2126 absent de la police
\DeclareSIUnit\division{div} % réglages d'oscilloscope (ms/div, V/div)
\DeclareSIUnit\tr{tr} % tours (vitesse de rotation, ex. \SI{1500}{\tr\per\minute})
\DeclareSIUnit\molar{M} % concentration molaire (ex. \SI{3}{\molar}, \SI{1}{\micro\molar})
\DeclareSIUnit\an{an} % année (le modèle confond parfois avec \year, un compteur TeX) — ex. \SI{5}{\kilo\meter\per\an}
\DeclareSIUnit\FCFA{FCFA} % franc CFA (ex. \SI{500}{\FCFA\per\kilo\gram})
\DeclareSIUnit\degreeBrix{\textdegree Brix} % teneur en sucre d'un sirop/jus (réfractométrie)
\DeclareSIUnit\month{mois} % le modèle utilise parfois \month, un compteur TeX, comme unité de durée
\DeclareSIUnit\microliter{\micro\liter} % le modèle écrit parfois \microliter en un seul token
\DeclareSIUnit\million{millions} % dénombrement (ex. \SI{15}{\million\per\mL} spermatozoïdes)
\usepackage{qrcode}
% \degree : commande fréquemment utilisée par les modèles pour un angle ou une
% température en dehors de \SI{}{} (non fournie par les paquets chargés).
\providecommand{\degree}{^{\circ}}
% \qed (fin de démonstration), utilisé par les modèles sans amsthm
\providecommand{\qed}{\hfill$\square$}
% \barre : double barre des valeurs interdites dans un tableau de variations (tkz-tab)
\providecommand{\barre}{\ensuremath{\|}}
% environnement proof (amsthm non chargé) utilisé par les modèles
\ifdefined\proof\else\newenvironment{proof}[1][Démonstration]{\par\noindent\textit{#1.}\ }{\qed\par}\fi
\usepackage{lastpage}
\usepackage{hyperref}
\hypersetup{hidelinks}

% ============================================================
% Palette « Pop » OnBuch+ — mêmes hex que la classe P (plus_ui.dart)
% ============================================================
\definecolor{poporange}{HTML}{F59321}
\definecolor{popdark}{HTML}{E07A0C}
\definecolor{popcream}{HTML}{FFF6E8}
\definecolor{popink}{HTML}{1C1714}
\definecolor{poppurple}{HTML}{7A5AE0}
\definecolor{popgreen}{HTML}{2E9E6A}
\definecolor{poppink}{HTML}{E0457B}
\definecolor{popblue}{HTML}{2F7DE1}
\definecolor{popgold}{HTML}{C98A12}
\definecolor{popmuted}{HTML}{8A7F76}
\definecolor{popline}{HTML}{EADFD2}
\definecolor{popgray}{HTML}{8A7F76}
\colorlet{popgrey}{popgray}
% Alias tolérants (le modèle nomme parfois les couleurs différemment)
\colorlet{popwhite}{white}
\colorlet{popblack}{popink}
\colorlet{popred}{poppink}
\colorlet{popyellow}{popgold}
% Teintes claires (fonds de tableaux/encadrés) — le modèle nomme parfois
% les couleurs avec un suffixe L (ex. popblueL) sans les avoir définies.
\colorlet{poporangeL}{poporange!20}
\colorlet{popdarkL}{popdark!20}
\colorlet{poppurpleL}{poppurple!20}
\colorlet{popgreenL}{popgreen!20}
\colorlet{poppinkL}{poppink!20}
\colorlet{popblueL}{popblue!20}
\colorlet{popgoldL}{popgold!20}
\colorlet{popredL}{popred!20}
\colorlet{popgrayL}{popgray!20}
% Teintes claires pour fonds de tableaux / zones
\colorlet{popblueL}{popblue!10!white}
\colorlet{poporangeL}{poporange!14!white}
\colorlet{popgreenL}{popgreen!12!white}
\colorlet{poppurpleL}{poppurple!12!white}
\colorlet{poppinkL}{poppink!12!white}
\colorlet{popgoldL}{popgold!14!white}

\pagecolor{popcream}

% ============================================================
% Typographie
% ============================================================
\defaultfontfeatures{Ligatures=TeX}
\setmainfont{PlusJakartaSans-Medium.ttf}[Path=../../../fonts/,BoldFont=PlusJakartaSans-Bold.ttf]
% Maths en sans-serif (Fira Math) avec chiffres et lettres droites de Plus
% Jakarta, pour que les formules se fondent dans le texte.
\usepackage[math-style=ISO,bold-style=ISO]{unicode-math}
\setmathfont{FiraMath-Regular.otf}[Path=../../../fonts/]
\setmathfont{PlusJakartaSans-Medium.ttf}[Path=../../../fonts/,range={up/num,up/{Latin,latin},\mathup}]
\usepackage{icomma} % virgule décimale sans espace en mode math
% Caractères mathématiques Unicode absents des polices de texte (Plus Jakarta,
% Archivo Black) : rendus par leur équivalent mathématique, en texte comme en
% formule — sinon ils s'affichent en carré vide (ex. « Arithmétique dans ℤ »).
\usepackage{newunicodechar}
\newunicodechar{ℝ}{\ensuremath{\mathbb{R}}}
\newunicodechar{ℤ}{\ensuremath{\mathbb{Z}}}
\newunicodechar{ℂ}{\ensuremath{\mathbb{C}}}
\newunicodechar{ℕ}{\ensuremath{\mathbb{N}}}
\newunicodechar{ℚ}{\ensuremath{\mathbb{Q}}}
\newunicodechar{𝔻}{\ensuremath{\mathbb{D}}}
\newunicodechar{α}{\ensuremath{\alpha}}
\newunicodechar{β}{\ensuremath{\beta}}
\newunicodechar{γ}{\ensuremath{\gamma}}
\newunicodechar{δ}{\ensuremath{\delta}}
\newunicodechar{ε}{\ensuremath{\varepsilon}}
\newunicodechar{θ}{\ensuremath{\theta}}
\newunicodechar{λ}{\ensuremath{\lambda}}
\newunicodechar{μ}{\ensuremath{\mu}}
\newunicodechar{σ}{\ensuremath{\sigma}}
\newunicodechar{τ}{\ensuremath{\tau}}
\newunicodechar{φ}{\ensuremath{\varphi}}
\newunicodechar{ψ}{\ensuremath{\psi}}
\newunicodechar{ω}{\ensuremath{\omega}}
\newunicodechar{Σ}{\ensuremath{\Sigma}}
% majuscules produites par \MakeUppercase dans les titres (α-aminés -> Α)
\newunicodechar{Α}{\ensuremath{\alpha}}
\newunicodechar{Β}{\ensuremath{\beta}}
\newunicodechar{Γ}{\ensuremath{\Gamma}}
\newunicodechar{Δ}{\ensuremath{\Delta}}
\newunicodechar{Ε}{\ensuremath{\varepsilon}}
\newunicodechar{Θ}{\ensuremath{\Theta}}
\newunicodechar{Λ}{\ensuremath{\Lambda}}
\newunicodechar{Μ}{\ensuremath{\mu}}
\newunicodechar{Τ}{\ensuremath{\tau}}
\newunicodechar{Φ}{\ensuremath{\Phi}}
\newunicodechar{Ψ}{\ensuremath{\Psi}}
\newunicodechar{Ω}{\ensuremath{\Omega}}
\newunicodechar{↦}{\ensuremath{\mapsto}}
\newunicodechar{∈}{\ensuremath{\in}}
\newunicodechar{∉}{\ensuremath{\notin}}
\newunicodechar{∧}{\ensuremath{\wedge}}
\newunicodechar{⇔}{\ensuremath{\Leftrightarrow}}
\newunicodechar{⟺}{\ensuremath{\Longleftrightarrow}}
\newunicodechar{⇒}{\ensuremath{\Rightarrow}}
\newunicodechar{⟹}{\ensuremath{\Longrightarrow}}
\newunicodechar{∘}{\ensuremath{\circ}}
\newunicodechar{∪}{\ensuremath{\cup}}
\newunicodechar{∩}{\ensuremath{\cap}}
\newunicodechar{≡}{\ensuremath{\equiv}}
\newunicodechar{⊂}{\ensuremath{\subset}}
\newunicodechar{⊥}{\ensuremath{\perp}}
\newunicodechar{∃}{\ensuremath{\exists}}
\newunicodechar{∀}{\ensuremath{\forall}}
\newunicodechar{∛}{\ensuremath{\sqrt[3]{\,}}}
\newunicodechar{∜}{\ensuremath{\sqrt[4]{\,}}}
\newunicodechar{⃗}{\ensuremath{{}^{\to}}}
\newunicodechar{ⁿ}{\ifmmode^{n}\else\textsuperscript{\ensuremath{n}}\fi}
\newunicodechar{ˣ}{\ifmmode^{x}\else\textsuperscript{\ensuremath{x}}\fi}
\newunicodechar{ᵇ}{\ifmmode^{b}\else\textsuperscript{\ensuremath{b}}\fi}
\newunicodechar{ʳ}{\ifmmode^{r}\else\textsuperscript{\ensuremath{r}}\fi}
\newunicodechar{ᵘ}{\ifmmode^{u}\else\textsuperscript{\ensuremath{u}}\fi}
\newunicodechar{ʸ}{\ifmmode^{y}\else\textsuperscript{\ensuremath{y}}\fi}
\newunicodechar{ᵃ}{\ifmmode^{a}\else\textsuperscript{\ensuremath{a}}\fi}
\newunicodechar{ᵉ}{\ifmmode^{e}\else\textsuperscript{\ensuremath{e}}\fi}
\newunicodechar{ᵏ}{\ifmmode^{k}\else\textsuperscript{\ensuremath{k}}\fi}
\newunicodechar{ᵖ}{\ifmmode^{p}\else\textsuperscript{\ensuremath{p}}\fi}
\newunicodechar{ᴺ}{\ifmmode^{N}\else\textsuperscript{\ensuremath{N}}\fi}
\newunicodechar{ᶜ}{\ifmmode^{c}\else\textsuperscript{\ensuremath{c}}\fi}
\newunicodechar{ʰ}{\ifmmode^{h}\else\textsuperscript{\ensuremath{h}}\fi}
\newunicodechar{ᵠ}{\ifmmode^{q}\else\textsuperscript{\ensuremath{q}}\fi}
\newunicodechar{ₙ}{\ifmmode_{n}\else\textsubscript{\ensuremath{n}}\fi}
\newunicodechar{ₐ}{\ifmmode_{a}\else\textsubscript{\ensuremath{a}}\fi}
\newunicodechar{ᵢ}{\ifmmode_{i}\else\textsubscript{\ensuremath{i}}\fi}
\newunicodechar{ᵣ}{\ifmmode_{r}\else\textsubscript{\ensuremath{r}}\fi}
\newunicodechar{ₖ}{\ifmmode_{k}\else\textsubscript{\ensuremath{k}}\fi}
\newunicodechar{ᵦ}{\ifmmode_{\beta}\else\textsubscript{\ensuremath{\beta}}\fi}
\newunicodechar{ₓ}{\ifmmode_{x}\else\textsubscript{\ensuremath{x}}\fi}
\newfontfamily\popdisplay{ArchivoBlack-Regular.ttf}[Path=../../../fonts/]
\newfontfamily\poplogo{SpaceGrotesk-Bold.ttf}[Path=../../../fonts/]
\newfontfamily\popsemi{PlusJakartaSans-SemiBold.ttf}[Path=../../../fonts/]
\newfontfamily\popxbold{PlusJakartaSans-ExtraBold.ttf}[Path=../../../fonts/]
\newfontfamily\popxboldls{PlusJakartaSans-ExtraBold.ttf}[Path=../../../fonts/,LetterSpace=3.0]

\setlist[enumerate]{leftmargin=1.7em,itemsep=5pt,topsep=4pt}
\setlist[itemize]{leftmargin=1.7em,itemsep=4pt,topsep=4pt,label={\color{poporange}\textbullet}}
\setlength{\parindent}{0pt}
\setlength{\parskip}{5pt}
\renewcommand{\arraystretch}{1.25}

% pgfplots : style Pop par défaut
\pgfplotsset{
  every axis/.append style={axis line style={popink,line width=1pt},tick style={popink},
    grid style={popline,line width=0.5pt},label style={font=\small},tick label style={font=\footnotesize}},
  popcurve/.style={poporange,line width=1.8pt,no markers,smooth},
}
% Même style disponible dans un \draw TikZ simple (hors pgfplots)
\tikzset{popcurve/.style={poporange,line width=1.8pt}}
\tikzset{popgrid/.style={popline,very thin}}
\colorlet{popcurve}{poporange} % le nom du style est parfois utilisé comme couleur

% ============================================================
% Pill / squiggle
% ============================================================
\newcommand{\poppill}[3]{%
  \tikz[baseline=-0.55ex]\node[draw=popink,line width=1.2pt,rounded corners=2.6pt,%
    fill=#2,inner xsep=6.5pt,inner ysep=3pt]{%
    \popxboldls\fontsize{7}{7}\selectfont\color{#3}\MakeUppercase{#1}};%
}
\newcommand{\popsquiggle}[1]{%
  \tikz[baseline=-0.4ex]{
    \draw[#1,line width=2.6pt,line cap=round]
      plot [smooth] coordinates {(0,0.09) (0.34,0.01) (0.68,0.21) (1.02,0.01) (1.36,0.10)};
  }%
}
% Mot-clé surligné (marqueur orange sous le texte)
\newcommand{\cle}[1]{{\popxbold\color{popdark}#1}}

% ============================================================
% En-tête / pied
% ============================================================
\pagestyle{fancy}
\fancyhf{}
\renewcommand{\headrulewidth}{0pt}
\renewcommand{\footrulewidth}{0pt}
\lhead{\raisebox{-0.15ex}{\poplogo\fontsize{13}{13}\selectfont\color{popink}OnBuch\color{poporange}+}}
\rhead{\poppill{\DOCMATIERE\ \textbullet\ \DOCNIVEAU\ \textbullet\ Leçon \DOCLECON}{white}{popink}}
\lfoot{\popxbold\fontsize{7.6}{7.6}\selectfont\color{popmuted}\MakeUppercase{Cours signé OnBuch+}\ \textbullet\ {\popsemi\DOCTITRE}}
\rfoot{\tikz[baseline=-0.6ex]\node[rounded corners=8.5pt,fill=popink,minimum height=17pt,minimum width=17pt,inner xsep=5pt,inner sep=0]{\popxbold\fontsize{8}{8}\selectfont\color{popcream}\thepage\,/\,\pageref*{LastPage}};}

% ============================================================
% Titres
% ============================================================
\newcommand{\chaptitre}[2]{%
  \begin{tcolorbox}[enhanced,colback=poporange,colframe=popink,boxrule=2.6pt,arc=11pt,
      drop shadow=popink,left=15pt,right=15pt,top=12pt,bottom=11pt,before skip=3pt,after skip=18pt]
    \poppill{#2}{popink}{popcream}\par\vspace{9pt}
    {\raggedright\hyphenpenalty=10000\exhyphenpenalty=10000\popdisplay\fontsize{22}{24}\selectfont\color{popcream}\MakeUppercase{#1}\par}
  \end{tcolorbox}
}
% Section numérotée : pastille orange avec le numéro + titre Archivo
\newcounter{popsec}
\newcommand{\coursec}[1]{%
  \stepcounter{popsec}\setcounter{popsub}{0}%
  \par\vspace{16pt}\noindent
  \begin{minipage}{\linewidth}
  \tikz[baseline=-0.7ex]\node[draw=popink,line width=1.6pt,fill=poporange,rounded corners=4pt,minimum size=22pt,inner sep=0]{\popdisplay\fontsize{12}{12}\selectfont\color{popcream}\thepopsec};%
  \hspace{8pt}{\popdisplay\fontsize{15}{17}\selectfont\color{popink}\MakeUppercase{#1}}\par
  \vspace{4pt}\noindent{\color{poporange}\rule{36pt}{3.6pt}}
  \end{minipage}\par\nopagebreak\vspace{8pt}
}
\newcounter{popsub}
\newcommand{\courssub}[1]{%
  \stepcounter{popsub}%
  \par\vspace{9pt}\noindent{\popxbold\fontsize{11.5}{13}\selectfont\color{popink}{\color{poporange}\thepopsec.\thepopsub}\hspace{6pt}#1}\par\nopagebreak\vspace{4pt}
}

% ============================================================
% Boîtes pédagogiques — même recette : fond blanc, bordure épaisse colorée,
% ombre dure encre, badge plein. Titre optionnel [#1].
% ============================================================
\tcbset{popbase/.style={breakable,enhanced,colback=white,boxrule=2.3pt,arc=9pt,
  shadow={3pt}{-3pt}{0pt}{popink},left=14pt,right=14pt,top=11pt,bottom=11pt,before skip=11pt,after skip=15pt,
  fontupper=\popsemi\fontsize{10.6}{14.5}\selectfont\color{popink},
  fontlower=\popsemi\fontsize{10.6}{14.5}\selectfont\color{popink}}}
% Coche dessinée (le glyphe ✓ n'existe pas dans les polices du document)
\newcommand{\popcheck}[1]{\tikz[baseline=-0.5ex]\draw[#1,line width=1.6pt,line cap=round,line join=round](-0.13,0.0)--(-0.04,-0.09)--(0.13,0.1);}
\providecommand{\coche}{\popcheck{popgreen}}
\providecommand{\resultat}[1]{\par\smallskip\noindent\fcolorbox{popgreen}{popgreenL}{\parbox{\dimexpr\linewidth-2\fboxsep-2\fboxrule\relax}{\textbf{Résultat.} #1}}\par\smallskip}
\newcommand{\popboxhead}[3]{\poppill{#1}{#2}{white}\ifx\relax#3\relax\else\hspace{6pt}{\popxbold\fontsize{10.6}{12}\selectfont\color{popink}#3}\fi\par\vspace{7pt}}

\newtcolorbox{aretenir}[1][]{popbase,colframe=popgreen,before upper={\popboxhead{À retenir}{popgreen}{#1}}}
\newtcolorbox{exemplebox}[1][]{popbase,colframe=poporange,before upper={\popboxhead{Exemple}{poporange}{#1}}}
\newtcolorbox{definition}[1][]{popbase,colframe=popblue,before upper={\popboxhead{Définition}{popblue}{#1}}}
\newtcolorbox{propriete}[1][]{popbase,colframe=poppurple,before upper={\popboxhead{Propriété}{poppurple}{#1}}}
\newtcolorbox{methode}[1][]{popbase,colframe=popgold,before upper={\popboxhead{Méthode}{popgold}{#1}}}
\newtcolorbox{attention}[1][]{popbase,colframe=poppink,before upper={\popboxhead{Attention}{poppink}{#1}}}
\newtcolorbox{experience}[1][]{popbase,colframe=popblue,colback=popblueL,before upper={\popboxhead{Expérience}{popblue}{#1}}}
% « Le savais-tu ? » : carte violette pleine (contexte, culture, Cameroun)
\newtcolorbox{savaistu}[1][]{popbase,colback=poppurple,colframe=popink,
  fontupper=\popsemi\fontsize{10.4}{14.2}\selectfont\color{white},
  before upper={\poppill{Le savais-tu\,?}{white}{poppurple}\ifx\relax#1\relax\else\hspace{6pt}{\popxbold\color{white}#1}\fi\par\vspace{7pt}}}
% Exercice résolu : énoncé en haut, \tcblower puis la solution sur fond crème
\newcounter{popexo}\newcounter{popexr}
\newtcolorbox{exoresolu}[1][]{popbase,colframe=poporange,colbacklower=popcream,
  segmentation style={popink,line width=1.2pt,dashed},
  before upper={\stepcounter{popexr}\popboxhead{Exercice résolu \thepopexr}{poporange}{#1}},
  before lower={\poppill{Solution}{popink}{popcream}\par\vspace{6pt}}}
% Exercice (non résolu en place) — la correction est dans la partie Corrigés
\newtcolorbox{exercice}[1][]{popbase,colframe=popink,
  before upper={\stepcounter{popexo}\popboxhead{Exercice \thepopexo}{popink}{#1}}}
% Corrigé
\newtcolorbox{corrige}[1][]{popbase,colframe=popgreen,colback=popgreenL,
  before upper={\popboxhead{Corrigé}{popgreen}{#1}}}
% Objectifs / prérequis / bilan
\newtcolorbox{objectifs}{popbase,colframe=popink,colback=poporangeL,
  before upper={\popboxhead{Objectifs de la leçon}{popink}{}}}
\newtcolorbox{prerequis}{popbase,colframe=popink,colback=popblueL,
  before upper={\popboxhead{Prérequis}{popblue}{}}}
\newtcolorbox{bilan}{popbase,colframe=popink,colback=white,boxrule=2.8pt,
  before upper={\popboxhead{Fiche bilan}{poporange}{L'essentiel à connaître pour l'examen}}}
% Figure encadrée avec légende
\newtcolorbox{popfigure}[1][]{enhanced,breakable=false,colback=white,colframe=popline,boxrule=1.4pt,arc=8pt,
  left=10pt,right=10pt,top=10pt,bottom=8pt,before skip=10pt,after skip=12pt,halign=center,
  lower separated=false,halign lower=center,
  fontlower=\popsemi\fontsize{9}{11.5}\selectfont\color{popmuted},
  #1}
\newcommand{\legende}[1]{\par\vspace{6pt}{\popsemi\fontsize{9}{11.5}\selectfont\color{popmuted}#1}}

% ============================================================
% Signature OnBuch+
% ============================================================
\newcommand{\poplogoinline}[1]{{\poplogo\fontsize{#1}{#1}\selectfont\color{popink}OnBuch\color{poporange}+}}
\newcommand{\signatureonbuch}{%
  \begin{tcolorbox}[enhanced,colback=popink,colframe=popink,boxrule=2.2pt,arc=10pt,
      drop shadow=poporange,left=14pt,right=14pt,top=10pt,bottom=10pt,before skip=0pt,after skip=0pt]
    \tikz[baseline=-0.7ex]\node[circle,fill=popgreen,draw=popcream,line width=1.3pt,minimum size=22pt,inner sep=0]{\popcheck{white}};%
    \hspace{9pt}\begin{minipage}[c]{0.62\linewidth}
      {\popxbold\fontsize{12}{13}\selectfont\color{popcream}Signé\ \textbullet\ {\poplogo OnBuch\color{poporange}+}}\par\vspace{2pt}
      {\raggedright\popsemi\fontsize{8}{10.5}\selectfont\color{popline}\MakeUppercase{Équipe pédagogique OnBuch+}\\\MakeUppercase{Conforme au programme officiel MINESEC}\par}
    \end{minipage}\hfill
    {\poplogo\fontsize{22}{22}\selectfont\color{popcream}OnBuch\color{poporange}+}
  \end{tcolorbox}%
}

% ============================================================
% Page de garde
% #1 titre · #2 matière · #3 niveau · #4 module · #5 n° leçon · #6 durée ·
% #7 description / objectifs (texte ou itemize)
% ============================================================
\newcommand{\pagegarde}[7]{%
  \thispagestyle{empty}
  \newgeometry{a4paper,margin=1.7cm,top=1.6cm,bottom=1.4cm}%
  \begin{tikzpicture}[remember picture,overlay]
    \fill[poporange!18] ([xshift=-3.2cm,yshift=-3.6cm]current page.north east) circle (4.6cm);
    \fill[poppurple!14] ([xshift=2.4cm,yshift=7.6cm]current page.south west) circle (3.8cm);
  \end{tikzpicture}%
  {\poplogo\fontsize{30}{30}\selectfont\color{popink}OnBuch\color{poporange}+}\hfill
  \raisebox{0.6ex}{\poppill{Cours signé \textbullet\ Premium}{popink}{popcream}}\par
  \vspace{1pt}\popsquiggle{poppurple}\par
  \vspace{8pt}
  \begin{tcolorbox}[enhanced,colback=poporange,colframe=popink,boxrule=2.8pt,arc=13pt,
      drop shadow=popink,left=18pt,right=18pt,top=12pt,bottom=13pt,after skip=10pt]
    \poppill{#2 \textbullet\ #3}{popink}{popcream}\hspace{5pt}\poppill{#4}{white}{popink}\par\vspace{8pt}
    {\popdisplay\fontsize{14}{14}\selectfont\color{popink}LEÇON #5}\par\vspace{4pt}
    {\raggedright\hyphenpenalty=10000\exhyphenpenalty=10000\popdisplay\fontsize{25}{27}\selectfont\color{popcream}\MakeUppercase{#1}\par}
    \vspace{8pt}
    {\popxbold\fontsize{9.5}{9.5}\selectfont\color{popink}\MakeUppercase{Durée conseillée : #6}}
  \end{tcolorbox}
  \begin{tcolorbox}[enhanced,colback=white,colframe=popink,boxrule=2.2pt,arc=10pt,
      drop shadow=popink,left=16pt,right=16pt,top=10pt,bottom=10pt,after skip=8pt]
    \poppill{Dans cette leçon}{poppurple}{white}\par\vspace{5pt}
    \popsemi\fontsize{9.3}{12.6}\selectfont\color{popink}#7
  \end{tcolorbox}
  \noindent\begin{minipage}[t]{0.487\linewidth}
    \begin{tcolorbox}[enhanced,colback=popgreen,colframe=popink,boxrule=2.2pt,arc=10pt,
        drop shadow=popink,left=11pt,right=11pt,top=10pt,bottom=10pt,height=3.1cm,valign=center]
      \tcbox[colback=white,colframe=popink,boxrule=1.3pt,arc=5pt,boxsep=0pt,nobeforeafter,
        left=3pt,right=3pt,top=3pt,bottom=3pt,valign=center]{\qrcode[height=1.9cm]{https://play.google.com/store/apps/details?id=com.luvvix.onbuch&hl=fr}}%
      \hspace{8pt}\begin{minipage}[c]{0.5\linewidth}
        {\popdisplay\fontsize{11}{12.5}\selectfont\color{white}\MakeUppercase{Télécharge OnBuch}}\par\vspace{4pt}
        {\raggedright\popsemi\fontsize{8.4}{11}\selectfont\color{white}Gratuit \textbullet\ Google Play\\Tuteur IA, annales, concours, résultats\par}
      \end{minipage}
    \end{tcolorbox}
  \end{minipage}\hfill
  \begin{minipage}[t]{0.487\linewidth}
    \begin{tcolorbox}[enhanced,colback=poppurple,colframe=popink,boxrule=2.2pt,arc=10pt,
        drop shadow=popink,left=11pt,right=11pt,top=10pt,bottom=10pt,height=3.1cm,valign=center]
      \tikz[baseline=-0.7ex]\node[circle,fill=white,draw=popink,line width=1.3pt,minimum size=1.6cm,inner sep=0]{\popdisplay\fontsize{24}{24}\selectfont\color{poppurple}+};%
      \hspace{8pt}\begin{minipage}[c]{0.6\linewidth}
        {\popdisplay\fontsize{11}{12.5}\selectfont\color{white}\MakeUppercase{Passe à OnBuch+}}\par\vspace{4pt}
        {\raggedright\popsemi\fontsize{8.4}{11}\selectfont\color{white}Tous les cours signés, Tuteur IA illimité, fiches PDF — onglet OnBuch+ de l'app\par}
      \end{minipage}
    \end{tcolorbox}
  \end{minipage}
  \vspace{8pt}
  \popcontenu
  \vfill
  \signatureonbuch
  \restoregeometry
  \newpage
}

% Rangée « Ce cours contient » (page de garde)
\newcommand{\poptile}[3]{% #1 couleur #2 gros texte #3 légende
  \begin{tcolorbox}[enhanced,colback=#1,colframe=popink,boxrule=2pt,arc=9pt,shadow={2.5pt}{-2.5pt}{0pt}{popink},
      left=6pt,right=6pt,top=9pt,bottom=9pt,halign=center,valign=center,height=1.75cm,width=\linewidth,nobeforeafter]
    {\popdisplay\fontsize{12.5}{13}\selectfont\color{white}\MakeUppercase{#2}}\par\vspace{3pt}
    {\centering\popsemi\fontsize{7.8}{9.5}\selectfont\color{white}#3\par}
  \end{tcolorbox}}
\newcommand{\popcontenu}{%
  \par\noindent{\popxboldls\fontsize{7.5}{7.5}\selectfont\color{popmuted}CE COURS SIGNÉ CONTIENT}\par\vspace{7pt}
  \noindent\begin{minipage}[t]{0.235\linewidth}\poptile{popblue}{Cours}{complet, illustré, pas à pas}\end{minipage}\hfill
  \begin{minipage}[t]{0.235\linewidth}\poptile{poporange}{Méthodes}{et exercices résolus}\end{minipage}\hfill
  \begin{minipage}[t]{0.235\linewidth}\poptile{poppink}{Exercices}{type examen + corrigés}\end{minipage}\hfill
  \begin{minipage}[t]{0.235\linewidth}\poptile{popgreen}{Fiche bilan}{l'essentiel à retenir}\end{minipage}\par}

% Sommaire
\newtcolorbox{sommaire}{enhanced,colback=white,colframe=popink,boxrule=2.2pt,arc=10pt,
  drop shadow=popink,left=18pt,right=18pt,top=13pt,bottom=13pt,before skip=6pt,after skip=20pt,
  before upper={\poppill{Sommaire}{poppurple}{white}\par\vspace{9pt}\popsemi\fontsize{10.6}{14.5}\selectfont\color{popink}}}

% Signature de fin de cours — poussée tout en bas de la dernière page
\newcommand{\findecours}{%
  \par\vfill\nopagebreak
  \begin{center}\popsquiggle{poporange}\end{center}\vspace{4pt}
  \signatureonbuch
}
\AtBeginDocument{\def\llbracket{\mathopen{[\mkern-3.5mu[}}\def\rrbracket{\mathclose{]\mkern-3.5mu]}}} % intervalles d'entiers (glyphe ⟦ absent de la police)
% Glyphes absents de Fira Math : redéfinis à partir de glyphes disponibles
\AtBeginDocument{%
  \def\setminus{\mathbin{\backslash}}\let\smallsetminus\setminus
  \def\vdots{\mathord{\vbox{\baselineskip3.2pt\lineskiplimit0pt\kern4pt\hbox{.}\hbox{.}\hbox{.}}}}%
  \def\ddots{\mathinner{\mkern1mu\raise6.5pt\hbox{.}\mkern2mu\raise3.6pt\hbox{.}\mkern2mu\raise0.7pt\hbox{.}\mkern1mu}}%
  \def\triangle{\mathord{\scalebox{0.9}{$\symup{\Delta}$}}}%
  \def\nabla{\mathord{\raisebox{\depth}{\scalebox{1}[-1]{$\symup{\Delta}$}}}}%
  \def\checkmark{\popcheck{popdark}}%
  \def\star{\mathbin{\ast}}\def\bigstar{\mathord{\ast}}%
  \def\diamond{\mathbin{\scalebox{0.6}{$\Diamond$}}}%
  \def\bigcup{\mathop{\mathchoice{\raisebox{-0.3em}{\scalebox{1.7}{$\cup$}}}{\scalebox{1.25}{$\cup$}}{\cup}{\cup}}\displaylimits}%
  \def\bigcap{\mathop{\mathchoice{\raisebox{-0.3em}{\scalebox{1.7}{$\cap$}}}{\scalebox{1.25}{$\cap$}}{\cap}{\cap}}\displaylimits}%
  \def\lesseqqgtr{\mathrel{\lessgtr}}\def\bigoplus{\mathop{\oplus}\displaylimits}%
}
\providecommand{\ssi}{si et seulement si} % abréviation fréquente
\AtBeginDocument{\providecommand{\wideparen}{\overparen}} % arc de cercle

%% FIGFIX-BEGIN v5 — correctifs globaux des figures (docs/onbuch-generation/tools/figfix)
\usepackage{adjustbox}\usepackage{etoolbox}\tcbuselibrary{raster}
% toute figure TikZ/pgfplots plus large que la ligne est réduite à la largeur disponible
\BeforeBeginEnvironment{tikzpicture}{\begin{adjustbox}{max width=\linewidth}}
\AfterEndEnvironment{tikzpicture}{\end{adjustbox}}
% légendes pgfplots lisibles par-dessus les courbes
\pgfplotsset{every axis/.append style={legend style={fill=white,fill opacity=0.92,text opacity=1,draw=black!15,inner sep=3pt}}}
% étiquettes d'arêtes (syntaxe edge["..."]) avec fond blanc
\tikzset{every edge quotes/.append style={fill=white,inner sep=1.2pt}}
%% FIGFIX-END
\begin{document}

\pagegarde{Transformation et conservation des fruits de saison}{SVTEEHB}{Tle C}{Module IV : La Biotechnologie}{N10}{3 h}{Cette leçon présente les principes biologiques et les techniques de transformation (jus, purées, concentrés, séchage) et de conservation (froid, traitements thermiques, additifs) des fruits de saison, illustrés par les cas de la mangue et de la tomate. Elle montre comment la maîtrise des causes d'altération des fruits permet de lutter contre les pertes post-récolte et d'améliorer la sécurité alimentaire au Cameroun.
\vspace{4pt}
\begin{multicols}{2}\small
\begin{itemize}
\item À la fin de cette leçon, je sais définir la transformation et la conservation des fruits et distinguer ces deux notions.
\item À la fin de cette leçon, je sais identifier les causes biologiques, chimiques et physiques de l'altération des fruits de saison.
\item À la fin de cette leçon, je sais interpréter des résultats expérimentaux montrant le rôle des micro-organismes et des enzymes dans la détérioration des fruits.
\item À la fin de cette leçon, je sais décrire les étapes de transformation de la mangue (jus, nectar, purée, mangue séchée).
\item À la fin de cette leçon, je sais décrire les étapes de transformation de la tomate (concentré, purée, jus, sauce).
\item À la fin de cette leçon, je sais décrire les principales techniques de conservation des fruits (froid, chaleur, séchage, sucre, additifs) et expliquer leur principe.
\end{itemize}\end{multicols}}

\begin{sommaire}
\begin{multicols}{2}
\begin{enumerate}
\item Généralités et problème des pertes post-récolte
\item Les causes d'altération des fruits
\item Principes généraux des techniques de conservation
\item Transformation et conservation de la mangue
\item Transformation et conservation de la tomate
\item Pratique, qualité et enjeux socio-économiques
\item Activité d'intégration
\item Méthodes et exercices résolus
\item Exercices
\item Corrigés des exercices
\item Fiche bilan
\end{enumerate}
\end{multicols}
\end{sommaire}

\begin{objectifs}
\begin{itemize}
\item À la fin de cette leçon, je sais définir la transformation et la conservation des fruits et distinguer ces deux notions.
\item À la fin de cette leçon, je sais identifier les causes biologiques, chimiques et physiques de l'altération des fruits de saison.
\item À la fin de cette leçon, je sais interpréter des résultats expérimentaux montrant le rôle des micro-organismes et des enzymes dans la détérioration des fruits.
\item À la fin de cette leçon, je sais décrire les étapes de transformation de la mangue (jus, nectar, purée, mangue séchée).
\item À la fin de cette leçon, je sais décrire les étapes de transformation de la tomate (concentré, purée, jus, sauce).
\item À la fin de cette leçon, je sais décrire les principales techniques de conservation des fruits (froid, chaleur, séchage, sucre, additifs) et expliquer leur principe.
\item À la fin de cette leçon, je sais pratiquer, à l'échelle artisanale, une technique de transformation et de conservation d'un fruit de saison.
\item À la fin de cette leçon, je sais évaluer l'intérêt économique et nutritionnel de la transformation des fruits pour la couverture des besoins alimentaires.
\end{itemize}
\end{objectifs}

\begin{prerequis}
\begin{itemize}
\item Les micro-organismes (bactéries, levures, moisissures) et leurs conditions de développement (notion de Seconde/Première).
\item Les enzymes : nature, rôle catalytique, sensibilité à la température et au pH.
\item Les constituants des aliments : eau, glucides, vitamines, pigments (notions de Première sur les biomolécules).
\item La respiration cellulaire et la fermentation (Module sur le métabolisme énergétique).
\item Notions d'hygiène alimentaire et de contamination.
\end{itemize}
\end{prerequis}

\newpage

\coursec{Généralités et problème des pertes post-récolte}

\textbf{Situation d'accroche.} Chaque année, entre mars et juin, les marchés de Douala, de Yaoundé et de Bafoussam débordent de mangues : les étals croulent sous les fruits, les prix s'effondrent — une bassine de mangues se négocie parfois à moins de 500 FCFA — et des tonnes entières finissent par pourrir sur place ou au bord des routes, attirant mouches et mauvaises odeurs. Quelques mois plus tard, en octobre, ces mêmes fruits deviennent introuvables ou hors de prix. Ce paradoxe de l'\emph{abondance suivie de la pénurie} n'est pas une fatalité : il traduit l'absence de techniques de transformation et de conservation adaptées. Comment expliquer ce phénomène ? Quelles techniques, simples ou industrielles, permettraient de transformer et de conserver ces fruits de saison afin de nourrir les populations toute l'année, de réduire les pertes et de créer de la richesse ? C'est à ces problèmes liés à la couverture des besoins alimentaires que répond la \cle{biotechnologie} appliquée aux fruits. Cette première section pose les définitions, décrit les caractéristiques des fruits qui expliquent leur fragilité, puis mesure l'ampleur et les conséquences des pertes post-récolte.

\courssub{Fruits de saison, transformation et conservation : définitions}

\begin{definition}[Fruit de saison]
Un \cle{fruit de saison} est un fruit dont la production, liée au cycle de floraison et de fructification de la plante, est concentrée sur une période limitée de l'année (quelques semaines à quelques mois), au cours de laquelle il est abondant et bon marché, et en dehors de laquelle il devient rare ou indisponible. Au Cameroun : la mangue (mars à juin), la goyave, l'ananas, la tomate fraîche en saison sèche, l'avocat, la papaye à certaines périodes.
\end{definition}

\begin{definition}[Transformation et conservation des fruits]
\begin{itemize}
\item La \cle{transformation} d'un fruit est l'ensemble des opérations (lavage, tri, épluchage, dénoyautage, broyage, cuisson, concentration, fermentation, séchage, conditionnement) qui modifient sa \textbf{forme}, sa \textbf{structure} ou sa \textbf{nature} pour obtenir un produit nouveau : jus, nectar, confiture, purée, concentré, fruits séchés, sirop.
\item La \cle{conservation} d'un fruit est l'ensemble des techniques visant à \textbf{maintenir ses qualités} (nutritionnelles, gustatives, sanitaires) dans le temps, en bloquant ou en ralentissant les causes d'altération : réfrigération, congélation, séchage, appertisation (stérilisation en boîte ou en bocal), pasteurisation, salaison, sucrage, fumage, atmosphère contrôlée.
\end{itemize}
\end{definition}

\begin{attention}[Ne pas confondre transformation et conservation]
Erreur fréquente au Baccalauréat : employer les deux termes comme des synonymes. Retiens la distinction :
\begin{itemize}
\item \textbf{Transformation} = \emph{changement de forme ou de nature} du produit (la mangue devient jus, la tomate devient concentré). Le produit final est différent du produit de départ.
\item \textbf{Conservation} = \emph{maintien des qualités dans le temps} (la mangue reste une mangue, mais réfrigérée, séchée ou appertisée).
\end{itemize}
Les deux notions se combinent souvent : transformer une mangue en jus, \emph{puis} conserver ce jus par pasteurisation. Mais une confiture est une transformation, tandis qu'une chambre froide est une conservation. Dans une copie, précise toujours de quelle opération il s'agit.
\end{attention}

\courssub{Caractéristiques des fruits : des aliments précieux mais fragiles}

Pour comprendre pourquoi les fruits s'altèrent si vite, il faut connaître leur composition. Un fruit mûr est un organe végétal vivant, qui continue à respirer après la récolte (on dit qu'il \emph{respire} : il consomme du dioxygène et rejette du dioxyde de carbone, de la chaleur et de la vapeur d'eau).

\begin{propriete}[Composition moyenne des fruits frais]
\begin{itemize}
\item Une \textbf{forte teneur en eau} : de \SI{80}{\percent} à \SI{95}{\percent} de la masse fraîche (tomate : environ \SI{94}{\percent} ; mangue : environ \SI{82}{\percent}). Cette eau libre est un milieu idéal pour le développement des micro-organismes.
\item Une \textbf{richesse en sucres} solubles (glucose, fructose, saccharose : 8 à \SI{20}{g} pour \SI{100}{g} selon le fruit et sa maturité), source d'énergie pour l'Homme... mais aussi pour les moisissures et les levures.
\item Une \textbf{richesse en vitamines} : vitamine C (acide ascorbique, antioxydant, fragile à la chaleur et à l'oxygène) et provitamine A ($\beta$-carotène, pigment orange de la mangue, précurseur de la vitamine A indispensable à la vision).
\item Une \textbf{acidité variable} : le pH des fruits varie d'environ 2,5 (citron) à 4,5 (tomate mûre). Cette acidité est une barrière partielle : la plupart des bactéries pathogènes ne se développent pas à pH $< 4,5$, ce qui explique que les fruits soient surtout attaqués par les moisissures et les levures, tolérantes aux milieux acides.
\end{itemize}
\end{propriete}

\begin{aretenir}[Pourquoi un fruit pourrit-il vite ?]
Eau abondante + sucres + températures tropicales (25 à \SI{35}{\celsius}) = conditions optimales de croissance microbienne. À cela s'ajoutent les \textbf{enzymes propres du fruit} (polyphénoloxydases responsables du brunissement enzymatique, pectinases qui ramollissent la chair) et la \textbf{respiration} qui épuise les réserves. Un fruit récolté est donc un aliment \emph{périssable} : sans technique de conservation, sa durée de vie se compte en jours sous climat tropical.
\end{aretenir}

\courssub{Le phénomène de saisonnalité : abondance ponctuelle puis pénurie}

La saisonnalité résulte du cycle biologique des arbres fruitiers : une floraison unique (souvent en saison sèche pour le manguier) entraîne une fructification concentrée, donc une récolte massive sur quelques semaines. L'offre explose brutalement, les prix s'effondrent ; puis la production cesse et les prix flambent.

\begin{experience}[Document : suivi de l'offre et de la demande de mangues au Cameroun]
Des relevés de marché effectués sur une année dans les grandes villes camerounaises donnent les volumes indicatifs de mangues fraîches écoulées chaque mois (en tonnes relatives, base 100 au pic), comparés à la demande des consommateurs, supposée régulière.
\textbf{Observations} (figure ci-dessous) : la production est quasi nulle de septembre à février, explose en mars, culmine en avril-mai, puis retombe en juillet. La demande, elle, reste stable toute l'année.
\textbf{Interprétation} : pendant le pic, l'offre dépasse largement la demande $\Rightarrow$ effondrement des prix et gaspillage ; hors saison, la demande n'est pas satisfaite $\Rightarrow$ pénurie et importations.
\textbf{Conclusion} : seule la \emph{transformation/conservation} du surplus de saison permet de lisser l'offre sur l'année.
\end{experience}

\begin{popfigure}
\begin{center}
\begin{tikzpicture}
\begin{axis}[
    ybar, bar width=7pt,
    width=\linewidth, height=6.5cm,
    xlabel={Mois}, ylabel={Volume (base 100 au pic)},
    symbolic x coords={Jan,Fev,Mar,Avr,Mai,Juin,Juil,Aou,Sep,Oct,Nov,Dec},
    xtick=data,
    xticklabels={J,F,M,A,M,J,J,A,S,O,N,D},
    ymin=0, ymax=115,
    legend style={at={(0.5,1.02)},anchor=south,legend columns=2,draw=none},
    ymajorgrids=true, grid style={popline!40},
]
\addplot[fill=popblue,draw=popblue] coordinates {(Jan,2) (Fev,10) (Mar,55) (Avr,95) (Mai,100) (Juin,70) (Juil,15) (Aou,3) (Sep,1) (Oct,1) (Nov,1) (Dec,2)};
\addplot[fill=poporange,draw=poporange] coordinates {(Jan,40) (Fev,40) (Mar,40) (Avr,40) (Mai,40) (Juin,40) (Juil,40) (Aou,40) (Sep,40) (Oct,40) (Nov,40) (Dec,40)};
\legend{Production, Consommation}
\end{axis}
\end{tikzpicture}
\end{center}
\legende{Diagramme comparatif de la production mensuelle de mangues au Cameroun (pic de mars à juin) et de la consommation, supposée régulière. L'écart entre les deux courbes matérialise le surplus perdu en saison et le déficit hors saison.}
\end{popfigure}

\begin{attention}[Lecture de graphique]
Ne confonds pas les deux séries : la production (barres bleues) est \emph{fortement saisonnière}, la consommation (barres oranges) est \emph{quasi constante}. Au Bac, quand on te donne un tel document, commence toujours par décrire les allures (\og la production présente un maximum en avril-mai \fg) avant d'interpréter.
\end{attention}

\courssub{Ampleur des pertes post-récolte : un gâchis mesurable}

\begin{experience}[Document : estimation des pertes le long de la chaîne]
Des enquêtes menées en Afrique subsaharienne (organisations agricoles internationales, services camerounais de l'agriculture) estiment que \textbf{30 à \SI{50}{\percent} des fruits produits sont perdus} entre le champ et le consommateur. Pour la mangue au Cameroun, on retient couramment l'ordre de grandeur suivant : sur \SI{100}{kg} de mangues récoltées, environ \SI{15}{kg} sont perdus à la récolte et au tri (fruits écrasés, blessés, trop mûrs), \SI{10}{kg} pendant le transport (routes dégradées, absence de caisses adaptées, chaleur), \SI{10}{kg} au stockage et sur les étals (mûrissement rapide, pourriture), et \SI{5}{kg} chez le détaillant et le consommateur. Au final, seulement \SI{60}{kg} environ sont réellement consommés.
\end{experience}

\begin{popfigure}
\begin{center}
\begin{tikzpicture}[
    node distance=0.35cm and 0.55cm,
    etape/.style={rectangle, rounded corners=3pt, draw=popblue, fill=popblueL, thick, minimum height=1.1cm, align=center, font=\small},
    perte/.style={rectangle, rounded corners=3pt, draw=poporange, fill=poporangeL, thick, align=center, font=\footnotesize},
    flux/.style={-{Stealth[length=3mm]}, very thick, popdark}
]
\node[etape, align=center] (rec){Récolte\\ \SI{100}{kg}};
\node[etape, right=of rec, align=center] (trans){Transport\\ \SI{85}{kg}};
\node[etape, right=of trans, align=center] (stock){Stockage /\\ marché \SI{75}{kg}};
\node[etape, right=of stock, align=center] (conso){Consommation\\ \SI{60}{kg}};
\node[perte, below=of rec, align=center] (p1){Pertes au tri\\ \SI{-15}{kg} (15 \%)};
\node[perte, below=of trans, align=center] (p2){Pertes au\\ transport \SI{-10}{kg} (10 \%)};
\node[perte, below=of stock, align=center] (p3){Pertes au\\ stockage \SI{-10}{kg} (10 \%)};
\node[perte, below=of conso, align=center] (p4){Pertes chez le\\ détaillant \SI{-5}{kg} (5 \%)};
\draw[flux] (rec) -- (trans);
\draw[flux] (trans) -- (stock);
\draw[flux] (stock) -- (conso);
\draw[flux, poporange] (rec) -- (p1);
\draw[flux, poporange] (trans) -- (p2);
\draw[flux, poporange] (stock) -- (p3);
\draw[flux, poporange] (conso) -- (p4);
\end{tikzpicture}
\end{center}
\legende{Bilan des pertes post-récolte de la mangue, de la récolte au consommateur (ordres de grandeur pour \SI{100}{kg} récoltés). En bleu : les étapes de la chaîne avec les quantités restantes ; en orange : les pertes à chaque étape. Au total, environ 40 \% de la récolte n'atteint jamais le consommateur.}
\end{popfigure}

\begin{exoresolu}[Calculer et interpréter des pertes post-récolte]
\textbf{Énoncé.} Un verger de la plaine des Mbo a produit \SI{12}{t} de mangues en avril. Les pertes sont estimées à 15 \% au tri, 10 \% au transport et 10 \% au stockage (pourcentages appliqués successivement à la quantité restante). 1) Calcule la masse de mangues effectivement commercialisées. 2) Calcule le taux de perte global et compare-le à la somme des trois taux.
\tcblower
\textbf{Solution.}
1) Après le tri : $12 \times (1 - 0{,}15) = 12 \times 0{,}85 = \SI{10,2}{t}$.
Après le transport : $10{,}2 \times 0{,}90 = \SI{9,18}{t}$.
Après le stockage : $9{,}18 \times 0{,}90 = \SI{8,26}{t}$.
La masse commercialisée est d'environ \SI{8,3}{t}.
2) Taux de perte global : $\dfrac{12 - 8{,}26}{12} \times 100 \approx 31{,}2\ \%$.
La somme des taux serait $15 + 10 + 10 = 35\ \%$ : elle \emph{surestime} la perte réelle, car chaque pourcentage s'applique à une quantité déjà diminuée. C'est un piège classique : les pertes successives ne s'additionnent pas directement.
\end{exoresolu}

\begin{savaistu}[Des montagnes de mangues perdues]
Au Cameroun, on estime que \textbf{plus du tiers des mangues produites chaque année est perdu} faute de transformation : cela représente des centaines de milliers de tonnes de fruits. Pourtant, la mangue camerounaise, très parfumée, est recherchée sur les marchés régionaux et européens sous forme de jus, de tranches séchées ou de purée. Quelques unités artisanales et industrielles (séchage solaire amélioré, production de jus dans la région du Littoral et de l'Adamaoua) montrent que ce gâchis peut devenir une source de revenus.
\end{savaistu}

\courssub{Conséquences des pertes post-récolte}

Les pertes post-récolte ne sont pas qu'un problème technique : elles ont des conséquences économiques, alimentaires et sociales majeures.

\begin{itemize}
\item \textbf{Pertes économiques pour les producteurs.} Le surplus invendu pourrit : le travail, les intrants et le temps investis sont perdus. En pleine saison, l'effondrement des prix (une mangue peut tomber à moins de 10 FCFA pièce) ne permet même pas de couvrir les coûts de récolte et de transport.
\item \textbf{Insécurité alimentaire et nutritionnelle.} Les fruits sont des sources irremplaçables de vitamine C et de provitamine A ; leur disparition hors saison aggrave les carences, en particulier chez les enfants (la carence en vitamine A est une cause de cécité et de fragilité face aux infections).
\item \textbf{Importations coûteuses.} Hors saison, jus, concentrés de tomate et conserves sont importés, ce qui coûte des devises au pays alors que la matière première locale a pourri faute de transformation.
\item \textbf{Nuisances environnementales et sanitaires.} Les montagnes de fruits pourris en ville attirent les mouches, encombrent les caniveaux et aggravent le problème de la gestion des déchets urbains.
\end{itemize}

\courssub{Intérêt de la transformation et de la conservation : la réponse de la biotechnologie}

\begin{propriete}[Intérêts de la transformation et de la conservation des fruits]
\begin{enumerate}
\item \textbf{Disponibilité toute l'année} : conserver le surplus de saison permet de lisser l'offre et de couvrir les besoins alimentaires en permanence.
\item \textbf{Création de valeur ajoutée} : un produit transformé se vend beaucoup plus cher que le fruit frais de saison.
\item \textbf{Création d'emplois} : unités de séchage, de production de jus, de confitures et de concentrés emploient main-d'œuvre locale, souvent des femmes et des jeunes.
\item \textbf{Réduction des pertes et des importations} : chaque tonne transformée est une tonne non perdue et une devise économisée.
\item \textbf{Sécurité sanitaire} : les techniques de conservation (chaleur, acidité, séchage) détruisent ou bloquent les micro-organismes responsables de toxi-infections alimentaires.
\end{enumerate}
\end{propriete}

\begin{exemplebox}[La valeur ajoutée de la transformation : le cas de la mangue]
Un kilogramme de mangues fraîches vendu \textbf{100 FCFA} en pleine saison peut valoir environ \textbf{1\,500 FCFA} sous forme de jus conditionné : la valeur est multipliée par 15. Même en retranchant les coûts (sucre, emballage, énergie, main-d'œuvre, estimés à 600--800 FCFA par litre), la marge dégagée rémunère le producteur, le transformateur et le commerçant. C'est toute la logique de la \cle{biotechnologie} alimentaire : utiliser des procédés biologiques et techniques (cuisson, fermentation, séchage, conditionnement) pour convertir une ressource périssable et surabondante en produit stable, sûr et rentable.
\end{exemplebox}

\begin{aretenir}[L'essentiel de la section]
\begin{itemize}
\item Un \cle{fruit de saison} est produit massivement sur une courte période ; il est riche en eau (80--95 \%), en sucres et en vitamines (C, provitamine A), donc très périssable.
\item \cle{Transformation} = changement de forme ou de nature (jus, confiture, concentré) ; \cle{conservation} = maintien des qualités dans le temps (froid, séchage, appertisation).
\item En Afrique subsaharienne, 30 à 50 \% des fruits sont perdus après récolte (tri, transport, stockage, vente), avec de lourdes conséquences économiques et nutritionnelles.
\item Transformer et conserver les fruits assure leur disponibilité toute l'année, crée de la valeur ajoutée et des emplois : c'est une application concrète de la biotechnologie aux besoins alimentaires.
\end{itemize}
\end{aretenir}

\begin{exercice}[Pour vérifier tes acquis]
1) Définis : fruit de saison ; transformation ; conservation. 2) Cite trois caractéristiques des fruits qui expliquent leur altération rapide sous climat tropical. 3) À partir du diagramme production/consommation, explique en quelques lignes l'origine des pertes de mangues en avril-mai. 4) Un producteur perd 20 \% de sa récolte au transport puis 25 \% du reste au stockage : quel est son taux de perte global ?
\end{exercice}

\begin{corrige}[Exercice : Pour vérifier tes acquis]
1) Voir les définitions encadrées : fruit produit abondamment sur une courte période ; transformation = modification de la forme/nature du fruit ; conservation = maintien de ses qualités dans le temps.
2) Forte teneur en eau (80--95 \%), richesse en sucres (nutriments des micro-organismes), respiration et enzymes propres du fruit, températures élevées (25--\SI{35}{\celsius}) favorisant la croissance microbienne.
3) En avril-mai, la production (base 95--100) dépasse largement la consommation (base 40) ; le surplus invendu ne peut être écoulé ni conservé, donc il s'altère : c'est la perte de saison.
4) Reste après transport : $0{,}80$ ; après stockage : $0{,}80 \times 0{,}75 = 0{,}60$. Taux de perte global : $1 - 0{,}60 = 0{,}40$, soit \textbf{40 \%} (et non $20 + 25 = 45\ \%$).
\end{corrige}

\coursec{Les causes d'altération des fruits}

Dans la section précédente, nous avons constaté l'ampleur des pertes post-récolte : au Cameroun, une fraction considérable des mangues et des tomates produites en saison n'atteint jamais le consommateur. Mais \emph{pourquoi} un fruit se dégrade-t-il ? Pour concevoir des techniques de conservation efficaces (section suivante), il faut d'abord identifier les \cle{agents} et les \cle{mécanismes} de l'altération. C'est l'objet de cette section, construite selon la démarche d'investigation scientifique : observation, hypothèse, expérience, interprétation, conclusion.

Trois grandes catégories de causes seront distinguées : les altérations \cle{biologiques} (micro-organismes), les altérations \cle{enzymatiques} (réactions internes au fruit) et les altérations \cle{physico-chimiques} (eau, respiration, oxydations).

\courssub{Les altérations biologiques : l'action des micro-organismes}

\textbf{Observation.} Une mangue intacte, cueillie saine, peut se conserver plusieurs jours à l'air libre. La même mangue, écrasée ou piquée par un insecte, se couvre en quelques jours d'un duvet blanchâtre ou verdâtre, puis pourrit. D'où viennent ces pourritures ?

\textbf{Hypothèse.} Des \cle{micro-organismes} (moisissures, levures, bactéries), présents dans l'air, sur le sol ou à la surface des fruits, pénètrent dans les tissus et les dégradent.

\begin{experience}[Mise en évidence du rôle des micro-organismes dans la pourriture des fruits]
\textbf{Matériel :} mangues ou tomates de même calibre et de même maturité, eau de Javel diluée (solution désinfectante), sacs ou cloches identiques, local à température ambiante (\SI{25}{\celsius} environ).

\textbf{Protocole :} on répartit les fruits en quatre lots :
\begin{enumerate}
\item \textbf{Lot A (témoin) :} fruits sains, simplement lavés à l'eau claire ;
\item \textbf{Lot B :} fruits sains, lavés puis désinfectés à l'eau de Javel diluée et rincés ;
\item \textbf{Lot C :} fruits blessés volontairement (entaille de \SI{1}{cm}), non traités ;
\item \textbf{Lot D :} fruits blessés puis désinfectés.
\end{enumerate}
Tous les lots sont conservés dans les mêmes conditions (même température, même humidité). On note chaque jour le nombre de fruits présentant des taches de pourriture.

\textbf{Observations typiques :}
\begin{center}
\begin{tabularx}{\linewidth}{|l|X|X|X|}
\hline
\rowcolor{popblueL} \textbf{Lot} & \textbf{J+2} & \textbf{J+5} & \textbf{J+8} \\
\hline
A (sain, lavé) & aucune pourriture & 1 fruit taché sur 10 & 3 fruits sur 10 \\
\hline
B (sain, désinfecté) & aucune pourriture & aucune pourriture & 1 fruit sur 10 \\
\hline
C (blessé, non traité) & 4 fruits sur 10 & 9 fruits sur 10 & 10 fruits sur 10 (pourris) \\
\hline
D (blessé, désinfecté) & aucune pourriture & 2 fruits sur 10 & 5 fruits sur 10 \\
\hline
\end{tabularx}
\end{center}

\textbf{Interprétation :} la blessure (comparaison A/C) accélère fortement l'apparition des pourritures : elle ouvre une porte d'entrée aux micro-organismes et supprime la barrière protectrice qu'est l'épiderme. La désinfection (comparaisons A/B et C/D) retarde l'altération : elle réduit le nombre de micro-organismes présents à la surface. Remarque : le lot D finit par s'altérer malgré la désinfection, car la blessure reste une porte d'entrée pour les micro-organismes de l'air qui contaminent le fruit \emph{après} le traitement.

\textbf{Conclusion :} les \cle{micro-organismes} sont les agents de la pourriture des fruits. Ils sont transportés par l'air, l'eau, le sol, les insectes et les mains, et pénètrent surtout par les blessures.
\end{experience}

\begin{definition}[Micro-organismes altérant les fruits]
Les \cle{micro-organismes} responsables de l'altération des fruits sont des êtres vivants microscopiques qui se nourrissent des tissus du fruit et les dégradent :
\begin{itemize}
\item les \cle{moisissures} (champignons filamenteux : \emph{Aspergillus}, \emph{Penicillium}, \emph{Rhizopus}), visibles sous forme de duvets colorés (vert, noir, blanc) ;
\item les \cle{levures} (champignons unicellulaires), qui provoquent la fermentation des sucres (production d'alcool et de gaz, goût et odeur désagréables) ;
\item les \cle{bactéries}, responsables des pourritures molles et visqueuses.
\end{itemize}
\end{definition}

\begin{attention}[Une blessure, même minime, suffit]
L'épiderme et la cuticule du fruit constituent une barrière naturelle. Une piqûre d'insecte, un choc lors de la récolte ou du transport, une coupure au couteau suffisent à permettre l'invasion microbienne. C'est pourquoi, en pratique, on trie les fruits avant conservation : \textbf{tout fruit blessé est écarté ou transformé immédiatement}. Ne dis jamais qu'un fruit pourrit « tout seul » : il y a toujours un agent vivant ou une réaction chimique derrière l'altération.
\end{attention}

\courssub{Les altérations enzymatiques : le brunissement enzymatique}

\textbf{Observation.} Coupe une mangue ou une banane bien mûre et laisse la tranche à l'air libre : en quelques dizaines de minutes, la chair exposée \emph{brunit}. Ce phénomène se produit même en l'absence totale de micro-organismes : ce n'est donc pas une pourriture. Quelle en est la cause ?

\begin{definition}[Brunissement enzymatique]
Le \cle{brunissement enzymatique} est un noircissement des tissus végétaux coupés ou blessés, dû à l'action d'enzymes appelées \cle{polyphénoloxydases} (PPO). Dans la cellule intacte, ces enzymes et leurs substrats (les composés phénoliques) sont séparés dans des compartiments différents. Lorsque les cellules sont rompues, enzymes et substrats entrent en contact et, \textbf{en présence du dioxygène de l'air}, les composés phénoliques sont oxydés en quinones, qui se transforment ensuite en pigments bruns (mélanines). C'est une réaction d'oxydation catalysée par une enzyme.
\end{definition}

\textbf{Hypothèse de travail :} si le brunissement est dû à une enzyme ayant besoin de dioxygène, alors il doit être bloqué par tout traitement qui détruit l'enzyme (chaleur), qui défavorise son activité (acidité) ou qui écarte le dioxygène.

\begin{experience}[Facteurs du brunissement : tranches témoin, citronnées et blanchies]
\textbf{Protocole :} on découpe une même mangue (ou banane) en tranches identiques réparties en trois lots :
\begin{itemize}
\item \textbf{Lot 1 (témoin) :} tranches exposées à l'air libre ;
\item \textbf{Lot 2 :} tranches arrosées de jus de citron (milieu acide, pH voisin de 2) puis exposées à l'air ;
\item \textbf{Lot 3 :} tranches plongées \SI{2}{min} dans l'eau bouillante (\emph{blanchiment}) puis exposées à l'air.
\end{itemize}

\textbf{Observations :}
\begin{itemize}
\item Lot 1 : brunissement net au bout de \SI{30}{min} à \SI{1}{h}, très marqué à J+1 ;
\item Lot 2 : chair claire, quasi inchangée à J+1 ;
\item Lot 3 : chair claire, pas de brunissement (texture légèrement ramollie).
\end{itemize}

\textbf{Interprétation :} le citron abaisse le pH ; or les polyphénoloxydases ont une activité optimale à pH voisin de la neutralité : l'acidité \emph{inhibe} l'enzyme. Le blanchiment soumet l'enzyme à une température proche de \SI{100}{\celsius} : la chaleur \emph{dénature} (détruit la structure spatiale de) la protéine enzymatique, qui devient définitivement inactive. Dans les deux cas, l'oxydation des composés phénoliques ne se produit plus.

\textbf{Conclusion :} le brunissement est une réaction \emph{enzymatique} nécessitant le dioxygène ; il est bloqué par l'\cle{acidification} et par la \cle{dénaturation thermique} des enzymes.
\end{experience}

\begin{popfigure}
\begin{center}
\begin{tikzpicture}
% Lot 1 : témoin
\draw[fill=popgoldL,draw=popdark] (0,0) ellipse (0.9 and 0.55);
\draw[fill=popgold!60!popdark,draw=popdark] (0,2.2) ellipse (0.9 and 0.55);
\node[font=\small\bfseries] at (0,-1.0) {Lot 1 : témoin (air)};
% Lot 2 : citronné
\draw[fill=popgoldL,draw=popdark] (3.6,0) ellipse (0.9 and 0.55);
\draw[fill=popgoldL,draw=popdark] (3.6,2.2) ellipse (0.9 and 0.55);
\draw[fill=popgreenL,draw=popgreen] (3.6,3.2) circle (0.35);
\node[font=\scriptsize] at (3.6,3.2) {citron};
\node[font=\small\bfseries] at (3.6,-1.0) {Lot 2 : citronné (pH acide)};
% Lot 3 : blanchi
\draw[fill=popgoldL,draw=popdark] (7.2,0) ellipse (0.9 and 0.55);
\draw[fill=popgoldL,draw=popdark] (7.2,2.2) ellipse (0.9 and 0.55);
\draw[fill=popblueL,draw=popblue] (6.3,3.0) rectangle (8.1,3.7);
\node[font=\scriptsize] at (7.2,3.35) {eau \SI{100}{\celsius}};
\node[font=\small\bfseries] at (7.2,-1.0) {Lot 3 : blanchi (chaleur)};
% Labels J0 / J+1
\node[font=\small\bfseries,popmuted] at (-1.6,2.2) {J0};
\node[font=\small\bfseries,popmuted] at (-1.6,0) {J+1};
% Annotations observations
\node[font=\scriptsize,align=center] at (0,-1.8) {J+1 : chair \textbf{brunie}};
\node[font=\scriptsize,align=center] at (3.6,-1.8) {J+1 : chair claire\\(enzyme inhibée)};
\node[font=\scriptsize,align=center] at (7.2,-1.8) {J+1 : chair claire\\(enzyme dénaturée)};
\end{tikzpicture}
\end{center}
\legende{Résultats schématisés de l'expérience du brunissement enzymatique : à J0 toutes les tranches sont claires ; à J+1, seule la tranche témoin exposée à l'air a bruni. L'acidification (citron) et le blanchiment (chaleur) bloquent le brunissement.}
\end{popfigure}

\begin{methode}[Interpréter une expérience à témoin (démarche à rédiger au Bac)]
Face à un tableau ou un schéma de résultats expérimentaux, procède toujours ainsi :
\begin{enumerate}
\item \textbf{Identifier la variable modifiée} (le facteur étudié) entre le lot expérimental et le lot témoin : ici, le pH (lot 2) ou la température préalable (lot 3) ; tout le reste doit être identique (même fruit, même durée, même air).
\item \textbf{Comparer les résultats} du lot expérimental à ceux du témoin, en citant les observations (ou les chiffres) : « à J+1, la tranche témoin est brune alors que la tranche citronnée est restée claire ».
\item \textbf{Conclure sur le facteur responsable} : « l'acidification empêche le brunissement, donc le brunissement dépend d'une enzyme sensible au pH ». Une seule variable modifiée à la fois : sinon aucune conclusion valable n'est possible.
\end{enumerate}
\end{methode}

\courssub{Les altérations physico-chimiques}

Même sans micro-organisme et sans coupure, un fruit finit par se dégrader : il est le siège de réactions physiques et chimiques liées au fait qu'\textbf{après la récolte, le fruit reste vivant} et continue de respirer.

\textbf{1. L'évaporation de l'eau (flétrissement).} Le fruit perd de l'eau par transpiration à travers son épiderme. Cette perte d'eau provoque le \cle{flétrissement} : la peau se ride, la masse diminue, la texture devient molle. Une tomate peut perdre plusieurs pourcents de sa masse en quelques jours à l'air chaud et sec. La perte d'eau est accélérée par la chaleur, le vent et la faible humidité de l'air.

\textbf{2. La respiration du fruit.} Le fruit respire : il consomme du dioxygène et des sucres de réserve et rejette du dioxyde de carbone, de l'eau et de la chaleur, selon le bilan général :
\[ \ce{C6H12O6 + 6O2 -> 6CO2 + 6H2O + energie} \]
Cette respiration épuise les réserves, ramollit les tissus et dégage de la chaleur qui, dans un tas de fruits mal aéré, accélère encore la dégradation. Chez certains fruits, la respiration présente un pic brutal au moment du mûrissement.

\begin{definition}[Fruit climatérique]
Un \cle{fruit climatérique} est un fruit dont la respiration et la production d'\cle{éthylène} (hormone végétale gazeuse du mûrissement) augmentent brusquement après la récolte, formant un \emph{pic climatérique} : le fruit peut alors poursuivre son mûrissement après cueillette. La \textbf{mangue}, la banane, la tomate et l'avocat sont climatériques. À l'inverse, l'orange et l'ananas sont \emph{non climatériques} : ils ne mûrissent plus une fois cueillis. Conséquences pratiques : les fruits climatériques peuvent être récoltés à maturité « verte » puis mûris, mais leur pic respiratoire les rend très périssables une fois mûrs ; de plus, l'éthylène dégagé par un fruit mûr accélère le mûrissement des fruits voisins (d'où l'accélération des pourritures dans un tas ou une caisse).
\end{definition}

\textbf{3. L'oxydation des vitamines.} La vitamine C (acide ascorbique), très abondante dans la mangue et la tomate, est fragile : elle s'oxyde à l'air, à la lumière et à la chaleur. Un jus de mangue exposé à l'air perd une part importante de sa vitamine C en quelques heures. Le fruit altéré perd ainsi non seulement son aspect et son goût, mais aussi sa \cle{valeur nutritionnelle}.

\courssub{Les facteurs qui favorisent l'altération}

L'ensemble des observations et expériences précédentes permet de dresser la liste des facteurs qui accélèrent l'altération des fruits :
\begin{itemize}
\item la \textbf{température élevée} : elle accélère la multiplication microbienne, l'activité enzymatique, la respiration et l'évaporation ;
\item l'\textbf{humidité} de l'air : elle favorise le développement des moisissures à la surface des fruits ;
\item la \textbf{lumière} et le \textbf{dioxygène} : ils accélèrent les oxydations (vitamines, pigments, lipides) ;
\item les \textbf{blessures} : portes d'entrée des micro-organismes et mise en contact enzymes/substrats ;
\item le \textbf{mûrissement excessif} : un fruit trop mûr a des tissus ramollis, une peau fragilisée et des sucres très disponibles : il est le substrat idéal des micro-organismes.
\end{itemize}

\begin{popfigure}
\begin{center}
\begin{tikzpicture}
\begin{axis}[
width=\linewidth, height=7cm,
xlabel={Température (\si{\celsius})},
ylabel={Activité (unités arbitraires)},
xmin=-5, xmax=100, ymin=0, ymax=1.25,
xtick={0,10,20,30,40,50,60,70,80,90},
ytick={0,0.5,1},
legend style={at={(0.02,0.97)},anchor=north west,font=\small},
grid=both, grid style={popline!40},
]
% Activité microbienne : optimum vers 30-37, nulle au-delà de 70
\addplot[popcurve,poporange,domain=-5:100,samples=200]
{1.05*exp(-((x-33)^2)/(2*17^2)) * (1/(1+exp((x-62)/3)))};
\addlegendentry{Développement microbien}
% Activité enzymatique : montée puis dénaturation brutale vers 60-70
\addplot[popcurve,popblue,domain=-5:100,samples=200]
{0.9*exp(-((x-40)^2)/(2*22^2)) * (1/(1+exp((x-66)/2.5)))};
\addlegendentry{Activité enzymatique}
% Zones
\draw[popgreen,dashed] (axis cs:5,0) -- (axis cs:5,1.2);
\draw[popgreen,dashed] (axis cs:70,0) -- (axis cs:70,1.2);
\node[font=\scriptsize,popgreen,align=center] at (axis cs:0.5,1.12) {zone du froid\\(ralentissement)};
\node[font=\scriptsize,popdark,align=center] at (axis cs:36,1.12) {optimum\\(25--40 \si{\celsius})};
\node[font=\scriptsize,popgreen,align=center] at (axis cs:85,1.12) {destruction\\(au-delà de 70 \si{\celsius})};
\end{axis}
\end{tikzpicture}
\end{center}
\legende{Vitesse de développement microbien et activité enzymatique en fonction de la température (courbes schématiques). Le froid (0--5 \si{\celsius}) ralentit sans détruire ; l'optimum se situe entre 25 et 40 \si{\celsius} (températures ambiantes tropicales !) ; la chaleur au-delà de 70 \si{\celsius} détruit les micro-organismes et dénature les enzymes. Cette double courbe fonde toutes les techniques de conservation.}
\end{popfigure}

\begin{propriete}[Fondement expérimental des techniques de conservation]
Les expériences précédentes établissent que les \cle{enzymes} et les \cle{micro-organismes} sont :
\begin{itemize}
\item \textbf{détruits ou inactivés par la chaleur} (dénaturation des protéines au-delà de 70--100 \si{\celsius} : blanchiment, pasteurisation, stérilisation) ;
\item \textbf{ralentis par le froid} (réfrigération vers 4 \si{\celsius}, congélation vers $-18$ \si{\celsius}), sans être détruits ;
\item \textbf{ralentis ou inhibés par l'acidité} (pH bas : citron, vinaigre) et par le manque d'eau disponible (dessiccation, forte concentration en sucre ou en sel).
\end{itemize}
Toute technique de conservation consiste donc à agir sur un ou plusieurs de ces facteurs : température, eau, pH, oxygène. Cette propriété sera exploitée dans les sections suivantes.
\end{propriete}

\begin{attention}[Le froid ne tue pas les micro-organismes !]
Erreur très fréquente au Bac : écrire que « le froid détruit les microbes ». \textbf{C'est faux.} Le froid ne fait que \emph{ralentir} leur développement et l'activité enzymatique : les micro-organismes restent vivants, en état de vie ralentie. Dès le retour à la température ambiante, la croissance reprend. C'est pourquoi un fruit sorti du réfrigérateur s'altère ensuite très vite, et pourquoi la chaîne du froid ne doit jamais être interrompue. Seule la \textbf{chaleur} (ou un antiseptique) \emph{détruit} les micro-organismes. Rédige toujours : « le froid bloque ou ralentit », jamais « le froid tue ».
\end{attention}

\begin{exoresolu}[Exploitation de la courbe « activité--température »]
\textbf{Énoncé.} Un producteur de mangues de la région du Nord conserve ses fruits dans un hangar où la température atteint \SI{35}{\celsius} en journée. Ses mangues pourrissent en moins d'une semaine. (1) À l'aide de la courbe activité--température, expliquer cette altération rapide. (2) Proposer deux mesures fondées sur la courbe pour réduire les pertes, en précisant le mécanisme invoqué.

\tcblower
\textbf{Solution détaillée.}
(1) La courbe montre que la zone de 25 à \SI{40}{\celsius} correspond à l'\emph{optimum} de développement microbien et d'activité enzymatique. À \SI{35}{\celsius}, les moisissures, levures et bactéries présentes sur les fruits se multiplient à vitesse maximale, et les enzymes du fruit (dont les polyphénoloxydases) fonctionnent à pleine activité ; de plus, la respiration du fruit (fruit climatérique) et l'évaporation de l'eau sont accélérées. Toutes les causes d'altération sont donc simultanément favorisées : d'où la pourriture en moins d'une semaine.

(2) Deux mesures possibles :
\begin{itemize}
\item \textbf{Réfrigérer} les mangues vers 10--13 \si{\celsius} (en dessous de 10 \si{\celsius}, la mangue, fruit tropical, subit des dommages dus au froid : brunissement de la peau, chair qui ne mûrit plus normalement) : on se place dans la zone de ralentissement de la courbe, ce qui freine microbes et enzymes sans les détruire ;
\item \textbf{Transformer rapidement} les fruits (séchage, jus bouilli puis mis en bouteille, confiture) : le chauffage au-delà de \SI{70}{\celsius} place le produit dans la zone de destruction des micro-organismes et de dénaturation des enzymes, ce qui stabilise le produit.
\end{itemize}
\end{exoresolu}

\begin{savaistu}[Le citron, conservateur de grand-mère... validé par la science]
Au Cameroun comme ailleurs, on arrose depuis longtemps les avocats, bananes et mangues coupés de jus de citron « pour qu'ils ne noircissent pas ». Cette pratique empirique repose exactement sur le mécanisme étudié ici : l'acide citrique abaisse le pH et inhibe les polyphénoloxydases. L'industrie agroalimentaire utilise le même principe avec l'acide ascorbique (vitamine C), à la fois acidifiant et antioxydant, dans les jus de fruits commerciaux.
\end{savaistu}

\begin{exercice}[À toi de jouer]
On conserve quatre lots de tomates mûres : lot 1 à \SI{30}{\celsius} à l'air libre ; lot 2 à \SI{4}{\celsius} ; lot 3 à \SI{30}{\celsius} après blanchiment de \SI{2}{min} ; lot 4 à \SI{30}{\celsius} sous film plastique fermé. (1) Prévoir, en le justifiant, le lot qui s'altérera le plus vite et celui qui se conservera le plus longtemps. (2) Expliquer pourquoi le lot 3, remis à \SI{30}{\celsius}, finit quand même par s'altérer. (3) Rédiger la conclusion générale sur les causes d'altération des fruits.
\end{exercice}

\begin{corrige}[Exercice]
(1) Le lot 1 s'altère le plus vite : \SI{30}{\celsius} est proche de l'optimum microbien et enzymatique, et l'air fournit le dioxygène des oxydations. Le lot 2 se conserve le plus longtemps : le froid ralentit simultanément microbes, enzymes, respiration et évaporation.
(2) Le blanchiment a détruit les enzymes et les micro-organismes \emph{présents au moment du traitement}, mais le fruit n'est pas protégé ensuite : de nouveaux micro-organismes de l'air le contaminent à nouveau et se développent à \SI{30}{\celsius}. D'où l'intérêt du conditionnement hermétique après chauffage.
(3) Les fruits s'altèrent sous l'action conjuguée des micro-organismes (pourritures), des enzymes propres au fruit (brunissement, ramollissement) et de réactions physico-chimiques (perte d'eau, respiration, oxydations), ces phénomènes étant accélérés par la chaleur, l'humidité, l'oxygène, la lumière et les blessures.
\end{corrige}

\begin{aretenir}[L'essentiel de la section]
\begin{itemize}
\item Les fruits s'altèrent pour trois grandes raisons : les \cle{micro-organismes} (moisissures, levures, bactéries, entrant surtout par les blessures), les \cle{enzymes} du fruit (brunissement enzymatique par les polyphénoloxydases en présence d'$\ce{O2}$) et les phénomènes \cle{physico-chimiques} (évaporation, respiration climatérique, oxydation des vitamines).
\item Le brunissement est bloqué par l'acidification (citron) et par la chaleur (blanchiment) : preuve expérimentale de son origine enzymatique.
\item Température, humidité, oxygène, lumière, blessures et mûrissement excessif favorisent l'altération.
\item Chaleur $\Rightarrow$ destruction ; froid $\Rightarrow$ simple ralentissement ; acidité et manque d'eau $\Rightarrow$ inhibition. C'est le fondement de toutes les techniques de conservation étudiées dans la suite de la leçon.
\end{itemize}
\end{aretenir}

\coursec{Principes généraux des techniques de conservation}

Dans la section précédente, nous avons identifié les \cle{agents responsables de l'altération des fruits} : les micro-organismes (bactéries, levures, moisissures), les enzymes propres du fruit (oxydases, pectinases) et les réactions chimiques d'oxydation. Nous savons donc désormais \emph{qui} attaque nos mangues et nos tomates pendant la saison d'abondance. Reste une question pratique, celle que se posent chaque année les ménagères de Ngaoundéré comme les industriels de Douala : \emph{comment les arrêter ?} C'est l'objet de cette section : dégager les \cle{principes scientifiques} qui fondent toutes les techniques de conservation, avant de les appliquer concrètement à la mangue et à la tomate dans les sections suivantes.

\courssub{Le principe général : bloquer les agents d'altération}

Toute technique de conservation repose sur une logique unique, qu'il faut bien comprendre : un fruit s'altère parce que des micro-organismes s'y développent et parce que des réactions enzymatiques s'y déroulent. Conserver, c'est donc agir sur les \cle{conditions du milieu} pour :

\begin{itemize}
\item \cle{détruire} les micro-organismes et dénaturer les enzymes (action curative : chaleur) ;
\item \cle{empêcher ou ralentir} leur développement (action préventive : froid, séchage, acidification, sucre, sel, additifs) ;
\item \cle{protéger} le produit contre une nouvelle contamination (emballage hermétique).
\end{itemize}

\begin{definition}[Conservation des fruits]
La \cle{conservation} est l'ensemble des techniques visant à maintenir les qualités nutritionnelles et organoleptiques (goût, couleur, texture) d'un fruit en empêchant ou en ralentissant le développement des micro-organismes et les réactions enzymatiques d'altération, et/ou en détruisant les agents d'altération déjà présents.
\end{definition}

\begin{popfigure}
\begin{tikzpicture}[
  grow=right,
  level distance=3.8cm,
  level 1/.style={sibling distance=10mm},
  every node/.style={font=\small},
  racine/.style={rectangle, rounded corners, draw=popdark, fill=popblueL, thick, align=center, text width=3.0cm},
  principe/.style={rectangle, rounded corners, draw=popblue, fill=popblue!15, align=center, text width=2.6cm},
  tech/.style={rectangle, draw=popmuted, fill=popcream, align=left, text width=4.8cm, font=\footnotesize}
]
\node[racine] {Techniques de conservation}
  child { node[principe] {Emballage / conditionnement} 
    child { node[tech] {Boîtes, bocaux hermétiques (appertisation), mise sous vide, atmosphère modifiée} } }
  child { node[principe] {Additifs} 
    child { node[tech] {Antioxydants (acide ascorbique), conservateurs (sorbate, benzoate, sulfites)} } }
  child { node[principe] {Acidification} 
    child { node[tech] {Vinaigre, citron, fermentation lactique (pH $< 4{,}5$)} } }
  child { node[principe] {Réduction de l'eau disponible} 
    child { node[tech] {Séchage, concentration, sucre (confitures), sel} } }
  child { node[principe] {Chaleur} 
    child { node[tech] {Pasteurisation (\SIrange{70}{90}{\celsius}), stérilisation (\SIrange{110}{121}{\celsius})} } }
  child { node[principe] {Froid} 
    child { node[tech] {Réfrigération (\SIrange{0}{4}{\celsius}), congélation (\SI{-18}{\celsius})} } };
\end{tikzpicture}
\legende{Classification des principales techniques de conservation des fruits selon leur principe d'action. Chaque principe agit sur un facteur indispensable à la vie microbienne : température, eau disponible, pH, oxygène ou contact avec le milieu extérieur.}
\end{popfigure}

\courssub{Conservation par le froid : ralentir sans détruire}

\textbf{Observation.} Une mangue mûre laissée à l'air libre à Maroua (température ambiante voisine de \SI{35}{\celsius}) se gâte en deux ou trois jours ; placée dans un réfrigérateur, elle reste consommable une à deux semaines. Le froid ne tue donc pas les microbes : il les \emph{endort}.

\textbf{Interprétation.} La vitesse des réactions enzymatiques et de la multiplication microbienne dépend fortement de la température : elle diminue quand la température baisse. Le froid est une technique de \cle{ralentissement}, réversible : dès que le produit se réchauffe, l'altération reprend.

\begin{definition}[Réfrigération et congélation]
\begin{itemize}
\item La \cle{réfrigération} consiste à maintenir le produit entre \SI{0}{\celsius} et \SI{4}{\celsius}. Le métabolisme microbien et les réactions enzymatiques sont fortement ralentis, mais pas arrêtés : la conservation est de courte durée (quelques jours à quelques semaines).
\item La \cle{congélation} consiste à porter le produit à environ \SI{-18}{\celsius}. L'eau du fruit cristallise et devient indisponible : le développement microbien est \cle{bloqué} (mais les microbes ne sont pas tués). La conservation peut atteindre plusieurs mois.
\end{itemize}
\end{definition}

\begin{attention}[Limites du froid]
Le froid n'est pas une panacée. D'une part, la congélation forme des cristaux de glace qui éclatent les cellules : à la décongélation, la \cle{texture} du fruit devient molle (purée). D'autre part, certaines \cle{vitamines} fragiles (vitamine C notamment) sont partiellement détruites, et les fruits tropicaux comme la mangue supportent mal les températures proches de \SI{0}{\celsius} (brunissement, « maladie du froid »). Enfin, toute rupture de la chaîne du froid relance immédiatement l'altération.
\end{attention}

\courssub{Conservation par la chaleur : détruire les agents d'altération}

\textbf{Un peu d'histoire pour comprendre.} Au XIX\textsuperscript{e} siècle, les vins et les laits « tournaient » régulièrement, causant des pertes énormes. Louis \cle{Pasteur} démontra que ces altérations sont dues à des micro-organismes vivants (théorie microbienne des fermentations), et qu'un chauffage modéré suffit à les éliminer sans dénaturer le produit : la \cle{pasteurisation} était née.

\begin{definition}[Pasteurisation]
La \cle{pasteurisation} est un traitement thermique modéré, entre \SI{70}{\celsius} et \SI{90}{\celsius} pendant quelques secondes à quelques minutes, qui détruit les micro-organismes pathogènes et la plupart des formes végétatives, ainsi que les enzymes, \emph{sans détruire les spores bactériennes}. Le produit pasteurisé doit donc être conservé au froid et consommé rapidement. Exemple : jus de mangue chauffé à \SI{85}{\celsius} pendant \SI{1}{\minute}.
\end{definition}

\begin{definition}[Stérilisation]
La \cle{stérilisation} est un traitement thermique intense, entre \SI{110}{\celsius} et \SI{121}{\celsius} (sous pression, en autoclave) pendant \SIrange{15}{30}{\minute}, qui détruit \emph{tous} les micro-organismes, \cle{y compris leurs spores}. Le produit stérilisé en récipient hermétique se conserve plusieurs années à température ambiante. Exemple : concentré de tomate en boîte.
\end{definition}

\textbf{Pourquoi un couple temps--température ?} La mort microbienne par la chaleur est progressive : plus la température est élevée, plus le temps nécessaire pour détruire une population microbienne est court. C'est ce que montre la courbe de survie ci-dessous.

\begin{popfigure}
\begin{center}
\begin{tikzpicture}
\begin{semilogyaxis}[
  width=0.9\linewidth, height=6.5cm,
  xlabel={Temps de chauffage (min)},
  ylabel={Nombre de micro-organismes survivants},
  xmin=0, xmax=12, ymin=1, ymax=100000000,
  legend style={at={(0.02,0.02)}, anchor=south west, font=\small},
  grid=major, grid style={popline!50},
  tick label style={font=\small}, label style={font=\small}
]
\addplot[popcurve, popblue, domain=0:12, samples=100] {100000000*exp(-0.35*x)};
\addlegendentry{\SI{65}{\celsius} (chauffage insuffisant)}
\addplot[popcurve, poporange, domain=0:12, samples=100] {100000000*exp(-1.1*x)};
\addlegendentry{\SI{75}{\celsius} (pasteurisation)}
\addplot[popcurve, poppink, domain=0:12, samples=100] {100000000*exp(-3*x)};
\addlegendentry{\SI{90}{\celsius}}
\end{semilogyaxis}
\end{tikzpicture}
\end{center}
\legende{Courbes de survie microbienne en fonction du temps de chauffage à trois températures (échelle logarithmique). À \SI{65}{\celsius}, la décroissance est lente : des microbes survivent encore après \SI{12}{\minute}. À \SI{75}{\celsius}, la population est pratiquement éliminée en quelques minutes : c'est le principe du couple \cle{temps--température} de la pasteurisation. À \SI{90}{\celsius}, la destruction est quasi immédiate.}
\end{popfigure}

\begin{attention}[Erreur fréquente : confondre pasteurisation et stérilisation]
Au Baccalauréat comme en pratique, ne confonds jamais ces deux traitements. La \cle{pasteurisation} (\SIrange{70}{90}{\celsius}) détruit les formes végétatives mais \emph{pas les spores} : un jus pasteurisé doit impérativement être conservé au réfrigérateur et consommé en quelques jours. Seule la \cle{stérilisation} (\SIrange{110}{121}{\celsius}) élimine les spores et permet une conservation longue à température ambiante. Écrire « le lait pasteurisé est stérile » est une faute scientifique grave.
\end{attention}

\courssub{Conservation par réduction de l'eau disponible}

\textbf{Observation.} Les moisissures se développent sur une mangue fraîche, jamais sur un fruit sec ni sur du sucre en poudre. Pourtant, le fruit sec contient encore de l'eau ! Comment l'expliquer ?

\textbf{Le concept clé : eau libre et eau liée.} L'eau d'un aliment existe sous deux formes :
\begin{itemize}
\item l'\cle{eau libre}, disponible pour les réactions chimiques et le métabolisme microbien ;
\item l'\cle{eau liée}, retenue par les molécules dissoutes (sucre, sel) ou les macromolécules, et \emph{indisponible} pour les micro-organismes.
\end{itemize}
Ce qui compte pour les microbes, ce n'est pas la quantité totale d'eau, mais la quantité d'eau \emph{libre} (activité de l'eau, notée $a_w$). Réduire l'eau libre, c'est affamer les microbes.

\begin{definition}[Techniques de réduction de l'eau disponible]
\begin{itemize}
\item Le \cle{séchage} élimine l'eau par évaporation : \emph{naturel} au soleil (mangues et tomates séchées, très pratiqué dans le Nord-Cameroun) ou \emph{artificiel} en séchoir (air chaud à \SIrange{50}{70}{\celsius}, plus rapide et plus hygiénique).
\item La \cle{concentration} élimine une partie de l'eau par chauffage (concentré de tomate).
\item L'ajout de \cle{sucre} en forte quantité lie l'eau : c'est le principe des \cle{confitures}, qui exigent au minimum \SI{60}{\%} de sucre pour être stables.
\item L'ajout de \cle{sel} agit de même (surtout pour les légumes et condiments).
\end{itemize}
\end{definition}

\begin{exemplebox}[La confiture de mangue]
Une confiture de mangue contient typiquement \SI{1}{\kg} de pulpe pour \SIrange{700}{1000}{\g} de sucre. Le sucre dissout capte l'eau libre de la pulpe : les levures et moisissures ne trouvent plus d'eau disponible pour se développer. Si l'on réduit trop le sucre (confiture « allégée » à \SI{30}{\%}), le produit moisit en quelques semaines : il faut alors le conserver au réfrigérateur après ouverture.
\end{exemplebox}

\courssub{Conservation par acidification : le pouvoir du pH}

\textbf{Observation.} Les citrons et les tamarinades se conservent bien mieux que les bananes mûres. Les cornichons plongés dans le vinaigre se gardent des mois. Le point commun : l'\cle{acidité}.

\begin{propriete}[Effet du pH sur la croissance microbienne]
Un pH inférieur à \num{4,5} empêche le développement de la quasi-totalité des bactéries, notamment les bactéries pathogènes (dont \emph{Clostridium botulinum}, agent du botulisme). Seules certaines levures et moisissures tolèrent encore ce milieu acide. C'est pourquoi les fruits naturellement acides (tomate : pH $\approx \num{4,2}$ ; mangue : pH $\approx \num{3,5}$--\num{4,0}) présentent une sécurité relative et se prêtent bien à la conservation par simple pasteurisation.
\end{propriete}

Les moyens pratiques d'abaisser le pH sont :
\begin{itemize}
\item l'ajout de \cle{vinaigre} (acide acétique) ou de jus de \cle{citron} (acide citrique) ;
\item la \cle{fermentation lactique} : des bactéries lactiques transforment les sucres en acide lactique, qui abaisse naturellement le pH (principe des légumes lacto-fermentés).
\end{itemize}

\courssub{Conservation par additifs chimiques}

\begin{definition}[Additifs de conservation]
\begin{itemize}
\item Les \cle{antioxydants}, comme l'\cle{acide ascorbique} (vitamine C), empêchent les réactions d'oxydation responsables du brunissement des pulpes de fruits exposées à l'air.
\item Les \cle{conservateurs chimiques} inhibent le développement microbien : \cle{sorbate} de potassium (antifongique), \cle{benzoate} de sodium (efficace en milieu acide), \cle{sulfites} (antioxydants et antimicrobiens, utilisés pour les fruits secs).
\end{itemize}
\end{definition}

\begin{attention}[Précautions d'emploi des additifs]
Les additifs ne sont efficaces et sûrs qu'aux \cle{doses réglementaires}. Un surdosage de sulfites peut déclencher des allergies (asthme) ; le benzoate n'agit qu'en milieu acide (pH $< \num{4,5}$) et est inutile en milieu neutre. Les additifs ne dispensent jamais des règles d'hygiène : ils complètent une bonne pratique, ils ne la remplacent pas.
\end{attention}

\courssub{Conservation par conditionnement : isoler le produit du milieu extérieur}

Détruire les microbes ne suffit pas si le produit est aussitôt recontaminé par l'air, les insectes ou les mains. Le \cle{conditionnement} protège le fruit transformé :
\begin{itemize}
\item emballage \cle{hermétique} (sachets, bocaux, boîtes métalliques) ;
\item \cle{mise sous vide} : l'air est retiré, privant les microbes aérobies d'oxygène ;
\item \cle{atmosphère modifiée} : l'air est remplacé par un mélange de gaz (diazote, dioxyde de carbone) défavorable aux microbes ;
\item \cle{appertisation} : chauffage du produit \emph{dans son récipient hermétiquement clos} (boîte ou bocal), qui combine destruction microbienne et protection contre la recontamination.
\end{itemize}

\begin{definition}[Appertisation]
L'\cle{appertisation} est un procédé de conservation consistant à enfermer l'aliment dans un récipient hermétique (boîte métallique, bocal de verre) puis à le chauffer à une température suffisante (\SIrange{110}{121}{\celsius}) pour détruire les micro-organismes et leurs spores. Le produit appertisé se conserve plusieurs années à température ambiante.
\end{definition}

\begin{savaistu}[Nicolas Appert, l'inventeur de la conserve]
En 1809, le confiseur français Nicolas \cle{Appert} remporte le prix offert par Napoléon Bonaparte pour nourrir ses armées en campagne : il a découvert que des aliments chauffés longuement dans des bouteilles de verre bouchées se conservent des mois sans se gâter. Fait remarquable : Appert appliquait un procédé \emph{avant} de le comprendre — il faudra attendre Pasteur, un demi-siècle plus tard, pour que la science explique que la chaleur détruit des micro-organismes. La conserve est née, et elle porte encore son nom : l'appertisation.
\end{savaistu}

\courssub{Comment choisir la bonne technique ?}

\begin{methode}[Choisir une technique de conservation adaptée à un fruit]
\begin{enumerate}
\item \textbf{Évaluer l'acidité du fruit.} Si le pH est inférieur à \num{4,5} (mangue, tomate), une simple pasteurisation suffit, car les bactéries dangereuses ne peuvent pas s'y développer. Si le pH est supérieur à \num{4,5}, la stérilisation s'impose.
\item \textbf{Tenir compte de la teneur en eau et de la texture.} Un fruit destiné à être séché doit supporter la déshydratation (mangue : excellente en lamelles séchées). Un fruit fragile se transforme plutôt en purée ou en jus.
\item \textbf{Définir la destination du produit.} Vente locale rapide : réfrigération ou confiture. Exportation ou stockage long : séchage, appertisation ou congélation.
\item \textbf{Prendre en compte les moyens disponibles.} En zone rurale sans électricité, le séchage solaire, la confiture et l'acidification sont les techniques les plus réalistes ; la congélation exige une chaîne du froid coûteuse.
\item \textbf{Vérifier l'hygiène et le conditionnement.} Quelle que soit la technique, un récipient propre et hermétique est indispensable pour éviter la recontamination.
\end{enumerate}
\end{methode}

\begin{exoresolu}[Choisir et justifier une technique de conservation]
\textbf{Énoncé.} Une coopérative de femmes de Garoua récolte chaque année un excédent de mangues très mûres (pH mesuré : \num{3,8}) qu'elle ne peut pas écouler sur le marché local. Elle dispose de marmites, de bocaux en verre, de sucre et d'un espace ensoleillé, mais pas d'électricité. Proposer deux techniques de conservation adaptées et justifier scientifiquement chaque choix.
\tcblower
\textbf{Solution.} \textbf{1. Confiture de mangue.} La mangue étant acide (pH \num{3,8} $<$ \num{4,5}), les bactéries pathogènes ne peuvent s'y développer. L'ajout de sucre à au moins \SI{60}{\%} lie l'eau libre, ce qui bloque levures et moisissures ; la cuisson (environ \SI{100}{\celsius}) détruit les formes végétatives et les enzymes. Le remplissage à chaud des bocaux propres assure un conditionnement hermétique. Conservation : plusieurs mois à température ambiante. \textbf{2. Séchage solaire en lamelles.} L'exposition au soleil (éventuellement sur claies surélevées, à l'abri des poussières et des mouches) évapore l'eau jusqu'à ce que l'eau libre restante soit insuffisante pour le développement microbien. Le produit sec, léger et stable, se conserve plusieurs mois en sachets hermétiques. \textbf{Justification du choix :} ces deux techniques ne demandent ni électricité ni chaîne du froid, exploitent l'acidité naturelle du fruit et correspondent aux équipements disponibles. La congélation et la stérilisation en autoclave seraient ici irréalistes.
\end{exoresolu}

\begin{aretenir}[L'essentiel de la section]
\begin{itemize}
\item Conserver, c'est \cle{détruire} les agents d'altération (chaleur) ou \cle{bloquer leur développement} (froid, séchage, sucre, sel, acidité, additifs), puis \cle{protéger} le produit (emballage hermétique).
\item Froid : réfrigération (\SIrange{0}{4}{\celsius}, ralentit) ; congélation (\SI{-18}{\celsius}, bloque sans tuer).
\item Chaleur : pasteurisation (\SIrange{70}{90}{\celsius}, sans les spores, conservation au froid) ; stérilisation (\SIrange{110}{121}{\celsius}, spores détruites).
\item Réduction de l'eau libre : séchage, concentration, sucre ($\geq \SI{60}{\%}$), sel.
\item Acidification : pH $< \num{4,5}$ bloque les bactéries pathogènes.
\item Appertisation : chauffage en récipient clos hermétique (Appert, 1809).
\end{itemize}
\end{aretenir}

Ces principes généraux étant posés, nous sommes maintenant armés pour les mettre en œuvre concrètement : la section suivante les appliquera pas à pas au fruit roi des savanes camerounaises, la mangue.

\coursec{Transformation et conservation de la mangue}

Nous venons de voir les principes généraux des techniques de conservation : abaisser l'\cle{eau disponible} ($a_w$), abaisser le \cle{pH}, détruire ou bloquer les micro-organismes par la chaleur ou le froid, et protéger le produit de toute recontamination par un emballage hermétique. Il est temps de mettre ces principes en application sur un fruit emblématique du Cameroun : la \cle{mangue}. Chaque année, entre avril et juillet, les marchés de Yaoundé, Douala, Garoua ou Ngaoundéré débordent de mangues... puis, quelques semaines plus tard, le fruit devient rare et cher. Transformer et conserver la mangue, c'est lutter contre ce gâchis et créer de la valeur.

\courssub{Choix et préparation des fruits : la première étape décisive}

\begin{definition}[Matière première de qualité]
Toute transformation réussie commence par le choix des fruits : on retient des mangues \cle{mûres} (bonne teneur en sucre et en arômes) mais \cle{saines}, c'est-à-dire sans blessures, sans coups, sans taches de pourriture ni piqûres de mouches des fruits. Le tri est suivi d'un \cle{lavage} à l'eau potable propre, éventuellement additionnée de quelques gouttes d'eau de Javel diluée, puis d'un égouttage.
\end{definition}

Pourquoi tant de rigueur ? Une blessure est une porte d'entrée pour les moisissures et les bactéries ; une zone pourrie concentre des millions de micro-organismes et parfois leurs toxines. Le lavage élimine les poussières, les spores et une grande partie de la flore microbienne de surface. Le tri, lui, garantit l'homogénéité du lot : des fruits de maturités très différentes donneraient un jus au goût instable.

\begin{attention}[La cuisson ne rattrape pas tout]
Erreur fréquente : utiliser des fruits pourris ou blessés « parce que de toute façon, la cuisson va tout stériliser ». C'est \textbf{faux}. Si la chaleur tue bien les micro-organismes vivants, les \cle{toxines} produites par certaines moisissures (mycotoxines, comme l'aflatoxine ou la patuline) sont \textbf{thermostables} : elles résistent aux températures de cuisson et de pasteurisation (savoir admis). Un fruit pourri doit être \textbf{jeté}, jamais transformé. De même, on ne « découpe » pas la partie pourrie pour garder le reste : les toxines diffusent au-delà de la zone visiblement altérée.
\end{attention}

\courssub{Fabrication du jus et du nectar de mangue}

\begin{definition}[Jus et nectar]
Le \cle{jus} de mangue est le produit obtenu à partir de la pulpe du fruit, sans dilution notable (100 \% fruit). Le \cle{nectar} est une boisson obtenue en diluant la pulpe ou le jus avec de l'eau et en ajoutant du sucre : il contient typiquement 25 à 50 \% de fruit. Dans les deux cas, une \cle{pasteurisation} (chauffage modéré puis refroidissement rapide) assure la conservation en détruisant les micro-organismes altérants et pathogènes.
\end{definition}

Le procédé suit une chaîne logique où chaque étape a une fonction précise :

\begin{itemize}
\item \textbf{Lavage} : élimination des salissures et de la flore de surface ;
\item \textbf{Épluchage} : la peau de la mangue est épaisse, fibreuse et contient des composés irritants ; elle est retirée au couteau ou à l'éplucheur ;
\item \textbf{Dénoyautage} : séparation de la pulpe du noyau aplati central ;
\item \textbf{Broyage} : la pulpe est transformée en purée fine au mixeur ou au moulin ;
\item \textbf{Raffinage} : passage au tamis pour éliminer les fibres et obtenir une texture homogène ;
\item \textbf{Formulation} : pour le nectar, ajout d'eau potable et de sucre (et éventuellement d'un peu de jus de citron pour abaisser le pH) ;
\item \textbf{Pasteurisation} : chauffage à environ \SI{85}{\celsius} pendant 5 à 10 minutes (ou \SI{90}{\celsius} pendant 1 à 2 minutes) ;
\item \textbf{Conditionnement à chaud} : remplissage immédiat des bouteilles ou bocaux préalablement lavés et ébouillantés, puis bouchage hermétique ;
\item \textbf{Refroidissement rapide} puis \textbf{stockage} à l'abri de la lumière et de la chaleur.
\end{itemize}

\begin{popfigure}
\begin{center}
\begin{tikzpicture}[
font=\small,
etape/.style={draw=popblue, fill=popblueL, rounded corners=2pt, minimum width=4.6cm, minimum height=0.85cm, align=center},
crit/.style={draw=poporange, fill=poporangeL, rounded corners=2pt, minimum width=3.6cm, minimum height=0.85cm, align=center},
fleche/.style={-{Stealth[length=2.5mm]}, thick, popdark}
]
\node[etape] (e1) at (0,0) {1. Tri et lavage des mangues};
\node[etape] (e2) at (0,-1.2) {2. Épluchage et dénoyautage};
\node[etape] (e3) at (0,-2.4) {3. Broyage (pulpe) et raffinage};
\node[etape] (e4) at (0,-3.6) {4. Formulation (eau + sucre pour le nectar)};
\node[etape] (e5) at (0,-4.8) {5. Pasteurisation};
\node[etape] (e6) at (0,-6.0) {6. Conditionnement à chaud, bouchage};
\node[etape] (e7) at (0,-7.2) {7. Refroidissement rapide et stockage};
\node[crit, align=center] (c5) at (5.6,-4.8){\SI{85}{\celsius} / 5--10 min\\ ou \SI{90}{\celsius} / 1--2 min};
\node[crit, align=center] (c6) at (5.6,-6.0){Remplissage $>$ \SI{80}{\celsius}\\ récipients ébouillantés};
\node[crit, align=center] (c7) at (5.6,-7.2){Eau froide / air\\ stockage frais, à l'ombre};
\foreach \a/\b in {e1/e2,e2/e3,e3/e4,e4/e5,e5/e6,e6/e7} \draw[fleche] (\a) -- (\b);
\draw[fleche, poporange] (c5) -- (e5);
\draw[fleche, poporange] (c6) -- (e6);
\draw[fleche, poporange] (c7) -- (e7);
\end{tikzpicture}
\end{center}
\legende{Organigramme de fabrication du jus (ou nectar) de mangue pasteurisé. Les étapes 1 à 7 (bleu) se succèdent sans interruption ; les paramètres critiques (orange) doivent être strictement respectés pour garantir la sécurité microbiologique du produit.}
\end{popfigure}

\begin{methode}[Produire artisanalement du jus de mangue pasteurisé en toute sécurité]
Protocole en 8 étapes, réalisable en cuisine familiale ou en club scientifique :
\begin{enumerate}
\item \textbf{Hygiène} : se laver soigneusement les mains au savon ; laver le plan de travail, les couteaux, le mixeur et les bouteilles à l'eau chaude savonneuse, puis rincer.
\item \textbf{Tri} : ne retenir que des mangues mûres, saines, sans blessure ni pourriture.
\item \textbf{Lavage} des fruits à l'eau potable propre ; égoutter.
\item \textbf{Épluchage et dénoyautage} : récupérer la pulpe dans un récipient propre.
\item \textbf{Broyage et raffinage} : mixer la pulpe, puis la passer au tamis fin pour éliminer les fibres.
\item \textbf{Formulation} : pour un nectar, diluer avec de l'eau potable (environ 1 volume d'eau pour 2 volumes de pulpe) et sucrer au goût ; quelques gouttes de jus de citron abaissent le pH.
\item \textbf{Pasteurisation} : porter le jus à \SI{85}{\celsius} et maintenir 5 à 10 minutes en remuant ; ne jamais faire bouillir violemment (perte d'arômes et de vitamine C).
\item \textbf{Conditionnement à chaud} : remplir immédiatement les bouteilles ébouillantées jusqu'à 1 cm du bord, boucher hermétiquement, retourner quelques secondes pour stériliser le bouchon, puis refroidir rapidement (bain d'eau froide) et étiqueter (date de fabrication).
\end{enumerate}
\end{methode}

\begin{attention}[Le couple temps--température]
La pasteurisation n'est efficace que si le \cle{couple temps--température} est respecté : une température trop basse ou un temps trop court laisse survivre des micro-organismes ; une température trop élevée ou trop longue détruit les arômes et la vitamine C. De même, un conditionnement à froid dans des récipients non ébouillantés recontamine le produit : la pasteurisation aurait été faite pour rien.
\end{attention}

\courssub{Fabrication de la purée et de la confiture de mangue}

\begin{definition}[Purée et confiture]
La \cle{purée} de mangue est la pulpe broyée et raffinée, éventuellement concentrée par évaporation d'une partie de son eau. La \cle{confiture} est obtenue par cuisson prolongée de la pulpe avec une forte proportion de \cle{sucre} (environ 50 à 60 g de sucre pour 100 g de pulpe) jusqu'à obtenir une consistance gélifiée ; on ajoute souvent de l'\cle{acide citrique} (jus de citron).
\end{definition}

Le principe de conservation repose sur deux mécanismes que tu dois savoir relier aux principes généraux de la section précédente :

\begin{itemize}
\item La \textbf{concentration par chauffage} évapore une partie de l'eau : l'activité de l'eau $a_w$ diminue, les micro-organismes ne peuvent plus se développer (dessiccation osmotique).
\item Le \textbf{sucre} n'est pas seulement un édulcorant : à forte concentration, il lie l'eau et réduit encore l'eau disponible. C'est un véritable \cle{conservateur} : plongées dans un milieu très sucré, les cellules microbiennes perdent leur eau par osmose et meurent ou restent dormantes.
\item L'\textbf{acide citrique} abaisse le \cle{pH} (vers 3--3,5) : or la plupart des bactéries pathogènes ne se développent pas en dessous de pH 4,5. L'acidité renforce donc la sécurité et améliore la prise de la gelée (activation de la pectine).
\end{itemize}

\begin{exemplebox}[Préparation d'une confiture de mangue]
Pour 1 kg de pulpe de mangue raffinée : ajouter 500 à 600 g de sucre et le jus de 2 citrons. Cuire à feu moyen en remuant régulièrement (pour éviter la caramélisation au fond) jusqu'à ce qu'une goutte déposée sur une assiette froide se fige (test de la goutte). Mettre en pots ébouillantés, à chaud, fermer hermétiquement et retourner les pots 5 minutes. Durée de conservation : plusieurs mois à température ambiante, à l'abri de la lumière.
\end{exemplebox}

\courssub{Fabrication de la mangue séchée}

\begin{definition}[Séchage]
Le \cle{séchage} est une conservation par déshydratation : on extrait l'eau du fruit jusqu'à ramener son humidité résiduelle en dessous de 15 \%, seuil au-dessous duquel moisissures, levures et bactéries ne peuvent plus se multiplier ($a_w < 0,6$).
\end{definition}

Le procédé comprend quatre étapes :

\begin{enumerate}
\item \textbf{Tranchage} : après lavage, épluchage et dénoyautage, la pulpe est coupée en tranches régulières de 5 à 8 mm d'épaisseur (des tranches trop épaisses sèchent mal et moisissent au cœur).
\item \textbf{Traitement anti-brunissement} : les tranches sont trempées dans du jus de citron dilué (l'acide citrique et la vitamine C inhibent les enzymes du brunissement, les polyphénoloxydases) ou \cle{blanchies} (plongées 1 à 2 minutes dans l'eau bouillante), ce qui détruit les enzymes responsables du brunissement enzymatique et réduit la flore microbienne de surface.
\item \textbf{Séchage} : au soleil (séchage solaire traditionnel, 2 à 4 jours selon l'ensoleillement) ou en \cle{séchoir} amélioré à 60--\SI{70}{\celsius} (6 à 12 heures), jusqu'à humidité inférieure à 15 \% : la tranche devient souple, non collante.
\item \textbf{Emballage hermétique} : sachets ou bocaux étanches, à l'abri de l'air, de l'humidité et des insectes.
\end{enumerate}

\begin{exemplebox}[Rendement du séchage]
La mangue fraîche contient environ 85 \% d'eau. Il faut environ \textbf{10 à 12 kg de mangues fraîches} (avec peau et noyau) pour obtenir \textbf{1 kg de mangue séchée}. Ce chiffre explique le prix plus élevé du produit séché, mais aussi sa valeur concentrée en sucres, en fibres et en carotènes (provitamine A).
\end{exemplebox}

\begin{popfigure}
\begin{center}
\begin{tikzpicture}[font=\small]
% Panneau gauche : séchage solaire traditionnel
\begin{scope}
\draw[fill=poppinkL, draw=poppink, rounded corners=4pt] (-4.6,-0.4) rectangle (-0.2,4.4);
\node[font=\small\bfseries, poppink!60!black] at (-2.4,4.0) {Séchage solaire traditionnel};
\draw[fill=popgold, draw=popgold!70!black] (-2.4,3.55) circle (0.28);
\foreach \a in {45,90,135,180,225,270,315,0} \draw[popgold!70!black] (-2.4,3.55) ++(\a:0.34) -- ++(\a:0.22);
\draw[fill=popcream, draw=popmuted] (-4.2,0.9) rectangle (-0.6,1.5);
\node[font=\footnotesize] at (-2.4,1.2) {Tranches sur claie / nappe au sol};
\draw[-{Stealth}, popmuted, dashed] (-3.6,3.1) -- (-3.6,1.7);
\draw[-{Stealth}, popmuted, dashed] (-1.2,3.1) -- (-1.2,1.7);
\node[font=\footnotesize, align=left, anchor=west] at (-4.4,0.55) {Durée : 2--4 jours ; exposé aux poussières,};
\node[font=\footnotesize, align=left, anchor=west] at (-4.4,0.15) {mouches, pluies ; qualité irrégulière.};
\end{scope}
% Panneau droit : séchoir amélioré
\begin{scope}
\draw[fill=popgreenL, draw=popgreen, rounded corners=4pt] (0.4,-0.4) rectangle (4.8,4.4);
\node[font=\small\bfseries, popgreen!50!black] at (2.6,4.0) {Séchoir amélioré (60--\SI{70}{\celsius})};
\draw[fill=popblueL, draw=popblue, rounded corners=2pt] (1.0,1.0) rectangle (4.2,3.3);
\foreach \y in {1.5,2.1,2.7} {\draw[popblue!60] (1.15,\y) -- (4.05,\y); \foreach \x in {1.6,2.4,3.2,3.8} \draw[fill=poporange, draw=poporange!70!black] (\x,\y) ellipse (0.16 and 0.05);}
\draw[-{Stealth}, thick, poppink] (2.6,0.75) -- (2.6,1.15);
\node[font=\footnotesize, poppink!70!black] at (2.6,0.5) {air chaud};
\draw[-{Stealth}, thick, popblue] (4.2,3.0) -- (4.6,3.35);
\node[font=\footnotesize, popblue!70!black, align=center] at (4.55,2.6) {air\\humide};
\node[font=\footnotesize, align=left, anchor=west] at (0.6,0.55) {Durée : 6--12 h ; à l'abri des poussières};
\node[font=\footnotesize, align=left, anchor=west] at (0.6,0.15) {et insectes ; produit homogène et sain.};
\end{scope}
\end{tikzpicture}
\end{center}
\legende{Comparaison des deux modes de séchage de la mangue. À gauche : séchage solaire traditionnel, gratuit mais lent (2 à 4 jours) et exposé aux contaminations (poussières, mouches, pluie). À droite : séchoir amélioré à 60--\SI{70}{\celsius} : l'air chaud traverse les claies de tranches (orange) et l'air humide est évacué ; séchage rapide (6 à 12 h), hygiénique et régulier.}
\end{popfigure}

\courssub{Conservation des mangues fraîches}

On ne transforme pas toujours la mangue : il est aussi possible de retarder son altération pour la commercialiser fraîche.

\begin{definition}[Fruit climatérique]
La mangue est un \cle{fruit climatérique} : elle \textbf{poursuit son mûrissement après la récolte}, grâce à une production d'éthylène (hormone végétale gazeuse) qui accélère le ramollissement, le jaunissement et l'édulcoration. C'est pourquoi on la récolte à \cle{maturité physiologique} (fruit formé, encore vert-ferme) et non à pleine maturité gustative.
\end{definition}

\begin{propriete}[Conditions de conservation des mangues fraîches]
\begin{itemize}
\item \textbf{Récolte} à maturité physiologique, avec précaution (éviter les chocs et blessures, portes d'entrée des pourritures) ;
\item \textbf{Stockage au frais} entre 10 et \SI{13}{\celsius} : le froid ralentit la respiration du fruit et le développement microbien ; en dessous de \SI{10}{\celsius}, la mangue subit des « dommages de froid » (taches brunes, mûrissement anormal, perte de saveur) ;
\item \textbf{Durée limitée} : même au frais, l'entreposage ne doit pas dépasser 2 à 4 semaines ; à température ambiante (25--\SI{30}{\celsius}), la mangue mûrit en quelques jours et se gâte en une à deux semaines.
\end{itemize}
\end{propriete}

\begin{exoresolu}[Choisir la bonne stratégie de conservation]
Énoncé. Une coopérative de l'Adamaoua récolte 500 kg de mangues à maturité physiologique, mais la piste vers la ville est coupée pour 10 jours. Proposer une stratégie de conservation en justifiant chaque choix par les principes étudiés.
\tcblower
Solution. Le transport étant impossible, il faut empêcher le mûrissement et l'altération pendant 10 jours :
\begin{enumerate}
\item \textbf{Trier} immédiatement : éliminer tout fruit blessé ou piqué (foyers de pourriture qui contamineraient le lot par contact et par éthylène).
\item \textbf{Stocker au frais} entre 10 et \SI{13}{\celsius} si une chambre froide est disponible : le froid ralentit la respiration du fruit climatérique et retarde le mûrissement ; 10 jours restent compatibles avec la durée limite de 2 à 4 semaines.
\item À défaut de chambre froide, stocker à l'\textbf{ombre}, en local aéré, en \textbf{couches fines} (pas de tas profonds qui échauffent et écrasent les fruits), et envisager de \textbf{transformer} une partie de la récolte (séchage, confiture) pour réduire les pertes.
\end{enumerate}
Le raisonnement repose sur deux principes : la mangue est climatérique (le froid retarde son mûrissement) et les blessures favorisent les contaminations (d'où le tri et la manutention soigneuse).
\end{exoresolu}

\courssub{Contrôle qualité et enjeux}

Quel que soit le produit (jus, confiture, mangue séchée), la qualité et la sécurité reposent sur trois piliers :

\begin{itemize}
\item \textbf{L'hygiène} : lavage des mains, du matériel et des locaux ; eau potable ; récipients ébouillantés. La plupart des échecs de conservation artisanale viennent d'une recontamination après traitement thermique.
\item \textbf{Le respect des couples temps--température} : pasteurisation à \SI{85}{\celsius} pendant 5 à 10 min, séchage à 60--\SI{70}{\celsius} jusqu'à moins de 15 \% d'humidité, stockage des fruits frais à 10--\SI{13}{\celsius}.
\item \textbf{L'étiquetage} : chaque pot, sachet ou bouteille doit porter la \cle{date de fabrication} et la \cle{durée de conservation} (DLC), ainsi que la nature du produit. C'est une exigence de traçabilité et de confiance pour le consommateur.
\end{itemize}

\begin{aretenir}[L'essentiel sur la mangue]
\begin{itemize}
\item Matière première : mangues mûres, saines, sans blessure ; tri et lavage obligatoires ; un fruit pourri se jette (toxines thermostables).
\item Jus / nectar : épluchage, dénoyautage, broyage, raffinage, pasteurisation (\SI{85}{\celsius}, 5--10 min), conditionnement à chaud.
\item Confiture : concentration + sucre (réduction de l'eau disponible) + acide citrique (abaissement du pH).
\item Mangue séchée : tranchage, anti-brunissement (citron ou blanchiment), séchage à 60--\SI{70}{\celsius} jusqu'à $< 15\,\%$ d'humidité, emballage hermétique ; rendement : 10 à 12 kg de frais pour 1 kg de sec.
\item Fruit frais : récolte à maturité physiologique (fruit climatérique), stockage à 10--\SI{13}{\celsius}, durée limitée.
\end{itemize}
\end{aretenir}

\begin{savaistu}[La mangue, un trésor du nord Cameroun]
Les régions de l'\textbf{Adamaoua} et du \textbf{Nord} concentrent l'essentiel de la production nationale de mangues, avec des vergers qui s'étendent à perte de vue autour de Ngaoundéré, Garoua et Maroua. Pendant la haute saison, une grande partie de la récolte pourrissait faute de débouchés et de routes praticables. Depuis quelques années, des \textbf{unités artisanales et semi-industrielles de séchage} se développent : la mangue séchée camerounaise s'exporte désormais vers l'Europe, créant des revenus pour les productrices et producteurs ruraux. Maîtriser les techniques de cette leçon, c'est donc aussi se donner les moyens d'un véritable projet d'entrepreneuriat agricole.
\end{savaistu}

\begin{exercice}[À toi de jouer]
Une élève de Terminale C prépare de la mangue séchée au club de biotechnologie. Elle coupe des tranches de 2 cm d'épaisseur, les expose au soleil sur une nappe posée au sol pendant une journée, puis les emballe encore souples et humides dans un sachet fermé. Trois jours plus tard, des moisissures apparaissent. Identifier \textbf{quatre erreurs} commises et proposer pour chacune la correction adaptée, en justifiant par les principes de conservation.
\end{exercice}

\begin{corrige}[Exercice : mangue séchée moisie]
\begin{enumerate}
\item \textbf{Tranches trop épaisses} (2 cm) : l'eau du cœur ne s'évacue pas assez vite. Correction : trancher à 5--8 mm.
\item \textbf{Pas de traitement anti-brunissement ni de blanchiment} : la flore microbienne de surface n'est pas réduite et les enzymes de brunissement restent actives. Correction : tremper dans du jus de citron dilué ou blanchir 1--2 min.
\item \textbf{Séchage insuffisant} (une seule journée, au sol) : l'humidité résiduelle reste supérieure à 15 \%, seuil de développement des moisissures. Correction : sécher 2--4 jours au soleil sur claie surélevée, ou 6--12 h en séchoir à 60--\SI{70}{\celsius}, jusqu'à texture souple non collante.
\item \textbf{Nappe au sol} : contamination par les poussières, les insectes et les animaux. Correction : utiliser des claies surélevées couvertes d'un filet, ou un séchoir fermé.
\end{enumerate}
Le principe général : les moisissures ne se développent que si l'eau disponible reste suffisante ($a_w > 0,6$) ; tout le procédé vise à abaisser l'humidité sous 15 \% et à éviter les recontaminations.
\end{corrige}

\coursec{Transformation et conservation de la tomate}

Après la mangue, fruit sucré et peu acide, abordons maintenant le second fruit-vedette de nos marchés camerounais : la \cle{tomate}. Si tu te promènes sur le marché Mokolo à Yaoundé ou au marché central de Douala en pleine saison sèche, tu verras des montagnes de tomates vendues à bas prix, dont une partie finit malheureusement à la poubelle. Quelques mois plus tard, la même cagette coûte trois à quatre fois plus cher. Comprendre comment transformer et conserver la tomate, c'est donc répondre à un véritable problème économique et alimentaire national. La démarche est la même que pour la mangue, mais tu vas voir que la tomate possède un atout majeur que la mangue n'a pas : son \cle{acidité naturelle}.

\courssub{Particularités de la tomate : un fruit fragile mais acide}

\begin{experience}[Observation : pourquoi la tomate s'abîme-t-elle si vite ?]
\textbf{Matériel :} tomates mûres intactes, tomates mûres légèrement écrasées, deux assiettes, thermomètre de cuisine.

\textbf{Protocole :} on dépose sur une assiette des tomates intactes et sur une autre des tomates écrasées. On laisse le tout à la température ambiante d'une cuisine (environ \SI{28}{\celsius}) pendant 4 à 5 jours. On observe chaque jour.

\textbf{Observations :} les tomates écrasées se couvrent de moisissures dès le 2\ieme{} ou 3\ieme{} jour et dégagent une odeur aigre ; les tomates intactes restent saines plus longtemps mais finissent par ramollir, se rider puis pourrir.

\textbf{Interprétation :} la tomate contient environ \cle{94 \% d'eau} : c'est un milieu idéal pour le développement des micro-organismes (moisissures, levures, bactéries). Sa \cle{peau très fine et fragile} est la seule barrière protectrice ; dès qu'elle est rompue (choc, écrasement, piqûre d'insecte), les microbes pénètrent et la pulpe, riche en eau et en nutriments, se dégrade rapidement. De plus, la tomate est un fruit \emph{climactérique} : elle continue à mûrir après la récolte grâce à sa respiration et à la production d'éthylène, ce qui accélère son ramollissement.
\end{experience}

La tomate est donc l'un des fruits les plus \cle{périssables} qui soient : à température ambiante tropicale, sa durée de conservation ne dépasse guère \textbf{3 à 7 jours}. Mais elle possède une propriété décisive pour sa transformation.

\begin{propriete}[Acidité naturelle de la tomate et traitement thermique]
Le jus de tomate a un pH compris entre \num{4,2} et \num{4,5}, donc \textbf{inférieur à 4,6}. Or, comme l'a montré l'expérience de la section 2 (germination de spores de \emph{Clostridium} en milieux de pH différents), la bactérie pathogène la plus redoutée en conserve — \emph{Clostridium botulinum}, agent du botulisme — \textbf{ne peut ni se multiplier ni produire de toxine à pH $< 4,6$}. L'industrie fixe souvent une marge de sécurité à pH 4,5.

\textbf{Conséquence pratique :} la tomate est un aliment \cle{acide} ; une simple \cle{pasteurisation} ou un traitement thermique à \SI{100}{\celsius} (ébullition à l'air libre) suffit à garantir une conservation sûre, sans recourir à la stérilisation à \SI{115}{\celsius}--\SI{121}{\celsius} en autoclave, obligatoire pour les aliments peu acides (haricots, maïs, viandes).
\end{propriete}

\begin{aretenir}[Les deux visages de la tomate]
\begin{itemize}
\item \textbf{Fragilité :} 94 \% d'eau, peau fine, fruit climactérique $\Rightarrow$ conservation fraîche très courte.
\item \textbf{Acidité :} pH $\approx 4,2$--$4,5$ $\Rightarrow$ traitement thermique simple à \SI{100}{\celsius} suffisant pour les conserves.
\end{itemize}
C'est cette acidité qui explique le succès mondial des conserves de tomate, du concentré au jus.
\end{aretenir}

\courssub{Fabrication du concentré de tomate}

\begin{definition}[Concentré de tomate et extrait sec]
Le \cle{concentré de tomate} est le produit obtenu par élimination partielle de l'eau d'un jus de tomate raffiné (débarrassé des peaux et des graines), jusqu'à atteindre une teneur minimale en matière sèche imposée par les normes (Codex Alimentarius).

L'\cle{extrait sec} (ou matière sèche soluble) est la fraction du produit qui subsiste après évaporation complète de l'eau : sucres, acides organiques, sels minéraux, pigments (lycopène), protéines. On l'exprime en pourcentage de la masse totale, souvent mesuré au réfractomètre en degrés Brix ($^\circ$Bx).
\end{definition}

Selon la teneur en extrait sec, on distingue usuellement (normes Codex / industrielles) :
\begin{itemize}
\item le \textbf{simple concentré} (tomate simple concentrée) : 16 à 20 \% d'extrait sec (souvent 18--20 \%) ;
\item le \textbf{double concentré} : 28 à 30 \% d'extrait sec (le plus courant dans le commerce camerounais, en boîtes de 70 g ou 400 g) ;
\item le \textbf{triple concentré} : 36 à 38 \% d'extrait sec.
\end{itemize}

\begin{methode}[Chaîne de fabrication du concentré de tomate (procédé \emph{Hot Break})]
\begin{enumerate}
\item \textbf{Lavage :} les tomates sont lavées à grande eau (par trempage puis aspersion) pour éliminer terre, poussières et résidus de pesticides.
\item \textbf{Tri et épluchage préalable :} on élimine les fruits pourris, verts ou écrasés, ainsi que les pédoncules. \emph{Point critique : un seul fruit pourri peut contaminer tout un lot par ses enzymes et mycotoxines.}
\item \textbf{Broyage :} les fruits sont écrasés en pulpe grossière.
\item \textbf{Chauffage rapide (\emph{Hot Break}) :} la pulpe est portée rapidement à \SI{85}{\celsius}--\SI{95}{\celsius}. Ce chauffage a deux rôles : rompre les tissus pour libérer le jus et \textbf{inactiver les enzymes} (notamment les pectinestérases et polygalacturonases) qui, sinon, dépolymériseraient la pectine et liquéfieraient le produit (perte de consistance).
\item \textbf{Raffinage :} la pulpe chaude traverse des tamis fins (raffineuse) qui retiennent les \textbf{peaux et les graines} ; on obtient un jus lisse.
\item \textbf{Concentration par évaporation :} le jus est chauffé dans des évaporateurs (souvent à effets multiples sous vide, ce qui permet d'évaporer l'eau vers \SI{60}{\celsius}--\SI{70}{\celsius} sans brûler le produit ni dégrader la couleur -- lycopène -- et le goût) jusqu'à la teneur en extrait sec visée.
\item \textbf{Conditionnement à chaud :} le concentré est rempli à plus de \SI{85}{\celsius} dans des boîtes ou sachets préalablement nettoyés, puis fermés hermétiquement.
\item \textbf{Pasteurisation :} les emballages fermés sont chauffés à \SI{90}{\celsius}--\SI{100}{\celsius} pendant 20 à 40 min selon le format, puis refroidis rapidement. Grâce au pH $< 4,6$, ce traitement suffit à détruire les micro-organismes d'altération (levures, moisissures, bactéries lactiques).
\end{enumerate}
\end{methode}

\begin{popfigure}
\begin{center}
\begin{tikzpicture}[
font=\small,
etape/.style={rectangle, rounded corners=2pt, draw=popblue, fill=popblueL, text width=3.4cm, align=center, minimum height=0.9cm},
crit/.style={rectangle, draw=poporange, fill=poporangeL, text width=3.1cm, align=center, minimum height=0.8cm, font=\footnotesize},
fleche/.style={-{Stealth[length=2.5mm]}, thick, popdark}
]
\node[etape] (lav) at (0,0) {1. Lavage};
\node[etape, right=0.55cm of lav] (tri) {2. Tri / Épluchage};
\node[etape, right=0.55cm of tri] (broy) {3. Broyage};
\node[etape, right=0.55cm of broy] (chauf) {4. Chauffage (Hot Break)};
\node[etape, below=1.5cm of chauf] (raf) {5. Raffinage};
\node[etape, left=0.55cm of raf] (conc) {6. Concentration (sous vide)};
\node[etape, left=0.55cm of conc] (cond) {7. Conditionnement à chaud};
\node[etape, left=0.55cm of cond] (past) {8. Pasteurisation};
\draw[fleche] (lav) -- (tri);
\draw[fleche] (tri) -- (broy);
\draw[fleche] (broy) -- (chauf);
\draw[fleche] (chauf) -- (raf);
\draw[fleche] (raf) -- (conc);
\draw[fleche] (conc) -- (cond);
\draw[fleche] (cond) -- (past);
\node[crit, above=0.35cm of chauf] (c1) {PCC : \SI{85}{\celsius}--\SI{95}{\celsius} (inactivation enzymatique)};
\draw[-{Stealth[length=2mm]}, poporange, thick] (c1) -- (chauf);
\node[crit, below=0.35cm of past] (c2) {PCC : \SI{90}{\celsius}--\SI{100}{\celsius}, 20--40 min (pH $< 4,6$)};
\draw[-{Stealth[length=2mm]}, poporange, thick] (c2) -- (past);
\node[crit, above=0.35cm of tri] (c3) {PCC : éliminer tout fruit pourri};
\draw[-{Stealth[length=2mm]}, poporange, thick] (c3) -- (tri);
\end{tikzpicture}
\end{center}
\legende{Organigramme de fabrication du concentré de tomate (procédé Hot Break). Les encadrés orange indiquent les \textbf{points critiques de contrôle} (PCC) : étapes où une température ou une opération mal maîtrisée compromettrait la sécurité ou la qualité du produit final.}
\end{popfigure}

\begin{savaistu}[Hot Break vs Cold Break]
Le \emph{Hot Break} (85--95 °C) inactiver les pectinases $\rightarrow$ concentré \textbf{visqueux}, stable, couleur rouge vif (lycopène préservé). Le \emph{Cold Break} (65--75 °C) préserve l'activité pectinase $\rightarrow$ jus plus fluide, mais risque de séparation de phase (sérum) et couleur plus terne. L'industrie du concentré utilise presque exclusivement le Hot Break.
\end{savaistu}

\begin{exemplebox}[Rendement de concentration]
La tomate fraîche contenant environ 94 \% d'eau (soit 6 \% d'extrait sec), il faut \textbf{5 à 6 kg de tomates fraîches} pour produire \textbf{1 kg de concentré à 28--30 \% d'extrait sec}. Ce chiffre explique à la fois le prix du concentré et l'intérêt de le fabriquer sur place, là où la matière première est abondante et bon marché en saison.
\end{exemplebox}

\begin{exoresolu}[Calcul d'un rendement de concentration]
\textbf{Énoncé.} Une coopérative de femmes de Bafia dispose de \SI{120}{kg} de tomates fraîches contenant 5,5 \% d'extrait sec. Elle veut fabriquer un concentré à 28 \% d'extrait sec. Quelle masse de concentré obtiendra-t-elle, et quelle masse d'eau devra-t-elle évaporer ?

\tcblower
\textbf{Solution.} L'extrait sec se conserve au cours de l'évaporation (seule l'eau s'en va).

Masse d'extrait sec disponible :
\[ m_{ES} = 120 \times \frac{5,5}{100} = \SI{6,6}{kg}. \]

Masse de concentré à 28 \% contenant ces \SI{6,6}{kg} d'extrait sec :
\[ m_{conc} = \frac{6,6}{0,28} \approx \SI{23,6}{kg}. \]

Masse d'eau évaporée :
\[ m_{eau} = 120 - 23,6 = \SI{96,4}{kg}. \]

\textbf{Conclusion :} la coopérative obtiendra environ \SI{23,6}{kg} de concentré (soit un rendement massique d'environ 20 \%) après évaporation d'environ \SI{96}{kg} d'eau. On vérifie le rapport : $120 / 23,6 \approx 5,1$ kg de tomates par kg de concentré, conforme au rendement usuel de 5 à 6.
\end{exoresolu}

\courssub{Purée, jus et sauce de tomate : des produits voisins mais distincts}

Ces trois produits partagent le début de la chaîne (lavage, tri, broyage, chauffage, raffinage) mais se distinguent par la \textbf{concentration finale} (teneur en extrait sec) et les \textbf{ingrédients ajoutés}.

\begin{center}
\begin{tabularx}{\linewidth}{|l|X|X|X|}
\hline
\rowcolor{popblueL} \textbf{Produit} & \textbf{Extrait sec (Codex / Usuel)} & \textbf{Ingrédients ajoutés} & \textbf{Utilisation} \\
\hline
Jus de tomate & 5 à 8 \% ($^\circ$Bx) & sel éventuel (0,5 à 1 \%) & boisson, base de sauces liquides \\
\hline
Purée de tomate & 8 à 24 \% (souvent 10--12 \% pour purée simple) & aucun ou peu (sel) & coulis, base de cuisine \\
\hline
Concentré (pâte) & $\ge$ 18 \% (simple), 28--30 \% (double), 36--38 \% (triple) & sel éventuel & assaisonnement concentré \\
\hline
Sauce de tomate & 12 à 25 \% & huile, oignon, ail, épices, sucre, vinaigre, herbes & produit prêt à l'emploi \\
\hline
\end{tabularx}
\end{center}

\begin{itemize}
\item Le \cle{jus de tomate} est le produit le moins concentré : c'est pratiquement le jus raffiné, légèrement salé, conditionné puis pasteurisé à \SI{90}{\celsius}--\SI{95}{\celsius}.
\item La \cle{purée} (ou coulis) est un jus partiellement concentré, sans peaux ni graines, de consistance onctueuse. Au-delà de 18--24 \% d'extrait sec, on bascule dans l'appellation « concentré ».
\item La \cle{sauce de tomate} est une préparation \emph{culinaire} : on y ajoute des ingrédients (huile, oignons, épices) et parfois des acidifiants autorisés (acide citrique, vinaigre) ; l'ajout de vinaigre abaisse encore le pH, ce qui renforce la sécurité microbiologique.
\end{itemize}

\begin{attention}[Ne pas confondre purée et concentré]
Au Baccalauréat comme à l'usine, le critère de distinction n'est \textbf{ni la couleur ni le goût}, mais la \textbf{teneur en extrait sec}. Une purée épaisse reste une purée si son extrait sec est inférieur à 18 \%. Retiens les seuils : jus $<$ purée $<$ concentré, dans l'ordre croissant de concentration.
\end{attention}

\courssub{Le séchage de la tomate}

Le séchage est la technique de conservation la plus ancienne et la plus accessible aux producteurs camerounais : il consiste à abaisser la teneur en eau du fruit jusqu'à \textbf{15--20 \%} au maximum, de sorte que l'\cle{activité de l'eau} ($a_w$) devienne inférieure à \textbf{0,6--0,7}, seuil en dessous duquel la plupart des moisissures, levures et bactéries ne peuvent plus se développer.

\begin{definition}[Activité de l'eau ($a_w$)]
L'activité de l'eau ($a_w$) est le rapport entre la pression de vapeur d'eau du produit et celle de l'eau pure à la même température. Elle représente l'eau \emph{libre}, disponible pour les réactions biochimiques et la croissance microbienne. $a_w = 1$ pour l'eau pure ; $a_w < 0,6$ assure une stabilité microbiologique longue.
\end{definition}

\begin{methode}[Fabrication de tomates séchées]
\begin{enumerate}
\item \textbf{Choisir} des tomates bien mûres, saines et fermes (les variétés allongées, type « Roma » ou « Plum », charnues et peu juteuses, conviennent le mieux car rapport chair/eau élevé).
\item \textbf{Laver} puis \textbf{couper} les fruits en deux ou en quatre dans le sens de la longueur ; on peut retirer les graines pour accélérer le séchage.
\item \textbf{Salage éventuel} léger : le sel attire l'eau par osmose (effet « saumure » de surface) et améliore la conservation en abaissant localement $a_w$.
\item \textbf{Séchage} :
\begin{itemize}
\item \emph{au soleil} : sur claies ou nattes propres, surélevées, couvertes d'un filet anti-insectes ; compter \textbf{3 à 5 jours} de plein soleil en saison sèche, en rentrant les claies la nuit pour éviter la réhumidification (rosée) ;
\item \emph{en séchoir} (solaire amélioré à convection forcée ou électrique) : à \SI{55}{\celsius}--\SI{65}{\celsius} pendant 8 à 15 h ; le séchage est plus rapide, plus hygiénique (pas de poussières, pas d'insectes) et indépendant de la météo.
\end{itemize}
\item \textbf{Conditionnement} : les tomates bien sèches (cassantes ou souples mais sans eau exsudée au pressage) sont mises en \textbf{sachets hermétiques} (type ZIP ou soudés) ou en bocaux, à l'abri de l'air, de l'humidité et de la lumière. Une variante appréciée : la conservation \textbf{dans l'huile} (avec ail et herbes), l'huile formant une barrière physique contre l'air et l'oxydation.
\end{enumerate}
\end{methode}

\begin{attention}[Un séchage incomplet est un danger sanitaire]
Si la tomate reste trop humide au cœur ($a_w > 0,7$), des moisissures se développent \emph{à l'intérieur} du sachet, parfois invisibles de l'extérieur. Certaines moisissures (\emph{Aspergillus flavus}, \emph{Penicillium}) produisent des \cle{mycotoxines} (aflatoxines, patuline) thermostables et dangereuses pour le foie (cancérigènes, mutagènes). Règle pratique : la tomate séchée ne doit plus laisser suinter de jus lorsqu'on la presse fermement entre les doigts.
\end{attention}

\courssub{Conservation des tomates fraîches}

Quand on ne peut pas transformer, on peut au moins \emph{ralentir} l'altération des fruits frais. Trois règles d'or basées sur la physiologie du fruit climactérique :

\begin{enumerate}
\item \textbf{Récolter au stade physiologique adapté.} Pour une consommation locale immédiate, on récolte à pleine maturité (fruit rouge, optimal gustatif). Pour le transport vers une ville lointaine ou un stockage de quelques jours, on récolte au stade « \emph{breaker} » ou « \emph{turning} » (début de coloration jaune-orangé sur 10--30 \% de la surface) : le fruit, encore ferme, supporte mieux les chocs mécaniques et achèvera sa maturation à destination grâce à sa production d'éthylène endogène.
\item \textbf{Stocker à température contrôlée.} La température optimale de stockage de la tomate mûre est de \SI{10}{\celsius} à \SI{12}{\celsius}, avec une humidité relative de 85 à 90 \% : on gagne ainsi 2 à 3 semaines de conservation en ralentissant la respiration et la sensibilité à l'éthylène. Pour les fruits verts/turning, on peut monter à \SI{12}{\celsius}--\SI{15}{\celsius}.
\item \textbf{Trier régulièrement.} Un fruit pourri libère de l'éthylène (accélérateur de maturation des voisins) et des spores de moisissures qui contaminent tout le lot : le tri périodique (tous les 2 à 3 jours) est indispensable dans les magasins et les marchés.
\end{enumerate}

\begin{attention}[Erreur fréquente : la tomate au réfrigérateur domestique]
Beaucoup de ménages rangent les tomates dans le réfrigérateur à \SI{4}{\celsius}--\SI{6}{\celsius}. C'est une erreur physiologique majeure ! En dessous de \SI{8}{\celsius} environ, la tomate, fruit tropical d'origine, subit des \cle{dommages du froid} (\emph{chilling injury}) : ses membranes cellulaires (riches en lipides insaturés) perdent leur fluidité et se dégradent, sa chair devient \textbf{farineuse} (méiolose), sa peau se tache (pitting), et surtout elle \textbf{perd son arôme} car la synthèse des composés volatils (C6 aldéhydes, alcools) est bloquée. Conserve tes tomates à l'ombre, dans un endroit aéré, ou entre 10 et \SI{12}{\celsius} si tu disposes d'une chambre froide réglable.
\end{attention}

\courssub{Le problème camerounais : gluts et flambée des prix}

La production camerounaise de tomate est fortement \textbf{saisonnière} : l'essentiel de la récolte arrive en saison sèche (décembre à mars), période favorable à la culture (peu de maladies fongiques, fort ensoleillement). Il en résulte un phénomène bien connu des économistes : le \cle{glut} (excédent ponctuel d'offre). L'offre explose, les prix s'effondrent, et faute de moyens de conservation et de transformation, \textbf{30 à 50 \% de la récolte} peut être perdue (pourriture, gaspillage). Quelques mois plus tard, en saison des pluies (juin--août), la tomate se raréfie (maladies, moins de soleil) et les prix flambent.

\begin{popfigure}
\begin{center}
\begin{tikzpicture}
\begin{axis}[
ybar, bar width=14pt,
width=0.9\linewidth, height=6.5cm,
xlabel={Mois}, ylabel={Prix moyen (FCFA/kg)},
symbolic x coords={Jan,Fév,Mar,Avr,Mai,Juin,Juil,Août,Sep,Oct,Nov,Déc},
xtick=data,
x tick label style={rotate=45, anchor=east, font=\footnotesize},
ymin=0, ymax=1000,
nodes near coords, every node near coord/.append style={font=\tiny, rotate=90, anchor=west, yshift=2pt},
axis lines=left,
ymajorgrids=true, grid style={popline!40},
enlarge x limits=0.15,
]
\addplot[fill=popblue, draw=popblue!70!black] coordinates {
(Jan,250) (Fév,200) (Mar,220) (Avr,350) (Mai,500) (Juin,650)
(Juil,800) (Août,850) (Sep,750) (Oct,600) (Nov,450) (Déc,300)};
\end{axis}
\end{tikzpicture}
\end{center}
\legende{Variation saisonnière indicative du prix de détail de la tomate sur un marché urbain camerounais : effondrement des prix en saison d'abondance (janvier--mars, période de \emph{glut}) et flambée en saison de pénurie (juin--août), où le kilogramme peut coûter 3 à 4 fois plus cher. Données moyennes observées (ex. marché Mfoundi, Yaoundé).}
\end{popfigure}

\begin{savaistu}[Des unités de transformation locales : une solution au gaspillage et à la souveraineté alimentaire]
Installer des \textbf{unités artisanales ou semi-industrielles de transformation} (concentré, purée, séchage) près des bassins de production — dans la plaine des Mbo, à Bafia, à Ngaoundéré ou dans le Mayo-Danay — permettrait de :
\begin{itemize}
\item \textbf{absorber les gluts} en achetant les tomates excédentaires aux producteurs à prix rémunérateur (évitant la vente à perte ou le jet) ;
\item \textbf{stabiliser l'offre} toute l'année et lisser les prix pour les consommateurs urbains ;
\item \textbf{créer des emplois ruraux}, notamment pour les femmes et les jeunes, et réduire la facture d'importation de concentré de tomate (le Cameroun importe encore des milliers de tonnes de double concentré par an).
\end{itemize}
C'est exactement la logique de la biotechnologie appliquée au développement : transformer localement ce que le pays produit, au lieu de le laisser pourrir ou de l'importer transformé.
\end{savaistu}

\begin{exercice}[Pour t'entraîner : bilan matière et valeur ajoutée]
Une unité de transformation de Ngaoundéré achète \SI{2}{t} de tomates à \SI{150}{FCFA/kg} en pleine saison (janvier). Elle produit une purée à 12 \% d'extrait sec à partir de tomates contenant 5,5 \% d'extrait sec. Les pertes de transformation (peaux, graines, résidus) représentent 5 \% de la masse initiale.
\begin{enumerate}
\item Calcule la masse de purée obtenue.
\item La purée est conditionnée en bocaux de 500 g vendus \SI{400}{FCFA} l'unité. Calcule le chiffre d'affaires prévisionnel.
\item En déduisant le coût de la matière première, estime la marge brute disponible pour couvrir énergie, main-d'œuvre, emballages, transport. Commente l'intérêt économique.
\end{enumerate}
\end{exercice}

\begin{corrige}[Exercice : unité de transformation de Ngaoundéré]
\begin{enumerate}
\item Masse initiale utile après pertes (peaux/graines) : $2000 \times (1 - 0,05) = \SI{1900}{kg}$ de pulpe/jus entrant en concentration.
Masse d'extrait sec disponible : $1900 \times 0,055 = \SI{104,5}{kg}$.
Masse de purée à 12 \% : $m_{puree} = \frac{104,5}{0,12} \approx \SI{871}{kg}$.
\item Nombre de bocaux de 500 g : $\frac{871}{0,5} = 1742$ bocaux.
Chiffre d'affaires : $1742 \times 400 = \SI{696800}{FCFA}$.
\item Coût matière première : $2000 \times 150 = \SI{300000}{FCFA}$.
Marge brute : $696800 - 300000 = \SI{396800}{FCFA}$.
\textbf{Commentaire :} Cette marge brute doit couvrir le bois/énergie (concentration), les bocaux/couvercles, la main-d'œuvre, le transport, l'amortissement du matériel. Même avec des charges estimées à \SI{200000}{FCFA}, le bénéfice net reste attractif (\textasciitilde \SI{200000}{FCFA} pour un lot de 2 t), tout en sauvant de la perte une récolte qui, non transformée, aurait en partie pourri. La transformation crée de la \textbf{valeur ajoutée} locale.
\end{enumerate}
\end{corrige}

\begin{aretenir}[L'essentiel sur la tomate]
\begin{itemize}
\item La tomate est très périssable (94 \% d'eau, peau fragile, climactérique) mais \textbf{acide} (pH $\approx 4,2$--$4,5 < 4,6$) : un traitement thermique à \SI{100}{\celsius} (pasteurisation) suffit pour la conserver sûrement (pas d'autoclave nécessaire).
\item Le concentré s'obtient par : lavage $\to$ tri $\to$ broyage $\to$ \textbf{Hot Break} (85--95 °C, inactivation pectinases) $\to$ raffinage $\to$ concentration sous vide $\to$ conditionnement à chaud $\to$ pasteurisation ; il faut 5 à 6 kg de tomates pour 1 kg de double concentré (28--30 \% ES).
\item Jus ($<$ 8 \% ES), purée (8--24 \% ES), concentré ($>$ 18 \% ES) et sauce (culinaire) se distinguent par la teneur en extrait sec et les ingrédients ajoutés.
\item Le séchage (soleil ou séchoir 55--65 °C) abaisse $a_w < 0,6$ ; danger mycotoxines si séchage incomplet.
\item Stockage frais optimal : \SI{10}{\celsius}--\SI{12}{\celsius} (jamais $< \SI{8}{\celsius}$ : dommages du froid, chair farineuse, perte arôme).
\item La transformation locale (concentré, purée, séché) est la réponse technique et économique aux pertes post-récolte (30--50 \%) et à la flambée des prix hors saison au Cameroun.
\end{itemize}
\end{aretenir}

\coursec{Pratique, qualité et enjeux socio-économiques}

Dans la section précédente, tu as appris à transformer la tomate en concentré, en purée et en jus, en appliquant les principes généraux de la conservation (destruction des micro-organismes et des enzymes par la chaleur, réduction de l'eau disponible, protection contre l'oxygène). Mais un savoir biotechnologique ne vaut vraiment que s'il est \emph{pratiqué} avec rigueur et s'il produit un bien \emph{de qualité}, vendable et bénéfique à la communauté. Cette dernière section te conduit du laboratoire scolaire à l'atelier de transformation : tu apprendras à conduire une séance pratique en toute sécurité, à juger de la qualité d'un produit fini, puis tu mesureras les enjeux économiques et nutritionnels de la transformation des fruits de saison au Cameroun.

\courssub{Mise en pratique encadrée : réaliser un produit simple}

\begin{experience}[TP : préparation d'un jus de mangue pasteurisé en classe ou en club scientifique]
\textbf{Matériel} : mangues mûres saines (environ \SI{2}{kg}), couteaux et planches lavés, mixeur ou moulin, tamis propre (toile fine ou passoire inox), marmite inox, thermomètre alimentaire (sonde), bouteilles en verre préalablement bouillies \SI{10}{min} et égouttées, entonnoir inox, eau potable, sucre (facultatif, \SIrange{50}{80}{g/L} selon acidité), balance de cuisine, gants propres, torchons propres, maniques.

\textbf{Protocole} :
\begin{enumerate}
\item \textbf{Hygiène préalable} : se laver les mains au savon (procédure chirurgicale : paumes, dos, interstices, pouces, ongles) ; laver les mangues à l'eau potable en les frottant ; ébouillanter le matériel de conditionnement (bouteilles, bouchons, entonnoir, louche) pendant \SI{10}{min} et laisser égoutter sur support propre.
\item \textbf{Préparation de la pulpe} : éplucher les fruits au couteau propre, détailler la chair en écartant le noyau, mixer avec un peu d'eau potable (environ \SI{0,5}{L} pour \SI{2}{kg} de pulpe) puis tamiser pour éliminer les fibres.
\item \textbf{Pasteurisation} : porter le jus à \SIrange{80}{85}{\celsius} et maintenir cette température pendant \SIrange{10}{15}{min} en remuant doucement (contrôle au thermomètre). Cette température assure la destruction des levures, moisissures et bactéries végétatives tout en limitant la dégradation de la vitamine C et des arômes.
\item \textbf{Conditionnement à chaud} : remplir les bouteilles encore chaudes (jus à \textgreater{} \SI{75}{\celsius}) à l'aide de l'entonnoir, en laissant un espace de tête d'environ \SI{1}{cm} ; boucher immédiatement et hermétiquement ; retourner les bouteilles \SI{5}{min} pour stériliser le bouchon par le jus chaud ; remettre à l'endroit et laisser refroidir rapidement à l'air libre (ou bain d'eau froide pour accélérer).
\item \textbf{Étiquetage et stockage} : étiqueter chaque bouteille (nom du produit, date de fabrication, liste ingrédients, conditions de conservation) ; stocker à l'ombre, au frais (\textless{} \SI{20}{\celsius}) et au sec.
\end{enumerate}

\textbf{Observations attendues} : jus homogène, couleur orange vif, odeur fruitée caractéristique ; après refroidissement, le couvercle est légèrement concave (dépression) signe du vide partiel. Après quelques jours, un jus correctement pasteurisé et embouteillé ne présente ni bulles, ni odeur aigre/alcoolique, ni gonflement du bouchon, ni dépôt anormal.

\textbf{Interprétation} : le chauffage à \SI{80}{\celsius} détruit la majorité des micro-organismes altérants (levures, moisissures, bactéries lactiques) et inactive les enzymes (polyphénoloxydases, pectinestérases) responsables du brunissement enzymatique et de la dépectinisation ; le remplissage à chaud et le bouchage immédiat créent une atmosphère pauvre en oxygène (vide partiel au refroidissement) et limitent la recontamination. La même démarche rigoureuse s'applique à une confiture (cuisson avec environ \SI{50}{\percent} de sucre en masse par rapport à la pulpe, qui abaisse l'activité de l'eau $a_w$ \textless{} 0,75) ou à des tranches séchées (déshydratation solaire sous moustiquaire ou en séchoir amélioré à \SIrange{50}{60}{\celsius} jusqu'à $a_w$ \textless{} 0,60).
\end{experience}

\begin{methode}[Conduire une séance pratique de transformation fruitière en toute sécurité : check-list hygiène]
\begin{enumerate}
\item \textbf{Avant la séance} : vérifier la disponibilité d'eau potable, de savon (liquide de préférence), de papier essuie-mains à usage unique et d'une source de chaleur stable ; trier rigoureusement les fruits (éliminer tout fruit pourri, moisi, meurtri ou atteint de mouches) ; ébouillanter bouteilles, bocaux, couvercles et ustensiles (louches, entonnoirs, cuillères) \SI{10}{min} ; préparer les étiquettes.
\item \textbf{Pendant la séance} : mains lavées au savon avant chaque manipulation et après tout contact avec une surface souillée (visage, cheveux, sol, déchets) ; cheveux attachés et couverts (charlotte ou foulard), tablier propre (coton ou jetable) ; ne jamais goûter avec la cuillère de service (utiliser une cuillère à usage unique) ; respecter scrupuleusement les couples température--durée validés (pasteurisation, cuisson confiture à \SI{103}{\celsius} \textasciitilde \SI{105}{\celsius} pour le point de gélification) ; éviter tout contact du produit chaud avec des surfaces non désinfectées (plan de travail, peau).
\item \textbf{Conditionnement} : récipients propres, secs (pas d'eau résiduelle qui diluerait le produit) et hermétiques ; remplissage à chaud pour les jus et confitures (température produit \textgreater{} \SI{75}{\celsius}) ; fermeture immédiate ; refroidissement rapide (bain d'eau froide ou étalement) pour limiter la cuisson résiduelle et fixer la couleur.
\item \textbf{Après la séance} : lavage, rinçage, séchage et rangement du matériel ; étiquetage obligatoire de chaque lot (nom du produit, date de fabrication, n° de lot, liste des ingrédients, DLUO indicative) ; suivi des lots pendant \SI{15}{jours} (test de stabilité : absence de fermentation, de moisissure, de bombage).
\item \textbf{Sécurité} : manipulation des liquides chauds et marmites avec des maniques sèches ; interdiction de courir ou de plaisanter près des sources de chaleur ; présence obligatoire de l'enseignant/encadrant pendant toute la manipulation ; extincteur ou couverture anti-feu accessible.
\end{enumerate}
\end{methode}

\begin{attention}[Un bocal bombé ou un jus fermenté est impropre à la consommation -- Danger toxines]
Erreur fréquente et dangereuse : croire qu'il suffit de « rebouillir » un produit avarié pour le rendre consommable. \textbf{C'est faux.} Le bombage d'un couvercle ou d'une bouteille traduit une production de gaz (\ce{CO2}, \ce{H2}) par des micro-organismes (levures, bactéries) qui se sont multipliés dans le produit par défaut de pasteurisation ou de conditionnement. Certains champignons microscopiques (moisissures) produisent des \cle{mycotoxines} (ex. patuline dans les jus de fruits, ochratoxine) qui sont \textbf{thermostables} : l'ébullition prolongée ne les détruit pas. De même, une moisissure visible en surface a souvent diffusé ses toxines dans toute la masse du produit par diffusion liquide. Règle absolue, sans exception : tout produit présentant bombage, bulles anormales, odeur aigre ou alcoolique, moisissure (même en surface), couleur modifiée (brunissement anormal) doit être \textbf{jeté immédiatement}, sans dégustation préalable, et le lot entier mis en quarantaine pour analyse de la cause.
\end{attention}

\courssub{Évaluer la qualité d'un produit transformé}

Comment savoir, sans laboratoire sophistiqué, si une confiture ou un jus est « bon » et commercialisable ? On utilise deux familles de critères complémentaires, accessibles par l'observation sensorielle et le contrôle visuel.

\begin{definition}[Critères de qualité d'un produit transformé]
Les \cle{critères organoleptiques} sont les caractéristiques perçues par les sens : la \textbf{couleur} (vive, homogène, proche du fruit frais, sans brunissement excessif), l'\textbf{odeur} (fruitée, typique, sans note aigre, alcoolique, de brûlé ou de moisi), le \textbf{goût} (équilibre sucre--acidité, saveur fruitée persistante, absence d'amertume ou de piquant) et la \textbf{texture} (consistance adaptée : nappante pour un jus, gélifiée pour une confiture, souple pour un fruit séché ; absence de grumeaux, de synérèse ou de cristallisation excessive du sucre). Les \cle{critères de conservation (stabilité microbienne)} traduisent la salubrité du produit : absence de fermentation (pas de bulles, pas de gaz à l'ouverture, pas de trouble anormal), absence de moisissures en surface ou dans la masse, absence de bombage des emballages (couvercle concave = bon signe), intégrité de la fermeture hermétique (joint intact, pas de fuite).
\end{definition}

\begin{exemplebox}[Grille simple d'évaluation d'une confiture de mangue (type Bac pratique)]
\begin{center}
\begin{tabularx}{\linewidth}{|l|X|X|}
\hline
\rowcolor{popblueL} \textbf{Critère} & \textbf{Produit conforme (Qualité supérieure)} & \textbf{Signe de défaut (Non-conformité)} \\
\hline
Couleur & Orange à dorée, homogène, brillante & Brunissement terne, taches sombres (surcuisson/oxydation) \\
\hline
Odeur & Fruitée, mangue mûre, agréable & Aigre (acide lactique), alcoolisée (levures), fermentée, brûlée \\
\hline
Goût & Sucré-acidulé équilibré, saveur fruitée & Acre, piquant (alcool), amer (brûlé), plat (manque d'acidité) \\
\hline
Texture & Gélifiée, nappante, tient sur la cuillère (test de la goutte) & Liquide (manque de cuisson / pectine) ou dure/cristallisée (excès sucre/cuisson) \\
\hline
Emballage & Couvercle concave (vide), joint intact, étiquette lisible & Couvercle bombé (gaz), fuite, moisissure au col, étiquette absente/illlisible \\
\hline
\end{tabularx}
\end{center}
\end{exemplebox}

\begin{savaistu}[La vitamine C, la plus fragile des vitamines]
La vitamine C (acide ascorbique) est la vitamine la plus fragile : elle est détruite par la chaleur (thermolabile), par l'oxygène de l'air (oxydation enzymatique et non enzymatique) et par la lumière (photodégradation). Une cuisson prolongée (confiture) ou un séchage solaire lent (exposition directe UV + chaleur + air) peut en détruire \SIrange{50}{80}{\percent}. D'où l'intérêt majeur des traitements \emph{rapides} et \emph{protégés} : pasteurisation courte à température contrôlée (\SI{85}{\celsius}, \SI{1}{min} en continu industriel ; \SI{80}{\celsius}, \SI{10}{min} en artisanal), séchage en séchoir amélioré (air chaud forcé, \SI{50}{--}\SI{60}{\celsius}, à l'abri du soleil direct) plutôt qu'exposition prolongée au soleil sur nattes, stockage des produits finis à l'abri de la lumière (bouteilles teintées, cartons). C'est pourquoi les jus industriels « flash-pasteurisés » sont chauffés très fort (\SI{95}{--}\SI{100}{\celsius}) mais très brièvement (\SI{15}{--}\SI{30}{s}) : on détruit les microbes en épargnant au maximum les vitamines et les arômes volatils.
\end{savaistu}

\courssub{Complément utile au Bac : unités de transformation et étiquetage (contexte professionnel)}

\begin{definition}[Unités de transformation et étiquetage réglementaire]
Une \cle{unité de transformation artisanale} est une petite structure (familiale, groupement, coopérative) utilisant des équipements simples (marmites, mixeurs, séchoirs solaires, remplisseuses manuelles) et une main-d'œuvre peu mécanisée ; sa capacité va de quelques kilogrammes à quelques centaines de kilogrammes de fruits par jour. Une \cle{unité semi-industrielle} emploie des équipements motorisés (pulpeuses, désaérateurs, pasteurisateurs tubulaires, remplisseuses doseuses semi-automatiques, étiqueteuses) et traite de l'ordre de la tonne à plusieurs tonnes par jour, avec un contrôle qualité formalisé (plan HACCP simplifié, analyses microbiologiques régulières). L'\cle{étiquetage} d'un produit alimentaire préemballé (norme CODEX / norme camerounaise NC 100) doit mentionner \textbf{obligatoirement} : 1) la dénomination de vente du produit, 2) la liste des ingrédients par ordre décroissant de masse (additifs inclus avec nom/fonction/n° CE), 3) la quantité nette (masse ou volume), 4) la date de fabrication et la Date Limite d'Utilisation Optimale (DLUO) ou Date Limite de Consommation (DLC), 5) les conditions particulières de conservation (ex. « À conserver au frais après ouverture »), 6) le nom ou la raison sociale et l'adresse du fabricant/conditionneur, 7) le numéro de lot, 8) le pays d'origine. L'apposition du logo « Made in Cameroon » ou de la marque de certification (ANOR) est un gage de confiance pour le consommateur.
\end{definition}

\courssub{Enjeux économiques : de la récolte à la valeur ajoutée}

Chaque année, au Cameroun comme dans toute l'Afrique tropicale, une fraction considérable des fruits de saison (souvent estimée entre \SI{30}{\percent} et \SI{50}{\percent} pour les fruits très périssables comme la mangue, la papaye, l'ananas) est perdue faute de débouchés rapides et de moyens de conservation : les fruits pourrissent sous les arbres, sur les étals ou le long des routes en pleine saison. La transformation est la réponse biotechnologique majeure à ce gâchis post-récolte.

\begin{exemplebox}[La transformation multiplie la valeur de la récolte : étude de cas chiffrée]
Une coopérative du Littoral transformant \SI{1}{t} de mangues par jour en jus pasteurisé peut multiplier par 5 à 10 la valeur de la récolte brute : des mangues achetées \SI{25}{FCFA/kg} en pleine saison (soit \SI{25000}{FCFA} la tonne à la collecte) fournissent, après tri, épluchage, pulpage (rendement \textasciitilde 50\% jus/pulpe), environ \SI{500}{L} de jus prêt à boire, vendus en moyenne \SI{500}{FCFA} le litre (soit \SI{250}{FCFA} le demi-litre), générant un chiffre d'affaires de l'ordre de \SI{250000}{FCFA} par tonne de mangues brutes transformées, avant déduction des coûts variables (sucre \textasciitilde \SI{15}{kg}, bouteilles/bouchons \textasciitilde \SI{500}{unités}, énergie bois/gaz/électricité, main-d'œuvre, transport, amortissement matériel). La transformation crée donc de la \cle{valeur ajoutée} à chaque étape de la chaîne : tri/calibrage (sélection qualité), pulpage/raffinage (extraction), pasteurisation (stabilisation), conditionnement (protection, image), distribution (logistique).
\end{exemplebox}

Les enjeux économiques se déclinent à trois niveaux majeurs pour le développement local :
\begin{itemize}
\item \textbf{Réduction drastique des pertes post-récolte} : chaque tonne transformée est une tonne sauvée de la pourriture, ce qui sécurise le revenu du producteur même en période de surproduction (effet « tampon » de marché) et évite la chute des prix au champ.
\item \textbf{Création de revenus et d'emplois décents ruraux} : la transformation crée des emplois non délocalisables (cueillette sélective, tri, conduite d'atelier, maintenance, vente, gestion administrative) et des revenus pour les producteurs (prix garanti), les coopératives (marge de transformation) et les PME agroalimentaires locales.
\item \textbf{Substitution aux importations et sécurité alimentaire} : produire localement jus, nectars, confitures, purées, concentrés de tomate et fruits séchés réduit structurellement l'achat de produits importés (souvent plus chers, moins frais, moins adaptés aux goûts locaux), économise les devises étrangères et renforce la souveraineté et la sécurité alimentaire nationale.
\end{itemize}

\begin{popfigure}
\begin{center}
\begin{tikzpicture}[
font=\small,
box/.style={draw=popblue, fill=popblueL, rounded corners=2pt, align=center, minimum height=0.9cm, text width=2.35cm},
va/.style={draw=poporange, fill=poporangeL, rounded corners=2pt, align=center, font=\scriptsize, text width=2.1cm},
fleche/.style={-{Stealth[length=3mm]}, thick, popdark}
]
\node[box, align=center] (prod) at (0,0){\textbf{Producteur}\\Récolte mangues};
\node[box, align=center] (collecte) at (3.4,0){\textbf{Collecte \& Tri}\\Transport, calibrage};
\node[box, align=center] (transfo) at (6.8,0){\textbf{Transformation}\\Pulpage, pasteurisation};
\node[box, align=center] (cond) at (10.2,0){\textbf{Conditionnement}\\Embouteillage, étiquetage};
\node[box, align=center] (distrib) at (13.6,0){\textbf{Distribution}\\Grossistes, détaillants};
\node[box, fill=popgreenL, draw=popgreen, align=center] (conso) at (13.6,-2.2){\textbf{Consommateur}\\Jus dispo. toute l'année};

\draw[fleche] (prod) -- (collecte);
\draw[fleche] (collecte) -- (transfo);
\draw[fleche] (transfo) -- (cond);
\draw[fleche] (cond) -- (distrib);
\draw[fleche] (distrib) -- (conso);

\node[va, align=center] at (1.7,-1.3){+ Valeur : prix garanti,\\pertes réduites};
\node[va, align=center] at (5.1,-1.3){+ Valeur : sélection\\fruits mûrs sains};
\node[va, align=center] at (8.5,-1.3){+ Valeur : produit stable,\\emplois atelier};
\node[va, align=center] at (11.9,-1.3){+ Valeur : conservation,\\traçabilité, image};
\end{tikzpicture}
\end{center}
\legende{Chaîne de valeur de la mangue, du producteur au consommateur. À chaque étape (flèches noires), la transformation et les services associés créent de la valeur ajoutée (encadrés orange) : le fruit frais, très périssable et saisonnier, devient un produit stable, transportable, traçable et vendable toute l'année.}
\end{popfigure}

\courssub{Enjeux nutritionnels : des vitamines toute l'année, avec lucidité sur les pertes}

Les fruits de saison sont des sources majeures de micronutriments : provitamine A (bêta-carotène) et vitamine C pour la mangue ; vitamine C et lycopène (antioxydant) pour la tomate ; sans oublier minéraux (potassium, magnésium) et fibres alimentaires. Or leur disponibilité fraîche est brève : la saison des mangues ne dure que 3 à 4 mois (mars-juin selon les zones). La transformation (jus, nectars, confitures, purées, concentrés, séchage) \textbf{étale la disponibilité de ces nutriments essentiels sur toute l'année}, ce qui est précieux pour l'équilibre alimentaire des populations, notamment des enfants en bas âge et des femmes enceintes, en zone rurale comme urbaine, où les carences en vitamine A et en fer restent un problème de santé publique majeur au Cameroun.

Mais il faut être scientifiquement honnête sur les limites : la vitamine C, thermosensible et oxydable, est partiellement détruite par la cuisson (confiture, concentré) et le séchage lent ; les traitements thermiques intenses dégradent aussi une partie des caroténoïdes (isomérisation trans/cis, oxydation). Les produits transformés ne remplacent donc \textbf{jamais totalement} le fruit frais sur le plan nutritionnel brut. Les technologies modernes (pasteurisation flash, séchage par air chaud déshumidifié, surgélation, haute pression) limitent drastiquement ces pertes. En contexte artisanal, les traitements rapides (pasteurisation courte \SI{85}{\celsius}/\SI{10}{min}, séchage rapide \SI{50}{--}\SI{60}{\celsius} à l'abri du soleil direct, conditionnement opaque) sont les leviers accessibles pour préserver le maximum de potentiel nutritionnel. C'est un \emph{savoir admis} au programme : on te demande de connaître l'existence et l'ordre de grandeur de ces pertes, pas leurs mécanismes biochimiques détaillés.

\courssub{Bilan de la leçon : relier chaque cause d'altération à la technique qui la neutralise}

Toute la logique de la biotechnologie de la conservation tient en une idée directrice, cœur exigible de la leçon : \textbf{à chaque cause d'altération des fruits correspond une technique de conservation qui la neutralise spécifiquement}. Le tableau de synthèse suivant formalise cette correspondance fondamentale.

\begin{popfigure}
\begin{center}
\begin{tikzpicture}[
font=\small,
cause/.style={draw=poppink, fill=poppinkL, rounded corners=2pt, align=center, minimum height=1.1cm, text width=3.8cm},
tech/.style={draw=popgreen, fill=popgreenL, rounded corners=2pt, align=center, minimum height=1.1cm, text width=4.8cm},
fleche/.style={-{Stealth[length=2.5mm]}, very thick, popdark}
]
\node[cause, align=center] (c1) at (0,0){\textbf{Micro-organismes}\\(bactéries, levures, moisissures)\\-- Dégradation, fermentation, toxines};
\node[tech, align=center] (t1) at (8.2,0){Chaleur : \textbf{pasteurisation} (\SI{70}{--}\SI{95}{\celsius}), \textbf{stérilisation} (\textgreater{} \SI{100}{\celsius}), \textbf{cuisson confiture}\\$\rightarrow$ Destruction cellules/virus};
\draw[fleche] (c1) -- (t1);

\node[cause, align=center] (c2) at (0,-1.85){\textbf{Enzymes endogènes du fruit}\\(polyphénoloxydases, pectinestérases,\\peroxydases) -- Brunissement, ramollissement};
\node[tech, align=center] (t2) at (8.2,-1.85){\textbf{Blanchiment} (\SI{80}{--}\SI{100}{\celsius}, \SI{1}{--}\SI{3}{min}) ou \textbf{chauffage} rapide\\$\rightarrow$ Inactivation enzymatique (dénaturation)};
\draw[fleche] (c2) -- (t2);

\node[cause, align=center] (c3) at (0,-3.7){\textbf{Eau disponible ($a_w$ élevée)}\\(milieu de vie indispensable aux microbes)};
\node[tech, align=center] (t3) at (8.2,-3.7){\textbf{Séchage} (solaire, séchoir), \textbf{concentration} (évaporation),\\ \textbf{ajout solutés} : sucre (confiture), sel (saumure)\\$\rightarrow$ Abaissement $a_w$ \textless{} 0,60--0,70};
\draw[fleche] (c3) -- (t3);

\node[cause, align=center] (c4) at (0,-5.55){\textbf{Oxygène de l'air}\\(oxydation vitamines, couleurs, arômes ;\\développement microbes aérobies)};
\node[tech, align=center] (t4) at (8.2,-5.55){\textbf{Conditionnement hermétique} (bouteilles, boîtes, sachets),\\ \textbf{remplissage à chaud} (création vide partiel),\\ \textbf{désaération} (industriel), récipients pleins\\$\rightarrow$ Exclusion O$_2$};
\draw[fleche] (c4) -- (t4);

\node[font=\bfseries\small, poppink] at (0,1.1) {Cause d'altération principale};
\node[font=\bfseries\small, popgreen] at (8.2,1.1) {Technique(s) de neutralisation (souvent combinées)};
\end{tikzpicture}
\end{center}
\legende{Tableau de synthèse de la leçon : chaque cause d'altération des fruits (encadrés roses) est combattue par une ou plusieurs techniques de conservation (encadrés verts). Les techniques industrielles ou artisanales combinent souvent plusieurs effets simultanés : une confiture associe la chaleur (effet anti-microbien et anti-enzymatique) et le sucre (abaissement de l'activité de l'eau $a_w$) ; un jus pasteurisé associe la chaleur et l'exclusion de l'oxygène par remplissage à chaud/bouchage hermétique.}
\end{popfigure}

\begin{exoresolu}[Analyser une situation de transformation à l'échelle d'une coopérative]
Énoncé. Une coopérative de la région du Littoral récolte des mangues vendues \SI{25}{FCFA/kg} en pleine saison. Elle transforme chaque jour \SI{1}{t} de mangues en jus pasteurisé et obtient \SI{500}{L} de jus, vendus \SI{500}{FCFA/L}. Les coûts de transformation (sucre, bouteilles, bouchons, étiquettes, énergie, main-d'œuvre, transport, amortissement) s'élèvent à \SI{120000}{FCFA} par tonne de mangues brutes traitées. 
\begin{enumerate}
\item Calcule la valeur de la mangue brute et le chiffre d'affaires du jus pour une tonne transformée. Quel est le coefficient de multiplication de la valeur ?
\item Calcule le bénéfice net journalier de la coopérative (marge brute).
\item Deux jours après embouteillage, quelques bouteilles présentent des bouchons bombés et une odeur alcoolique marquée. Que doit faire la coopérative ? Justifie ta réponse en te référant aux risques toxicologiques.
\end{enumerate}
\tcblower
Solution détaillée.
\begin{enumerate}
\item \textbf{Valeur matière première} : $1000\,\text{kg} \times \SI{25}{FCFA/kg} = \SI{25000}{FCFA}$.\\
\textbf{Chiffre d'affaires (CA) jus} : $500\,\text{L} \times \SI{500}{FCFA/L} = \SI{250000}{FCFA}$.\\
\textbf{Coefficient de multiplication} : $\dfrac{\SI{250000}{FCFA}}{\SI{25000}{FCFA}} = 10$. La transformation a multiplié par 10 la valeur marchande de la récolte brute.
\item \textbf{Bénéfice net journalier (marge brute)} : $\text{CA} - \text{Coût matière première} - \text{Coûts de transformation} = \SI{250000}{FCFA} - \SI{25000}{FCFA} - \SI{120000}{FCFA} = \SI{105000}{FCFA}$ par tonne transformée.
\item \textbf{Conduite à tenir} : Les bouteilles bombées avec odeur alcoolique témoignent d'une fermentation active (production de \ce{CO2} et d'éthanol par des levures) due à une pasteurisation insuffisante (température/temps non respectés) ou une recontamination au conditionnement (bouteilles/bouchons non stériles, remplissage tiède). Ces bouteilles doivent être \textbf{immédiatement écartées, détruites (vidées, bouteilles relavées/rebouillies) et en aucun cas rebouillies pour être vendues ou consommées}. Justification : les levures et bactéries ont pu produire des métabolites indésirables ; surtout, si des moisissures se sont développées (même invisibles), elles ont pu synthétiser des \textbf{mycotoxines thermostables} (patuline, etc.) que l'ébullition ne détruit pas. La coopérative doit simultanément ouvrir une enquête interne : vérifier les courbes de pasteurisation (thermomètre étalonné), la propreté du conditionnement, l'hygiène du personnel, et corriger le protocole avant le lot suivant.
\end{enumerate}
\end{exoresolu}

\begin{aretenir}[L'essentiel de la section -- Fiche de révision Bac]
\begin{itemize}
\item Une séance pratique de transformation exige une hygiène stricte, non négociable : mains et matériel lavés/ébouillantés, eau potable, fruits triés (zéro tolérance pourri/moisi), couples température/durée respectés (pasteurisation \SI{80}{--}\SI{85}{\celsius} \SI{10}{--}\SI{15}{min}), conditionnement propre/chaud/hermétique, étiquetage complet de chaque lot (traçabilité).
\item La qualité d'un produit transformé se juge par des \textbf{critères organoleptiques} (couleur, odeur, goût, texture) et des \textbf{critères de conservation/stabilité} (absence fermentation, moisissure, bombage, intégrité fermeture).
\item \textbf{Règle de sécurité vitale} : tout produit bombé, fermenté, moisi, à odeur anormale est impropre à la consommation et doit être jeté. Le rebouillage \textbf{ne garantit pas} l'élimination des mycotoxines thermostables.
\item La transformation crée de la \textbf{valeur ajoutée} à chaque étape de la chaîne (coefficient 5 à 10 fréquent), réduit les pertes post-récolte (30--50\% sauvées), crée des emplois ruraux durables et substitue des produits locaux aux importations (souveraineté alimentaire).
\item Elle rend les vitamines (A, C) et antioxydants (lycopène) disponibles toute l'année, mais la vitamine C est partiellement détruite par la chaleur et le séchage : les traitements \textbf{rapides} et \textbf{protégés de l'air/lumière} limitent ces pertes.
\item \textbf{Idée directrice exigible} : \og À chaque cause d'altération (microbes, enzymes, eau libre, oxygène) correspond une technique de conservation qui la neutralise \fg (chaleur, abaissement $a_w$, exclusion O$_2$, acification). Les techniques réelles combinent ces leviers (ex. confiture = chaleur + sucre ; jus = chaleur + vide partiel).
\end{itemize}
\end{aretenir}

\newpage

\coursec{Activité d'intégration}

\courssub{Objectifs de l'activité}

À l'issue de cette activité d'intégration, tu seras capable de :

\begin{itemize}
\item interpréter des documents chiffrés (tableaux de résultats expérimentaux, courbes de survie microbienne) relatifs à la conservation des fruits ;
\item identifier les causes d'altération des fruits de saison à partir d'observations et de données expérimentales ;
\item choisir et justifier des techniques de transformation et de conservation adaptées à un contexte réel ;
\item établir l'organigramme de fabrication d'un produit transformé en précisant les paramètres critiques (températures, durées, hygiène) ;
\item rédiger une fiche-conseil argumentée intégrant une estimation économique simple.
\end{itemize}

\courssub{Situation}

\textbf{Sauver la récolte de mangues de la coopérative de Mora.}

Dans le département du Mayo-Sava (région de l'Extrême-Nord du Cameroun), la coopérative des producteurs de mangues de Mora regroupe 120 exploitants. Chaque année, au mois d'avril, la coopérative récolte environ \SI{20}{t} de mangues (variétés locales et variété améliorée « Julie »). Faute de débouchés suffisants au moment du pic de production, la coopérative perd habituellement \SI{40}{\%} de sa récolte : les fruits surmûris pourrissent sur place ou sont bradés.

Le document suivant présente l'évolution des prix de vente de la mangue fraîche sur le marché de Maroua au cours de l'année.

\begin{center}
\begin{tabularx}{\linewidth}{|l|X|X|X|X|}
\hline
\rowcolor{popblueL} \textbf{Période} & \textbf{Avril (pic)} & \textbf{Juillet} & \textbf{Octobre} & \textbf{Janvier} \\
\hline
Prix du kg de mangue fraîche (FCFA) & 100 & 250 & 400 & 500 \\
\hline
\end{tabularx}
\end{center}

Un technicien de l'IRAD (Institut de Recherche Agricole pour le Développement) a réalisé deux expériences pour aider la coopérative.

\textbf{Expérience 1.} Des mangues de même variété et de même maturité sont réparties en quatre lots de 100 fruits : lot A (non lavées, stockées à l'air ambiant, \SI{30}{\celsius}), lot B (lavées à l'eau propre additionnée d'eau de Javel diluée, stockées à l'air ambiant), lot C (non lavées, stockées au frais à \SI{12}{\celsius}), lot D (lavées et stockées au frais à \SI{12}{\celsius}). Le pourcentage de fruits sains est relevé au jour 0 (J0), à J+7 et à J+14.

\begin{center}
\begin{tabularx}{\linewidth}{|l|X|X|X|}
\hline
\rowcolor{popblueL} \textbf{Lot} & \textbf{J0} & \textbf{J+7} & \textbf{J+14} \\
\hline
A : non lavées, air ambiant & 100 \% & 55 \% & 20 \% \\
\hline
B : lavées, air ambiant & 100 \% & 75 \% & 45 \% \\
\hline
C : non lavées, frais (\SI{12}{\celsius}) & 100 \% & 85 \% & 65 \% \\
\hline
D : lavées, frais (\SI{12}{\celsius}) & 100 \% & 95 \% & 85 \% \\
\hline
\end{tabularx}
\end{center}

\textbf{Expérience 2.} Du jus de mangue est chauffé à différentes températures pendant \SI{1}{min}, puis le nombre de micro-organismes vivants par millilitre est dénombré. Les résultats sont traduits par la courbe de survie microbienne ci-dessous.

\begin{popfigure}
\begin{tikzpicture}
\begin{axis}[
width=0.9\linewidth, height=6cm,
xlabel={Température de chauffage (\si{\celsius})},
ylabel={Micro-organismes survivants (\%)},
xmin=20, xmax=100, ymin=0, ymax=105,
grid=both, grid style={popline!40},
legend style={at={(0.03,0.97)}, anchor=north west, font=\small}
]
\addplot[popcurve, domain=20:100, samples=200] {100/(1+exp((x-60)/3))};
\addlegendentry{Survie microbienne}
\addplot[dashed, poporange, thick] coordinates {(72,0) (72,50)};
\addplot[dashed, poporange, thick] coordinates {(20,2) (72,2)};
\node[poporange, font=\small] at (axis cs:78,10) {\SI{72}{\celsius} : \textasciitilde 2\%};
\end{axis}
\end{tikzpicture}
\legende{Courbe de survie des micro-organismes du jus de mangue en fonction de la température de chauffage (durée fixe : \SI{1}{min}). La survie chute fortement entre \SI{60}{\celsius} et \SI{75}{\celsius} ; un chauffage à \SI{72}{\celsius} ne laisse plus que quelques survivants (\textasciitilde 2\%).}
\end{popfigure}

Tu es membre du groupe d'élèves de Terminale chargé par la coopérative de proposer une stratégie pour réduire les pertes et augmenter les revenus des producteurs.

\courssub{Consignes}

\begin{enumerate}
\item À partir de l'expérience 1, comparer l'évolution des quatre lots et identifier les deux facteurs qui limitent l'altération des mangues. En déduire deux causes d'altération des fruits.
\item Expliquer, en t'appuyant sur tes connaissances, pourquoi le lavage et le stockage au frais ralentissent l'altération.
\item À partir de la courbe de l'expérience 2, déterminer graphiquement la température minimale de chauffage permettant de détruire la quasi-totalité des micro-organismes. Nommer la technique de conservation qui utilise ce principe.
\item Proposer deux techniques de transformation adaptées à la coopérative (par exemple jus pasteurisé et mangue séchée) et justifier chaque choix par un argument scientifique et un argument économique.
\item Établir l'organigramme de fabrication du jus de mangue pasteurisé en précisant les paramètres critiques (hygiène, température, durée, conditionnement).
\item Calculer : (a) la masse de mangues perdue chaque année ; (b) le revenu que rapporterait la vente de \SI{8}{t} de mangues transformées en jus, sachant que \SI{1}{kg} de mangues donne \SI{0,5}{L} de jus vendu \SI{500}{FCFA} le litre ; (c) comparer avec la vente de ces mêmes \SI{8}{t} en fruits frais au prix d'avril.
\item Rédiger une fiche-conseil d'une quinzaine de lignes à destination de la coopérative, comportant : le diagnostic, les techniques recommandées, les précautions d'hygiène et l'intérêt économique.
\end{enumerate}

\courssub{Ressources à mobiliser}

\begin{aretenir}[Ce qu'il faut avoir en tête]
\begin{itemize}
\item Les causes d'altération des fruits : causes \cle{microbiennes} (moisissures, levures, bactéries), causes \cle{enzymatiques} (enzymes propres du fruit : brunissement enzymatique, ramollissement pectique), causes \cle{physico-chimiques} (chaleur, lumière, oxygène) et causes \cle{biologiques} (insectes, rongeurs).
\item Les principes de conservation : \cle{abaissement de la température} (réfrigération : ralentit les réactions enzymatiques et la multiplication microbienne), \cle{élévation de la température} (pasteurisation, stérilisation : détruit les micro-organismes et inactive les enzymes), \cle{réduction de l'activité de l'eau} (séchage, concentration en sucre : empêche le développement microbien), \cle{moyens chimiques} (vinaigre, sel, antiseptiques alimentaires).
\item Pasteurisation : chauffage modéré (environ \SI{72}{\celsius} pendant \SI{15}{s} à \SI{1}{min} selon le produit) suivi d'un refroidissement rapide ; elle détruit la flore pathogène et altérative sans altérer notablement les qualités organoleptiques et nutritionnelles.
\item Calculs : pourcentage de pertes $= \dfrac{\text{masse perdue}}{\text{masse totale}} \times 100$ ; rendement en jus $= \dfrac{\text{volume de jus}}{\text{masse de fruits}}$.
\end{itemize}
\end{aretenir}

\begin{savaistu}[Le cas de la tomate au Cameroun]
La tomate, autre fruit de saison majeur (zones de l'Ouest, du Littoral, du Sud), subit des pertes post-récolte encore plus élevées (jusqu'à 50 \%). Les techniques de transformation locales incluent : la \emph{concentrée de tomate} (cuisson longue à feu vif pour évaporer l'eau, conditionnement en bocaux stérilisés) et le \emph{séchage solaire} (tranches sur nattes ou séchoirs améliorés). La concentration en solides (sucre, sel, acides) abaisse l'activité de l'eau ($a_w < 0,7$), assurant une conservation de plusieurs mois sans chaîne du froid, ce qui est crucial dans les zones rurales non électrifiées.
\end{savaistu}

\courssub{Démarche et corrigé commenté}

\begin{exoresolu}[Sauver la récolte de Mora : résolution complète]
\textbf{Énoncé résumé.} Une coopérative perd \SI{40}{\%} de ses \SI{20}{t} de mangues. Interpréter les expériences, proposer deux techniques de transformation, construire l'organigramme du jus pasteurisé et estimer le gain économique.

\tcblower

\textbf{1. Interprétation de l'expérience 1.} Comparons les lots deux à deux. À J+14, le lot B (lavé, ambiant, 45 \% de fruits sains) est mieux conservé que le lot A (non lavé, ambiant, 20 \%) : le \cle{lavage} améliore la conservation. De même, le lot C (non lavé, frais, 65 \%) est mieux conservé que le lot A : le \cle{froid} améliore la conservation. Le lot D (lavé + frais, 85 \%) cumule les deux effets : c'est le meilleur résultat. Les deux facteurs limitant l'altération sont donc l'hygiène (réduction de la charge microbienne initiale) et la basse température. On en déduit deux causes majeures d'altération : le développement des \cle{micro-organismes} (moisissures, bactéries) présents à la surface des fruits, et l'activité des \cle{enzymes} propres du fruit (mûrissement, ramollissement pectique, brunissement), accélérée par la chaleur ambiante.

\textbf{2. Explication mécaniste.} Le lavage à l'eau javellisée (solution d'hypochlorite de sodium) élimine ou détruit une grande partie des micro-organismes déposés à la surface du fruit lors de la récolte et du transport : la contamination initiale étant plus faible, le pourrissement microbien démarre plus tard. Le stockage à \SI{12}{\celsius} ralentit cinétiquement l'ensemble des réactions biochimiques : la respiration du fruit (consommation d'O$_2$, production de CO$_2$ et de chaleur), l'activité des enzymes pectiques (ramollissement) et polyphénoloxydases (brunissement), ainsi que la vitesse de multiplication des micro-organismes résidents (dont la croissance est optimale vers \SI{30}{\celsius}). D'où l'effet cumulatif observé sur le lot D.

\textbf{3. Lecture de la courbe (Expérience 2).} La courbe de survie microbienne présente une décroissance sigmoïde. Graphiquement, la proportion de survivants devient très faible (inférieure à 5 \%) à partir de \SI{72}{\celsius} maintenue \SI{1}{min} (environ 2 \% d'après le modèle). La technique de conservation qui repose sur ce principe (chauffage modéré, durée courte, refroidissement rapide, conditionnement hermétique) est la \cle{pasteurisation}. Elle vise une réduction logarithmique de la flore (cible souvent 5 décimales, soit 99,999 \%) tout en préservant les qualités gustatives et vitaminique (vitamine C thermosensible).

\textbf{4. Choix de deux techniques de transformation.}
\begin{itemize}
\item \emph{Jus de mangue pasteurisé} : \textbf{Argument scientifique} : la pasteurisation (\SI{72}{\celsius}, \SI{1}{min}) détruit les micro-organismes pathogènes et altératifs et inactive les enzymes (pectinestérase, polyphénoloxydase) responsables de la trouble et du brunissement ; le conditionnement à chaud hermétique empêche toute recontamination. \textbf{Argument économique} : le jus se conserve \SI{6}{à 12}{mois} à température ambiante, permettant une vente différée (juillet-janvier) aux prix élevés (250 à 500 FCFA/kg équivalent), lissant le revenu annuel.
\item \emph{Mangue séchée (chips ou quartiers)} : \textbf{Argument scientifique} : le séchage (solaire sur claies hygiéniques ou séchoir artificiel à \SIrange{50}{60}{\celsius}) abaisse l'activité de l'eau ($a_w$) sous le seuil critique de \num{0,65}, inhibant tout développement microbien et ralentissant drastiquement l'activité enzymatique résiduelle. \textbf{Argument économique} : produit à haute valeur ajoutée, léger (réduction masse/volume \textasciitilde 80 \%), non périssable, transportable vers les grands centres urbains (Yaoundé, Douala) ou exportable, se vendant \SIrange{2500}{3500}{\FCFA\per\kilo\gram}.
\end{itemize}

\textbf{5. Organigramme de fabrication du jus de mangue pasteurisé.}

\begin{center}
\begin{tabularx}{\linewidth}{|l|X|}
\hline
\rowcolor{popblueL} \textbf{Étape} & \textbf{Opérations et paramètres critiques} \\
\hline
1. Réception & Tri manuel : élimination des fruits pourris, blessés, verts. \textbf{Critique :} hygiène du personnel (lavage mains, tabliers propres). \\
\hline
2. Lavage & Bain eau potable + hypochlorite \SIrange{50}{100}{ppm} (\SIrange{2}{5}{min}), rinçage eau potable. \textbf{Critique :} concentration chlore, temps de contact. \\
\hline
3. Préparation & Épluchage (vapeur ou manuel), dénoyautage. \textbf{Critique :} propreté couteaux/plan de travail. \\
\hline
4. Extraction & Pulpage / mixage pulpe. Ajout eau potable, sucre (sirop \SIrange{15}{20}{\%}), correction pH (\numrange{3,5}{4,0} acide citrique). \textbf{Critique :} standardisation °Brix, pH. \\
\hline
5. \textbf{Pasteurisation} & \textbf{Chauffage \SI{72}{\celsius} pendant \SI{1}{min} (échangeur tubulaire ou bain-marie agitée).} \textbf{Paramètre CRITIQUE :} couple Temps/Température (traçabilité). \\
\hline
6. Remplissage & Remplissage à chaud (\textasciitilde \SI{85}{\celsius}) dans bouteilles/verrines \textbf{préalablement stérilisées} (vapeur \SI{100}{\celsius} \SI{10}{min}). \\
\hline
7. Bouchage & Bouchage hermétique immédiat (capsules neuves, contrôle couple de serrage). \textbf{Critique :} étanchéité (test inversion). \\
\hline
8. Refroidissement & Refroidissement rapide par aspersion/eau glacée à \textless \SI{40}{\celsius} (éviter thermorésistance, trouble). \\
\hline
9. Étiquetage / Stockage & Étiquetage (DLUO, lot, ingrédients). Stockage local frais, sec, à l'abri de la lumière. \\
\hline
\end{tabularx}
\end{center}

\textbf{6. Calculs économiques.}
\begin{align*}
\text{(a) Masse perdue} &= \SI{20}{t} \times \frac{40}{100} = \SI{8}{t} = \SI{8000}{kg}.\\
\text{(b) Volume de jus obtenu} &= \SI{8000}{kg} \times \SI{0,5}{L/kg} = \SI{4000}{L}.\\
\text{Revenu (Chiffre d'affaires) du jus} &= \num{4000} \times \num{500} = \num{2000000} \text{ FCFA}.\\
\text{(c) Revenu vente fraîche (avril)} &= \SI{8000}{kg} \times \num{100} = \num{800000} \text{ FCFA}.\\
\text{Gain brut additionnel} &= \num{2000000} - \num{800000} = \num{1200000} \text{ FCFA}.
\end{align*}
La transformation en jus permet de multiplier le chiffre d'affaires par 2,5 par rapport à la vente en frais au prix du pic de saison (\num{1200000} FCFA de gain brut), avant déduction des coûts de transformation (intrants, énergie, main-d'œuvre, emballages, amortissement).

\textbf{7. Fiche-conseil (modèle pour la coopérative).}

\textbf{Diagnostic :} Perte annuelle de \SI{8}{t} de mangues (\SI{40}{\%} récolte) due au pourrissement microbien (contamination de surface) et au surmûrissement enzymatique accéléré par la chaleur (\SI{30}{\celsius} ambiant). Vente contrainte au prix plancher (\num{100} FCFA/kg) en avril.

\textbf{Techniques recommandées :}
\begin{enumerate}
\item \textbf{Conservation courte durée (vente fraîche) :} Lavage systématique eau javellisée (\SI{50}{ppm} Cl$_2$) + stockage ventilé à \SI{12}{\celsius} (chambre froide solaire ou cave aménagée) $\rightarrow$ gain de 2 à 3 semaines, vente à \num{250} FCFA/kg (juillet).
\item \textbf{Transformation (valeur ajoutée) :}
      \begin{itemize}
      \item Jus pasteurisé : \SI{72}{\celsius}/\SI{1}{min}, remplissage à chaud bouteilles stérilisées, DLUO 12 mois.
      \item Mangue séchée : séchoir solaire hygiénique (température \textgreater \SI{50}{\celsius}, $a_w < 0,65$), conditionnement sachets alimentaires.
      \end{itemize}
\end{enumerate}

\textbf{Précautions d'hygiène (Bonnes Pratiques de Fabrication) :} Eau potable certifiée, locaux fermés (anti-insectes/rongeurs), personnel formé (hygiène mains, tenues), nettoyage/désinfection quotidien du matériel, traçabilité lots (date, variété, opérateur).

\textbf{Intérêt économique :} Sur les \SI{8}{t} habituellement perdues (0 FCFA), la transformation en jus génère \num{2000000} FCFA de chiffre d'affaires (vs \num{800000} FCFA en frais avril). Même après déduction des charges estimées (\textasciitilde \num{600000} FCFA), le bénéfice net dépasse \num{1000000} FCFA/an, sécurisant le revenu des 120 producteurs hors saison.
\end{exoresolu}

\courssub{Bilan de l'activité}

Cette activité t'a permis de mettre en évidence plusieurs acquis essentiels de la leçon. D'abord, les pertes post-récolte ne sont pas une fatalité : l'interprétation rigoureuse d'un tableau expérimental à quatre lots a montré que deux facteurs simples — l'\cle{hygiène} (lavage désinfectant) et le \cle{froid} (réfrigération) — suffisent à faire passer la proportion de fruits sains à J+14 de 20 \% à 85 \%. Ensuite, la lecture d'une courbe de survie microbienne a permis de justifier scientifiquement le choix du couple température/durée de la \cle{pasteurisation} (\SI{72}{\celsius}, \SI{1}{min}), au lieu de l'apprendre mécaniquement. Enfin, tu as relié la démarche scientifique à la réalité socio-économique camerounaise : transformer les fruits de saison (jus pasteurisé, séchage, concentré de tomate) réduit les pertes, stabilise les revenus des producteurs entre les saisons et crée de la valeur ajoutée locale. C'est exactement la compétence visée par le programme : mobiliser des savoirs et savoir-faire de biotechnologie pour résoudre un problème concret lié à la couverture des besoins alimentaires.

\begin{attention}[Pour aller plus loin sans se tromper]
\begin{itemize}
\item Dans tes calculs économiques, n'oublie jamais de préciser les unités (t, kg, L, FCFA) et de distinguer le \emph{chiffre d'affaires} (recettes totales) du \emph{bénéfice net} (CA $-$ charges variables $-$ charges fixes). La transformation a un coût (sucre, bouteilles, énergie, main-d'œuvre, amortissement matériel) qu'il faudrait déduire pour obtenir le gain réel.
\item Au Baccalauréat, une réponse chiffrée sans unité ni justification expérimentale (référence au tableau ou à la courbe) perd des points.
\item Pour l'organigramme, les \textbf{paramètres critiques} (Points Critiques de Maîtrise - PCM au sens HACCP) doivent être explicites : température \textbf{ET} durée de pasteurisation, stérilisation des contenants, étanchéité du bouchage, pH final \textless 4,5.
\item Ne confonds pas \emph{stérilisation} (destruction de toutes formes microbiennes, \textgreater \SI{100}{\celsius}, conservation années) et \emph{pasteurisation} (destruction flore pathogène/altérative, \SIrange{65}{85}{\celsius}, conservation mois/année au frais ou ambiant si pH acide).
\end{itemize}
\end{attention}

\newpage

\coursec{Méthodes et exercices résolus}

\courssub{Fiche méthode 1 : Décrire une technique de transformation ou de conservation d'un fruit}

\begin{methode}[Décrire une technique de transformation/conservation]
Pour décrire correctement une technique (savoir-faire exigible), suis toujours le même plan en 5 étapes :
\begin{enumerate}
\item \textbf{Nommer la technique} et la classer (transformation : séchage, jus, confiture, purée, concentré ; conservation : froid, chaleur, produits chimiques, atmosphère contrôlée).
\item \textbf{Décrire le protocole} dans l'ordre chronologique : réception et tri des fruits $\rightarrow$ lavage $\rightarrow$ épluchage/dénoyautage $\rightarrow$ opération principale (chauffage, séchage, ajout de sucre...) $\rightarrow$ conditionnement $\rightarrow$ stockage.
\item \textbf{Donner les paramètres chiffrés} : températures, durées, concentrations (ex. pasteurisation à environ \SI{70}{\celsius} pendant quelques minutes ; confiture à environ \SI{65}{\percent} d'extrait sec sucré).
\item \textbf{Justifier scientifiquement} chaque étape clé : quel agent d'altération est éliminé ou bloqué (micro-organismes, enzymes, oxygène, eau) et par quel mécanisme.
\item \textbf{Conclure} par le produit obtenu, sa durée de conservation et ses conditions de stockage.
\end{enumerate}
\textbf{Réflexe Bac :} une description sans justification scientifique (principe de conservation) ne vaut que la moitié des points.
\end{methode}

\begin{exoresolu}[Description de la fabrication de jus de mangue pasteurisé]
Lors d'un TP au lycée de Ngaoundéré, un groupe d'élèves dispose de mangues mûres saines et souhaite produire un jus conservable pendant plusieurs semaines à température ambiante. Décris la technique à suivre en justifiant chaque étape essentielle.
\tcblower
\textbf{1. Nom et principe de la technique.} Il s'agit d'une \cle{transformation} en jus, associée à une \cle{conservation par la chaleur} (pasteurisation). Le principe est la destruction des micro-organismes et l'inactivation des enzymes par un traitement thermique, suivie d'un conditionnement hermétique qui empêche toute recontamination.

\textbf{2. Protocole chronologique.}
\begin{enumerate}
\item \textbf{Tri} : éliminer les mangues pourries ou blessées, car elles sont des foyers de contamination microbienne.
\item \textbf{Lavage} à l'eau potable (éventuellement additionnée d'une solution chlorée) : réduction de la flore microbienne de surface.
\item \textbf{Épluchage et dénoyautage}, puis \textbf{mixage} de la pulpe : obtention d'une purée homogène.
\item \textbf{Éventuellement tamisage et ajout d'eau potable et de sucre} pour ajuster la consistance et le goût.
\item \textbf{Pasteurisation} : chauffer le jus à environ \SI{70}{\celsius} pendant quelques minutes. Ce traitement détruit les formes végétatives des micro-organismes (bactéries, levures, moisissures) et dénature les enzymes responsables du brunissement et de l'altération du goût.
\item \textbf{Remplissage à chaud} des bouteilles préalablement lavées et ébouillantées, puis \textbf{bouchage immédiat hermétique} : le vide partiel qui se forme au refroidissement limite l'oxygène et garantit l'étanchéité.
\item \textbf{Refroidissement rapide} puis \textbf{stockage} à l'abri de la lumière et de la chaleur.
\end{enumerate}

\textbf{3. Justifications scientifiques.} La chaleur (\SI{70}{\celsius}) provoque la dénaturation des protéines microbiennes et enzymatiques ; le bouchage hermétique empêche l'entrée d'air (oxygène) et de nouveaux germes ; le stockage frais et sombre ralentit les réactions chimiques résiduelles (oxydation de la vitamine C, brunissement).

\textbf{4. Conclusion.} On obtient un jus de mangue pasteurisé, stable plusieurs semaines à température ambiante, et plusieurs mois au frais. Cette technique valorise les surplus de mangues de saison et réduit les pertes post-récolte.
\end{exoresolu}

\courssub{Fiche méthode 2 : Pratiquer la fabrication d'une confiture (transformation par le sucre)}

\begin{methode}[Réussir une confiture de mangue ou de tomate]
\begin{enumerate}
\item \textbf{Retenir le principe} : le sucre à haute concentration abaisse l'\cle{activité de l'eau} ($a_w$) du milieu ; les micro-organismes, privés d'eau disponible, ne peuvent plus se multiplier (plasmolyse).
\item \textbf{Respecter les proportions} : environ \SI{1}{kg} de sucre pour \SI{1}{kg} de pulpe (ratio ajustable entre 0,8 et 1), pour atteindre un extrait sec final d'environ \SI{65}{\percent}.
\item \textbf{Conduire la cuisson} : mélanger pulpe et sucre, porter à ébullition en remuant (éviter la caramélisation au fond), cuire jusqu'à la \textbf{prise} (test de la goutte sur assiette froide : la goutte se fige).
\item \textbf{Conditionner à chaud} dans des pots en verre stérilisés (ébouillantés \SI{10}{min}), fermer immédiatement, retourner les pots quelques minutes pour traiter thermiquement le couvercle.
\item \textbf{Stocker} au frais, au sec et à l'obscurité.
\end{enumerate}
\textbf{Formule à retenir :} masse de sucre $=$ masse de pulpe $\times$ ratio sucre/pulpe.
\end{methode}

\begin{exoresolu}[Fabrication de confiture de mangues et calcul du rendement]
Une coopérative de femmes de Garoua transforme \SI{40}{kg} de mangues. Après épluchage et dénoyautage, les déchets représentent \SI{35}{\percent} de la masse des fruits. La pulpe obtenue est mélangée au sucre dans le ratio 0,9 kg de sucre pour 1 kg de pulpe, puis cuite. La cuisson évapore \SI{20}{\percent} de la masse du mélange pulpe-sucre.
\begin{enumerate}
\item Calcule la masse de pulpe et la masse de sucre utilisée.
\item Calcule la masse de confiture obtenue.
\item Explique le principe de conservation de la confiture.
\end{enumerate}
\tcblower
\textbf{1. Masse de pulpe et masse de sucre.}

Masse des déchets : $40 \times \dfrac{35}{100} = \SI{14}{kg}$.

Masse de pulpe : $40 - 14 = \SI{26}{kg}$.

Masse de sucre : $26 \times 0{,}9 = \SI{23,4}{kg}$.

\textbf{2. Masse de confiture obtenue.}

Masse du mélange avant cuisson : $26 + 23{,}4 = \SI{49,4}{kg}$.

Masse d'eau évaporée : $49{,}4 \times \dfrac{20}{100} = \SI{9,88}{kg}$.

Masse de confiture : $49{,}4 - 9{,}88 = \SI{39,52}{kg}$, soit environ \SI{39,5}{kg}.

\textbf{3. Principe de conservation.} La confiture contient environ \SI{65}{\percent} d'extrait sec sucré. Cette forte concentration en sucre abaisse l'activité de l'eau ($a_w$) : l'eau n'est plus disponible pour le métabolisme microbien. Les micro-organismes subissent une plasmolyse (perte d'eau par osmose) et ne peuvent ni se multiplier ni altérer le produit. Le conditionnement à chaud en pots stérilisés et fermés hermétiquement détruit les germes résiduels et empêche la recontamination. L'acidité naturelle de la mangue renforce l'effet conservateur.

\textbf{Conclusion.} À partir de \SI{40}{kg} de mangues fraîches, la coopérative obtient environ \SI{39,5}{kg} de confiture, produit stable à température ambiante pendant plusieurs mois.
\end{exoresolu}

\courssub{Fiche méthode 3 : Identifier les causes d'altération et choisir la technique adaptée}

\begin{methode}[Raisonner en « agent d'altération $\rightarrow$ barrière »]
\begin{enumerate}
\item \textbf{Identifier l'agent d'altération} évoqué dans l'énoncé : micro-organismes (bactéries, levures, moisissures), enzymes endogènes (brunissement, ramollissement), réactions chimiques (oxydation par le dioxygène \ce{O2}), insectes, perte d'eau (flétrissement).
\item \textbf{Associer à chaque agent la barrière adaptée} :
\begin{itemize}
\item chaleur (pasteurisation, stérilisation, ébullition) $\rightarrow$ détruit germes et enzymes ;
\item froid (réfrigération vers \SIrange{0}{4}{\celsius}) $\rightarrow$ ralentit sans détruire ;
\item sucre ou sel concentré $\rightarrow$ abaisse l'activité de l'eau ;
\item séchage/déshydratation $\rightarrow$ retire l'eau ($a_w$ faible) ;
\item acidification (vinaigre, jus de citron) $\rightarrow$ pH défavorable aux germes ;
\item conditionnement hermétique/sous vide $\rightarrow$ supprime l'oxygène et la recontamination.
\end{itemize}
\item \textbf{Vérifier la cohérence} : le froid ne stérilise pas ; le séchage exige une protection contre l'humidité ; une boîte mal fermée annule la pasteurisation.
\item \textbf{Rédiger} la réponse en reliant explicitement agent $\rightarrow$ technique $\rightarrow$ mécanisme.
\end{enumerate}
\end{methode}

\begin{exoresolu}[Diagnostic d'altérations au marché de Mokolo]
Au marché de Mokolo (Yaoundé), on observe : (a) des tomates molles présentant des taches noires avec duvet gris ; (b) des mangues tranchées qui brunissent à l'air en quelques heures ; (c) des mangues séchées devenues moisies après stockage dans des sacs en plastique fermés. Pour chaque cas, identifie la cause d'altération et propose une technique de prévention justifiée.
\tcblower
\textbf{(a) Tomates molles, taches noires à duvet gris.} Le duvet est un mycélium : il s'agit d'une \cle{contamination par des moisissures} (pourriture grise, type \emph{Botrytis}), favorisée par les blessures et l'humidité. \textbf{Prévention :} trier les fruits blessés, éviter les chocs pendant le transport, conserver les tomates à température fraîche (environ \SIrange{10}{12}{\celsius} pour ne pas les blesser par le froid) et dans un local aéré ; à long terme, les transformer en purée ou concentré pasteurisé, la chaleur détruisant le mycélium et la plupart des spores.

\textbf{(b) Brunissement des mangues tranchées.} C'est le \cle{brunissement enzymatique} : au contact du dioxygène de l'air, des enzymes (polyphénoloxydases) oxydent les composés phénoliques de la pulpe en pigments bruns. \textbf{Prévention :} limiter le contact avec l'air (film alimentaire, conditionnement hermétique), acidifier la surface (jus de citron : le pH acide inhibe les enzymes), ou chauffer (dénaturation des enzymes).

\textbf{(c) Moisissure des mangues séchées en sacs fermés.} Le séchage a été \cle{insuffisant} : l'humidité résiduelle (activité de l'eau encore trop élevée) a permis le développement de moisissures, d'autant que le sac fermé emprisonne la vapeur d'eau. \textbf{Prévention :} prolonger le séchage jusqu'à une humidité finale suffisamment basse (produit souple mais sec au toucher), laisser refroidir avant ensachage, utiliser des emballages étanches à l'humidité et stocker au sec.

\textbf{Conclusion.} Toute technique de conservation efficace repose sur l'identification préalable de l'agent d'altération : germes, enzymes, oxygène ou eau.
\end{exoresolu}

\courssub{Fiche méthode 4 : Pratiquer la transformation de la tomate (purée et concentré)}

\begin{methode}[Transformer la tomate en purée ou concentré]
\begin{enumerate}
\item \textbf{Tri et lavage} des tomates bien mûres (éliminer les fruits verts ou pourris).
\item \textbf{Ébouillantage} (\SIrange{1}{2}{min} dans l'eau bouillante) : facilite l'\textbf{épluchage} et inactive partiellement les enzymes.
\item \textbf{Broyage puis tamisage} : élimination des peaux et graines, obtention du jus ou de la purée brute.
\item \textbf{Concentration par évaporation} : chauffage à feu doux en remuant jusqu'à la consistance voulue (purée, puis concentré si l'évaporation est poussée). Plus le produit est concentré, plus son activité de l'eau est faible, meilleure est sa conservation.
\item \textbf{Pasteurisation et conditionnement à chaud} en bocaux ou sachets hermétiques ; fermeture immédiate.
\item \textbf{Stockage} à l'abri de la chaleur et de la lumière.
\end{enumerate}
\textbf{Réflexe :} distinguer toujours purée (peu concentrée) et concentré (fortement concentré, extrait sec plus élevé, conservation plus longue).
\end{methode}

\begin{exoresolu}[Production de concentré de tomate par un club de jeunes]
Un club de jeunes de Bafoussam récolte \SI{60}{kg} de tomates. Après tri, \SI{10}{\percent} des fruits sont écartés. L'épluchage et le tamisage éliminent ensuite \SI{8}{\percent} de la masse restante. La purée obtenue est concentrée par évaporation jusqu'à perdre \SI{60}{\percent} de sa masse en eau.
\begin{enumerate}
\item Calcule la masse de concentré obtenue.
\item Cite deux précautions à prendre lors du conditionnement et justifie-les.
\end{enumerate}
\tcblower
\textbf{1. Masse de concentré.}

Après tri : $60 \times \left(1 - \dfrac{10}{100}\right) = 60 \times 0{,}9 = \SI{54}{kg}$.

Après épluchage et tamisage : $54 \times (1 - 0{,}08) = 54 \times 0{,}92 = \SI{49,68}{kg}$ de purée.

Après concentration (perte de \SI{60}{\percent}) : $49{,}68 \times (1 - 0{,}60) = 49{,}68 \times 0{,}40 = \SI{19,872}{kg}$.

Soit environ \SI{19,9}{kg} de concentré de tomate.

\textbf{2. Précautions de conditionnement.}
\begin{itemize}
\item \textbf{Remplir les récipients à chaud et les fermer immédiatement} : la chaleur détruit les germes présents et le refroidissement crée une dépression qui assure l'étanchéité, empêchant l'entrée d'air et de micro-organismes.
\item \textbf{Utiliser des récipients propres et préalablement stérilisés (ébouillantés)} : sinon, les germes déposés sur les parois recontaminent le concentré et annulent l'effet de la pasteurisation.
\end{itemize}
(On peut aussi citer : stockage à l'abri de la lumière pour limiter l'oxydation des pigments et des vitamines.)

\textbf{Conclusion.} De \SI{60}{kg} de tomates fraîches, on tire environ \SI{19,9}{kg} de concentré, produit de forte valeur ajoutée, transportable et conservable plusieurs mois.
\end{exoresolu}

\courssub{Fiche méthode 5 : Pratiquer le séchage des fruits (cas de la mangue séchée)}

\begin{methode}[Réussir le séchage de la mangue]
\begin{enumerate}
\item \textbf{Choisir} des mangues mûres mais fermes, saines ; laver, éplucher, dénoyauter.
\item \textbf{Trancher} en lamelles régulières (environ \SIrange{5}{10}{mm}) : plus les tranches sont fines et régulières, plus le séchage est rapide et homogène.
\item \textbf{Prétraiter éventuellement} (citronnage ou trempage rapide dans un sirop léger) pour limiter le brunissement enzymatique.
\item \textbf{Sécher} au soleil sur claies surélevées et couvertes d'un filet (protection contre poussières, mouches et autres insectes), ou en séchoir solaire/électrique à environ \SIrange{50}{60}{\celsius}, jusqu'à un produit souple et sec au toucher.
\item \textbf{Laisser refroidir} avant emballage (sinon condensation dans le sachet $\rightarrow$ moisissures).
\item \textbf{Conditionner} en sachets étanches à l'humidité, stocker au sec et à l'abri de la lumière.
\end{enumerate}
\textbf{Principe à citer :} le retrait de l'eau abaisse l'activité de l'eau ($a_w$) sous le seuil de développement des micro-organismes.
\end{methode}

\begin{exoresolu}[Séchage de mangues et calcul du taux d'humidité]
Une productrice de Maroua sèche des lamelles de mangue. Une lamelle fraîche de \SI{50}{g} contient \SI{82}{\percent} d'eau. Après séchage, la lamelle ne contient plus que \SI{15}{\percent} d'eau, la matière sèche étant conservée.
\begin{enumerate}
\item Calcule la masse de matière sèche de la lamelle fraîche.
\item Calcule la masse finale de la lamelle séchée.
\item Pourquoi le séchage permet-il la conservation ?
\end{enumerate}
\tcblower
\textbf{1. Masse de matière sèche.}

Masse d'eau initiale : $50 \times \dfrac{82}{100} = \SI{41}{g}$.

Matière sèche : $50 - 41 = \SI{9}{g}$ (soit \SI{18}{\percent} de \SI{50}{g} : $50 \times 0{,}18 = \SI{9}{g}$).

\textbf{2. Masse finale.} La matière sèche (\SI{9}{g}) est inchangée et représente désormais $100\,\% - 15\,\% = 85\,\%$ de la masse finale $m$ :
\[ m = \dfrac{9}{0{,}85} \approx \SI{10,6}{g}. \]
La lamelle a donc perdu $50 - 10{,}6 \approx \SI{39,4}{g}$ d'eau.

\textbf{3. Principe de conservation.} Le séchage retire l'eau libre du fruit : l'activité de l'eau ($a_w$) descend sous le seuil nécessaire à la croissance des bactéries, levures et moisissures, et les réactions enzymatiques, qui exigent un milieu hydraté, sont bloquées. Le fruit séché devient ainsi stable, à condition d'être protégé de la réhumidification par un emballage étanche.

\textbf{Conclusion.} La lamelle passe de \SI{50}{g} à environ \SI{10,6}{g} ; le séchage conserve le fruit en privant les agents d'altération de l'eau indispensable à leur activité.
\end{exoresolu}

\courssub{Fiche méthode 6 : Analyser les enjeux socio-économiques et la qualité des produits}

\begin{methode}[Construire une argumentation « qualité et enjeux »]
\begin{enumerate}
\item \textbf{Dégager le problème} : saisonnalité et abondance $\rightarrow$ surproduction ponctuelle $\rightarrow$ pertes post-récolte pouvant atteindre \SIrange{30}{50}{\percent} des fruits au Cameroun $\rightarrow$ baisse des prix et appauvrissement des producteurs.
\item \textbf{Montrer le rôle de la transformation/conservation} : réduction des pertes, disponibilité toute l'année, création de revenus (coopératives, PME agroalimentaires, emploi des jeunes et des femmes), amélioration de la sécurité alimentaire.
\item \textbf{Citer les critères de qualité et d'hygiène} : matières premières saines, eau potable, matériel propre, respect des paramètres (température, durée, dosage en sucre), emballage propre et étiqueté (nom du produit, date de fabrication, durée de conservation).
\item \textbf{Conclure} par un jugement équilibré : la transformation artisanale est rentable mais exige rigueur hygiénique et contrôle qualité pour être durable et sûre.
\end{enumerate}
\end{methode}

\begin{exoresolu}[Étude de cas : rentabilité d'une unité de transformation de mangues]
Pendant la saison, le kilogramme de mangues se vend \num{100}~FCFA à Ngaoundéré, mais environ \SI{40}{\percent} de la récolte est perdue faute de débouchés. Une coopérative achète \SI{500}{kg} de mangues et en extrait \SI{60}{\percent} de pulpe. Elle produit de la confiture avec \SI{1}{kg} de sucre (à \num{700}~FCFA le kg) par kg de pulpe. On néglige ici l'évaporation : la masse de confiture obtenue est prise égale à la masse du mélange pulpe-sucre. La confiture est vendue \num{2500}~FCFA le kg.
\begin{enumerate}
\item Calcule le coût total des matières premières (mangues + sucre).
\item Calcule la recette et le bénéfice brut.
\item Dégage deux avantages socio-économiques de cette activité.
\end{enumerate}
\tcblower
\textbf{1. Coût des matières premières.}

Masse de pulpe : $500 \times 0{,}60 = \SI{300}{kg}$.

Coût des mangues : $500 \times 100 = \num{50000}$~FCFA.

Masse de sucre : \SI{300}{kg} ; coût : $300 \times 700 = \num{210000}$~FCFA.

Coût total : $\num{50000} + \num{210000} = \num{260000}$~FCFA.

\textbf{2. Recette et bénéfice brut.}

Masse de confiture (évaporation négligée selon l'énoncé) : $300 + 300 = \SI{600}{kg}$.

Recette : $600 \times \num{2500} = \num{1500000}$~FCFA.

Bénéfice brut : $\num{1500000} - \num{260000} = \num{1240000}$~FCFA.

\textbf{3. Avantages socio-économiques.}
\begin{itemize}
\item \textbf{Réduction des pertes post-récolte} : les mangues autrefois perdues (jusqu'à \SI{40}{\percent} de la récolte) sont valorisées, ce qui augmente le revenu des producteurs.
\item \textbf{Création d'emplois et de revenus} pour les femmes et les jeunes (transformation, vente), et \textbf{sécurité alimentaire} : la confiture, riche en énergie, est disponible toute l'année, y compris hors saison.
\end{itemize}

\textbf{Conclusion.} Avec un bénéfice brut d'environ \num{1240000}~FCFA pour \SI{500}{kg} de mangues, la transformation artisanale est une voie concrète de lutte contre les pertes post-récolte et la pauvreté rurale, à condition de respecter les règles d'hygiène et de qualité.
\end{exoresolu}

\begin{attention}[Les 5 erreurs qui coûtent des points au Bac]
\begin{enumerate}
\item \textbf{Confondre transformation et conservation.} La transformation change la nature du produit (pulpe $\rightarrow$ jus, confiture, concentré) ; la conservation maintient le fruit en l'état (froid, chaleur, séchage). Une même opération (pasteurisation, séchage) peut servir les deux : précise toujours de quoi tu parles.
\item \textbf{Décrire sans justifier.} Écrire « on chauffe le jus » sans dire que la chaleur détruit les micro-organismes et dénature les enzymes fait perdre la moitié des points : chaque étape du protocole doit être reliée à l'agent d'altération qu'elle neutralise.
\item \textbf{Affirmer que le froid ou le sucre « tuent » les microbes.} La réfrigération et le sucre \emph{bloquent} le développement microbien (ralentissement, plasmolyse) sans stériliser : dès le retour à température ambiante ou après dilution, l'altération reprend. Seule la chaleur suffisante détruit les germes.
\item \textbf{Erreurs de calcul sur les pourcentages successifs.} Appliquer les pertes l'une après l'autre sur la masse \emph{restante} (ex. $-10\,\%$ puis $-8\,\%$ ne font pas $-18\,\%$), et ne jamais confondre masse d'eau évaporée et masse finale. Pose chaque étape avec son calcul et son unité.
\item \textbf{Oublier les conditions d'application.} Un protocole sans paramètres chiffrés (température de pasteurisation, ratio sucre/pulpe, épaisseur des tranches, emballage étanche, refroidissement avant ensachage) est incomplet ; de même, oublier le tri initial des fruits avariés invalide toute la chaîne de conservation.
\end{enumerate}
\end{attention}

\newpage

\coursec{Exercices}

\courssub{Je vérifie mes connaissances}

\begin{exercice}[QCM : les principes de conservation]
Pour chaque question, une seule réponse est exacte. Recopie la lettre correspondante et justifie brièvement ton choix.
\begin{enumerate}
\item Le brunissement des tranches de mangue exposées à l'air est dû principalement :
\begin{itemize}
\item[a)] à l'action des levures ;
\item[b)] à l'action d'enzymes (polyphénoloxydases) en présence d'oxygène ;
\item[c)] à la chaleur ambiante ;
\item[d)] à la teneur en eau du fruit.
\end{itemize}
\item La pasteurisation d'un jus de mangue consiste à :
\begin{itemize}
\item[a)] détruire tous les micro-organismes et leurs spores ;
\item[b)] détruire les micro-organismes pathogènes et altérants par un chauffage modéré suivi d'un refroidissement rapide ;
\item[c)] congeler le jus ;
\item[d)] ajouter du sucre au jus.
\end{itemize}
\item Le séchage conserve les fruits parce qu'il :
\begin{itemize}
\item[a)] abaisse l'activité de l'eau ($a_w$), ce qui bloque le développement microbien ;
\item[b)] tue toutes les bactéries par le froid ;
\item[c)] stérilise le fruit ;
\item[d)] augmente la teneur en sucre du fruit.
\end{itemize}
\item Dans la fabrication d'une confiture, le sucre joue le rôle de :
\begin{itemize}
\item[a)] simple édulcorant ;
\item[b)] agent conservateur qui rend l'eau indisponible pour les micro-organismes ;
\item[c)] agent colorant ;
\item[d)] enzyme de maturation.
\end{itemize}
\end{enumerate}
\end{exercice}

\begin{exercice}[Vrai ou faux justifié]
Indique si chaque affirmation est vraie ou fausse, puis justifie ta réponse en une ou deux phrases.
\begin{enumerate}
\item La stérilisation et la pasteurisation sont deux traitements thermiques équivalents.
\item Le citronnage des tranches de mangue ralentit le brunissement enzymatique.
\item La congélation détruit définitivement tous les micro-organismes présents dans un fruit.
\item Le blanchiment (ébouillantage) des tomates avant transformation inactive leurs enzymes.
\item Un fruit récolté à maturité complète se conserve plus longtemps qu'un fruit récolté au stade de début de maturation.
\end{enumerate}
\end{exercice}

\begin{exercice}[Définitions]
Définis en deux ou trois lignes chacun des termes suivants, en donnant un exemple tiré de la transformation de la mangue ou de la tomate :
\begin{enumerate}
\item perte post-récolte ;
\item pasteurisation ;
\item activité de l'eau ($a_w$) ;
\item brunissement enzymatique ;
\item appertisation.
\end{enumerate}
\end{exercice}

\begin{exercice}[Calculs directs]
Un groupe de femmes transformatrices de Garoua récolte \SI{200}{kg} de mangues fraîches.
\begin{enumerate}
\item Après épluchage et dénoyautage, il reste \SI{130}{kg} de pulpe. Calcule le rendement en pulpe (en \%).
\item Après séchage, la pulpe fournit \SI{26}{kg} de mangue séchée. Calcule le rendement global en mangue séchée par rapport à la masse de fruits frais (en \%).
\item Sachant que la mangue fraîche se vend \SI{150}{FCFA/kg} et la mangue séchée \SI{2500}{FCFA/kg}, calcule la valeur marchande des \SI{200}{kg} vendus frais, puis celle des \SI{26}{kg} de mangue séchée. Conclus en une phrase.
\end{enumerate}
\end{exercice}

\courssub{Je m'entraîne}

\begin{exercice}[Classer les techniques de conservation]
On donne cinq techniques de conservation : congélation, confiture, séchage, vinaigrage (conservation dans le vinaigre), appertisation.
\begin{enumerate}
\item Pour chaque technique, indique le ou les facteurs d'altération qu'elle neutralise (micro-organismes, enzymes, oxygène, eau disponible) et le principe physique, chimique ou biologique mis en jeu.
\item Présente tes réponses sous forme d'un tableau à trois colonnes : technique ; facteur neutralisé ; principe.
\item Cite, pour chaque technique, un fruit de saison camerounais auquel elle peut s'appliquer.
\end{enumerate}
\end{exercice}

\begin{exercice}[Organigramme du concentré de tomate]
Voici, dans le désordre, les étapes de fabrication d'un concentré de tomate :
\begin{center}
mise en boîte et fermeture hermétique — tri et lavage — concentration par évaporation — traitement thermique de conservation — broyage et raffinage — refroidissement et stockage
\end{center}
\begin{enumerate}
\item Remets ces étapes dans l'ordre chronologique correct.
\item Pour chaque étape, précise en une phrase le principe biologique, physique ou chimique mis en jeu.
\item Explique pourquoi la concentration par évaporation améliore à la fois la conservation et le coût de transport du produit.
\end{enumerate}
\end{exercice}

\begin{exercice}[Choisir le bon traitement thermique]
Une unité de transformation de Yaoundé produit du jus de mangue conditionné en bouteilles de verre. On dispose des données suivantes sur la flore microbienne du jus :
\begin{itemize}
\item les levures et moisissures sont détruites à \SI{65}{\celsius} en \SI{2}{min} ;
\item les bactéries pathogènes sont détruites à \SI{75}{\celsius} en \SI{1}{min} ;
\item le jus a un pH de \num{4,0} (milieu acide).
\end{itemize}
\begin{enumerate}
\item Explique pourquoi une pasteurisation à \SI{85}{\celsius} pendant \SI{5}{min} est suffisante pour ce jus.
\item Pourquoi une stérilisation à \SI{121}{\celsius} serait-elle ici inutile et même défavorable à la qualité du produit ?
\item Rappelle la différence fondamentale entre pasteurisation et stérilisation.
\end{enumerate}
\end{exercice}

\begin{exercice}[Étude de cas : séchage de la mangue à Maroua]
Une coopérative de Maroua transforme chaque semaine \SI{500}{kg} de mangues fraîches achetées \SI{100}{FCFA/kg}. Le séchage donne un rendement de \SI{12}{\percent} en mangue séchée, vendue \SI{3000}{FCFA/kg}. Les frais de transformation (énergie, emballage, main-d'œuvre) s'élèvent à \SI{20000}{FCFA} par semaine.
\begin{enumerate}
\item Calcule la masse hebdomadaire de mangue séchée obtenue.
\item Calcule le chiffre d'affaires hebdomadaire de la coopérative.
\item Calcule le bénéfice hebdomadaire, puis compare-le au bénéfice qu'aurait procuré la vente directe des fruits frais (marge de \SI{50}{FCFA/kg}).
\item Dégage deux avantages socio-économiques supplémentaires de la transformation locale des fruits de saison.
\end{enumerate}
\end{exercice}

\courssub{Je me prépare au Bac}

\begin{exercice}[Type Bac : le brunissement des tranches de mangue]
Au laboratoire d'un lycée de Douala, on réalise l'expérience suivante : des tranches de mangue de même variété et de même épaisseur sont réparties en quatre lots, puis exposées à l'air pendant \SI{6}{h}. L'intensité du brunissement est notée de 0 (aucun) à 4 (brunissement total). Les résultats sont consignés dans le tableau ci-dessous.

\begin{center}
\begin{tabularx}{\linewidth}{|l|X|c|}
\hline
\rowcolor{popblueL} Lot & Traitement & Intensité du brunissement après \SI{6}{h} \\
\hline
A (témoin) & Aucun & 4 \\
\hline
B & Tranches arrosées de jus de citron & 1 \\
\hline
C & Tranches blanchies \SI{2}{min} dans l'eau bouillante & 0 \\
\hline
D & Tranches placées au réfrigérateur à \SI{4}{\celsius} & 2 \\
\hline
\end{tabularx}
\end{center}

\begin{enumerate}
\item Identifie la variable expérimentale (facteur modifié) et la variable mesurée dans cette expérience. Quel est le rôle du lot A ?
\item Formule l'hypothèse que cette expérience permet de tester.
\item Interprète les résultats des lots B, C et D en faisant intervenir la nature enzymatique du brunissement.
\item Déduis de cette expérience trois principes de conservation applicables à la transformation de la mangue, en précisant pour chacun le facteur d'altération neutralisé.
\item Propose une expérience complémentaire permettant de vérifier que l'oxygène de l'air intervient dans le brunissement ; décris le protocole et le résultat attendu.
\end{enumerate}
\end{exercice}

\begin{exercice}[Type Bac : choix d'un barème de pasteurisation pour un jus de mangue]
Une PME agroalimentaire de Bafoussam souhaite conditionner un jus de mangue. Elle fait réaliser une étude de survie d'une bactérie pathogène de référence dans le jus, à différentes températures. Le temps nécessaire pour diviser par $10^{6}$ la population bactérienne (temps de réduction décimal à 6 cycles, $6D$) est mesuré :

\begin{center}
\begin{tabular}{|c|c|c|c|c|}
\hline
\rowcolor{popblueL} Température (\si{\celsius}) & 60 & 70 & 80 & 90 \\
\hline
Temps $6D$ (min) & 60 & 6 & \num{0,6} & \num{0,06} \\
\hline
\end{tabular}
\end{center}

\begin{enumerate}
\item Trace, sur papier millimétré ou à l'aide d'un repère semi-logarithmique, la courbe du temps $6D$ en fonction de la température (échelle logarithmique pour le temps).
\item À partir de la courbe, détermine graphiquement le temps $6D$ à \SI{85}{\celsius}.
\item L'entreprise choisit un barème de \SI{85}{\celsius} pendant \SI{5}{min}. Montre, à l'aide du graphique et des données, que ce barème offre une marge de sécurité confortable.
\item Explique pourquoi l'acidité naturelle du jus de mangue (pH $\approx$ 4) renforce l'efficacité de ce traitement thermique.
\item Distingue pasteurisation et stérilisation (températures, objectifs, effets sur le produit), puis justifie le choix de la pasteurisation pour ce jus.
\item Cite deux précautions à prendre après le traitement thermique pour garantir la conservation du jus conditionné.
\end{enumerate}
\end{exercice}

\begin{exercice}[Type Bac : filière tomate et pertes post-récolte dans la plaine des Mbo]
Dans la plaine des Mbo (région du Littoral), la production de tomates fraîches atteint \SI{4000}{tonnes} par campagne, concentrée sur trois mois. On estime que \SI{35}{\percent} de la récolte est perdue avant consommation. Une coopérative envisage d'installer une unité de fabrication de concentré de tomate. Données :
\begin{itemize}
\item le rendement de transformation est de \SI{6}{kg} de tomates fraîches pour \SI{1}{kg} de concentré ;
\item le concentré se vend \SI{1800}{FCFA/kg} ; la tomate fraîche se vend en moyenne \SI{200}{FCFA/kg} en pleine saison ;
\item l'unité peut traiter \SI{300}{tonnes} de tomates fraîches par campagne ;
\item les charges de transformation représentent \SI{40}{\percent} du chiffre d'affaires.
\end{itemize}

\begin{enumerate}
\item Calcule la masse de tomates perdue chaque campagne dans la plaine.
\item Explique, en t'appuyant sur tes connaissances, trois causes biologiques ou physiques de ces pertes post-récolte chez la tomate (fruit climactérique, riche en eau).
\item Décris l'organigramme complet de fabrication du concentré de tomate (au moins six étapes ordonnées), en précisant pour deux étapes au choix le principe scientifique mis en jeu.
\item Calcule la masse de concentré produite par campagne, le chiffre d'affaires et le bénéfice de la coopérative.
\item Compare la valeur marchande des \SI{300}{tonnes} vendues fraîches à celle du concentré obtenu, puis rédige un paragraphe de conclusion (8 à 10 lignes) sur l'intérêt de la transformation et de la conservation des fruits de saison pour la sécurité alimentaire et l'économie rurale au Cameroun.
\end{enumerate}
\end{exercice}

\newpage

\coursec{Corrigés des exercices}

\begin{corrige}[Exercice 1 — QCM : les principes de conservation]
\begin{enumerate}
\item \textbf{Réponse b}. Le brunissement est un phénomène enzymatique : les polyphénoloxydases (PPO), libérées lors de la déstructuration des cellules, oxydent les composés phénoliques du fruit en présence de l'oxygène de l'air, formant des pigments bruns (mélanines). Ni les levures, ni la chaleur seule, ni la teneur en eau ne sont en cause.
\item \textbf{Réponse b}. La pasteurisation est un chauffage modéré (typiquement $60$ à \SI{95}{\celsius}) suivi d'un refroidissement rapide : elle détruit les micro-organismes pathogènes et altérants \emph{sans} détruire toutes les spores (c'est la différence avec la stérilisation). La congélation (c) et l'ajout de sucre (d) sont d'autres techniques de conservation.
\item \textbf{Réponse a}. Le séchage élimine une grande partie de l'eau du fruit, ce qui abaisse l'activité de l'eau ($a_w$) en dessous du seuil de développement des micro-organismes (en général $a_w < \num{0,6}$). Le séchage ne stérilise pas le fruit et n'utilise pas le froid.
\item \textbf{Réponse b}. À forte concentration (environ $60$ à \SI{65}{\percent} d'extrait sec dans une confiture), le sucre exerce une pression osmotique élevée qui rend l'eau indisponible pour les micro-organismes : c'est un véritable agent conservateur, pas seulement un édulcorant.
\end{enumerate}
\end{corrige}

\begin{attention}[Pièges du QCM]
Ne confonds pas \emph{pasteurisation} (modérée, $< \SI{100}{\celsius}$, spores survivantes) et \emph{stérilisation} (forte, $\geq \SI{115}{\celsius}$ sous pression, spores détruites). Le brunissement enzymatique nécessite trois acteurs : enzyme (PPO), substrat (phénols) et oxygène ; supprimer l'un des trois l'empêche. Le séchage abaisse l'$a_w$, il ne la supprime pas totalement ($a_w > 0$).
\end{attention}

\begin{corrige}[Exercice 2 — Vrai ou faux justifié]
\begin{enumerate}
\item \textbf{Faux.} La pasteurisation (chauffage modéré, $60$ à \SI{95}{\celsius}) détruit les formes végétatives des micro-organismes pathogènes et altérants mais épargne certaines spores ; la stérilisation (environ \SI{121}{\celsius} sous pression) détruit tous les micro-organismes \emph{et} leurs spores. Ce ne sont donc pas des traitements équivalents.
\item \textbf{Vrai.} Le jus de citron apporte de l'acide citrique (abaissement du pH) et de la vitamine C (antioxydant), deux facteurs qui inhibent l'activité des polyphénoloxydases responsables du brunissement.
\item \textbf{Faux.} La congélation bloque le développement microbien mais ne tue pas tous les micro-organismes : beaucoup survivent à l'état dormant et reprennent leur activité à la décongélation. Le froid est bactériostatique, non bactéricide.
\item \textbf{Vrai.} Le blanchiment (immersion de $1$ à \SI{3}{min} dans l'eau bouillante) dénature les protéines, donc inactive les enzymes (PPO, pectinases) responsables du brunissement et de la dégradation de la texture ; il facilite aussi l'épluchage.
\item \textbf{Faux.} Un fruit à maturité complète a une respiration intense, une chair ramollie et une sensibilité maximale aux chocs et aux germes : il se conserve mal. On récolte au stade de début de maturation (stade « mûr-vert » ou « tournant ») pour allonger la durée de conservation.
\end{enumerate}
\end{corrige}

\begin{attention}[Vocabulaire de rigueur]
« La congélation \emph{tue} les microbes » $\rightarrow$ \textbf{Erreur}. Écris : « La congélation \emph{bloque} ou \emph{inhibe} le développement microbien ». « Pasteurisation = stérilisation douce » $\rightarrow$ \textbf{Imprécis}. Ce sont deux opérations distinctes aux objectifs et résultats microbiologiques différents.
\end{attention}

\begin{corrige}[Exercice 3 — Définitions]
\begin{enumerate}
\item \textbf{Perte post-récolte} : quantité de fruits (ou dégradation de leur qualité) perdue entre la récolte et la consommation, par pourrissement microbien, respiration, blessures ou mauvaise manutention. \emph{Exemple :} environ $30$ à \SI{40}{\percent} des mangues produites dans l'Extrême-Nord sont perdues en pleine saison faute de transformation.
\item \textbf{Pasteurisation} : traitement thermique modéré (ex. \SI{85}{\celsius} pendant \SI{5}{min}) suivi d'un refroidissement rapide, qui détruit les micro-organismes pathogènes et altérants sans stériliser le produit. \emph{Exemple :} pasteurisation du jus de mangue en bouteille.
\item \textbf{Activité de l'eau ($a_w$)} : fraction d'eau « disponible » pour les réactions chimiques et les micro-organismes dans un aliment, comprise entre 0 et 1. En dessous de $a_w \approx \num{0,6}$, aucun micro-organisme ne se développe. \emph{Exemple :} la mangue séchée ($a_w$ faible) se conserve plusieurs mois.
\item \textbf{Brunissement enzymatique} : réaction d'oxydation des composés phénoliques par les polyphénoloxydases en présence d'oxygène, produisant des pigments bruns (mélanines). \emph{Exemple :} brunissement des tranches de mangue coupées exposées à l'air.
\item \textbf{Appertisation} : conservation par stérilisation d'un aliment conditionné dans un récipient hermétique (boîte, bocal), chauffé vers \SI{121}{\celsius} pour détruire micro-organismes et spores. \emph{Exemple :} concentré de tomate en boîte stérilisée.
\end{enumerate}
\end{corrige}

\begin{attention}[Rigueur attendue au Bac]
Une définition doit être \emph{autosuffisante} : elle nomme le concept, en donne la nature (phénomène, traitement, grandeur), le mécanisme ou le principe, les valeurs seuils s'il y a lieu. L'exemple contextualisé (Cameroun si possible) valide la compréhension.
\end{attention}

\begin{corrige}[Exercice 4 — Calculs directs]
\begin{enumerate}
\item Rendement en pulpe :
\[ R_{\text{pulpe}} = \frac{130}{200}\times 100 = \SI{65}{\percent}. \]
\item Rendement global en mangue séchée par rapport aux fruits frais :
\[ R_{\text{sec}} = \frac{26}{200}\times 100 = \SI{13}{\percent}. \]
(Remarque : par rapport à la pulpe, le rendement de séchage est $\frac{26}{130}\times 100 = \SI{20}{\percent}$.)
\item Valeur marchande :
\begin{itemize}
\item Vente fraîche : $200 \times 150 = \SI{30000}{FCFA}$.
\item Vente séchée : $26 \times 2500 = \SI{65000}{FCFA}$.
\end{itemize}
\textbf{Conclusion :} la transformation en mangue séchée rapporte $\num{65000} - \num{30000} = \SI{35000}{FCFA}$ de plus, soit plus du double de la valeur de la vente fraîche : la transformation ajoute une forte valeur ajoutée tout en réduisant les pertes.
\end{enumerate}
\end{corrige}

\begin{attention}[Unités et conversions]
Vérifie toujours l'unité du prix (FCFA/kg ou FCFA/tonne) et la masse correspondante. Ici prix au kg $\times$ masse en kg. Ne mélange pas tonnes et kilogrammes sans conversion explicite ($1\ \text{tonne} = 1000\ \text{kg}$).
\end{attention}

\begin{corrige}[Exercice 5 — Classer les techniques de conservation]
\begin{enumerate}
\item[2.] Tableau récapitulatif (questions 1 et 3 incluses) :
\end{enumerate}
\begin{center}
\begin{tabularx}{\linewidth}{|l|X|X|l|}
\hline
\rowcolor{popblueL} Technique & Facteur(s) neutralisé(s) & Principe & Exemple de fruit \\
\hline
Congélation & Micro-organismes (développement bloqué), enzymes (ralenties) & Physique : le froid ralentit les réactions biochimiques et cristallise l'eau & Mangue (pulpe congelée) \\
\hline
Confiture & Eau disponible (micro-organismes privés d'eau) & Physico-chimique : forte concentration en sucre $\Rightarrow$ pression osmotique élevée, $a_w$ abaissée & Mangue, goyave \\
\hline
Séchage & Eau disponible & Physique : évaporation de l'eau, $a_w < \num{0,6}$ & Mangue, banane plantain \\
\hline
Vinaigrage & Micro-organismes (inhibés par l'acidité) & Chimique : abaissement du pH par l'acide acétique & Piment, jeunes mangues (pickles) \\
\hline
Appertisation & Micro-organismes et spores, enzymes, oxygène (boîte hermétique) & Physique : stérilisation à \SI{121}{\celsius} en récipient clos & Tomate (concentré en boîte) \\
\hline
\end{tabularx}
\end{center}
\end{corrige}

\begin{attention}[Point de vigilance]
La congélation ne \emph{détruit} pas les micro-organismes : elle bloque seulement leur développement. Écris « neutralise » ou « bloque », pas « tue ». De même, le séchage n'élimine pas \emph{toute} l'eau ($a_w > 0$), il la rend indisponible.
\end{attention}

\begin{corrige}[Exercice 6 — Organigramme du concentré de tomate]
\begin{enumerate}
\item Ordre chronologique correct :
\begin{enumerate}
\item tri et lavage ;
\item broyage et raffinage ;
\item concentration par évaporation ;
\item mise en boîte et fermeture hermétique ;
\item traitement thermique de conservation ;
\item refroidissement et stockage.
\end{enumerate}
\item Principes mis en jeu :
\begin{itemize}
\item \textbf{Tri et lavage} : élimination des fruits avariés et de la charge microbienne de surface (principe hygiénique).
\item \textbf{Broyage et raffinage} : libération de la pulpe et séparation des peaux et graines (principe mécanique) ; le chauffage éventuel à ce stade (\"hot break\") inactive les enzymes (pectinases).
\item \textbf{Concentration par évaporation} : élimination d'une partie de l'eau $\Rightarrow$ abaissement de l'$a_w$ et augmentation de l'extrait sec (principe physique).
\item \textbf{Mise en boîte hermétique} : protection contre l'oxygène et les recontaminations (barrière physique).
\item \textbf{Traitement thermique} : destruction des micro-organismes et de leurs spores, inactivation des enzymes (appertisation).
\item \textbf{Refroidissement et stockage} : évite la cuisson excessive et limite les réactions chimiques résiduelles.
\end{itemize}
\item La concentration retire de l'eau : l'$a_w$ diminue, ce qui freine le développement microbien et améliore la conservation. Comme la masse et le volume du produit diminuent fortement (environ \SI{6}{kg} de tomates pour \SI{1}{kg} de concentré), les coûts de transport et de stockage par unité de matière sèche sont fortement réduits.
\end{enumerate}
\end{corrige}

\begin{attention}[Ordre des opérations]
Le broyage/raffinage précède la concentration. Si l'on concentre avant de raffiner, les peaux et graines brûlent et bouchent les échangeurs. Le traitement thermique final se fait \emph{après} la mise en boîte (appertisation), pas avant.
\end{attention}

\begin{corrige}[Exercice 7 — Choisir le bon traitement thermique]
\begin{enumerate}
\item Le barème le plus exigeant est celui des bactéries pathogènes : \SI{75}{\celsius} pendant \SI{1}{min}. Le traitement choisi (\SI{85}{\celsius} pendant \SI{5}{min}) est très supérieur en température et en durée : il détruit donc largement les levures, moisissures et bactéries pathogènes. De plus, le pH de \num{4,0} (milieu acide) inhibe la germination des spores et augmente la thermosensibilité des micro-organismes : la pasteurisation suffit.
\item Une stérilisation à \SI{121}{\celsius} est inutile car, en milieu acide (pH $< \num{4,5}$), les spores de bactéries dangereuses (ex. \emph{Clostridium botulinum}) ne peuvent pas germer. Elle serait défavorable : chauffage prolongé $\Rightarrow$ destruction des vitamines (vitamine C), perte d'arômes, brunissement non enzymatique (réaction de Maillard) et coût énergétique plus élevé.
\item \textbf{Différence fondamentale :} la pasteurisation (chauffage modéré, $< \SI{100}{\celsius}$) détruit les formes végétatives pathogènes et altérantes mais pas toutes les spores ; le produit doit souvent être conservé au frais ou être naturellement acide. La stérilisation ($\geq \SI{115}{\celsius}$ sous pression) détruit tous les micro-organismes et leurs spores : le produit est stable à température ambiante, mais sa qualité nutritionnelle et organoleptique est davantage altérée.
\end{enumerate}
\end{corrige}

\begin{attention}[Argumentaire « pH < 4,5 »]
C'est l'argument clé pour justifier une pasteurisation au lieu d'une stérilisation sur produit acide. \emph{Clostridium botulinum} ne germe pas en dessous de pH 4,5 : le danger botulique est écarté, il ne reste que les formes végétatives (tuées par pasteurisation). Cite ce seuil explicitement.
\end{attention}

\begin{corrige}[Exercice 8 — Étude de cas : séchage de la mangue à Maroua]
\begin{enumerate}
\item Masse hebdomadaire de mangue séchée :
\[ m = 500 \times \frac{12}{100} = \SI{60}{kg}. \]
\item Chiffre d'affaires hebdomadaire :
\[ CA = 60 \times 3000 = \SI{180000}{FCFA}. \]
\item Bénéfice hebdomadaire de la transformation :
\begin{align*}
\text{Charges} &= \text{achat des fruits} + \text{frais de transformation} \\
&= 500 \times 100 + \num{20000} = \num{50000} + \num{20000} = \SI{70000}{FCFA}. \\
\text{Bénéfice} &= \num{180000} - \num{70000} = \SI{110000}{FCFA}.
\end{align*}
Vente directe des fruits frais : marge de $\SI{50}{FCFA/kg} \times 500 = \SI{25000}{FCFA}$.
La transformation procure un bénéfice $110\,000 / 25\,000 = \num{4,4}$ fois supérieur à la vente directe.
\item Avantages socio-économiques supplémentaires (deux parmi) :
\begin{itemize}
\item création d'emplois locaux et de revenus pour les femmes et les jeunes ;
\item réduction des pertes post-récolte et disponibilité des fruits hors saison (sécurité alimentaire) ;
\item possibilité d'exportation et d'entrées de devises ; développement de l'artisanat et des PME rurales.
\end{itemize}
\end{enumerate}
\end{corrige}

\begin{attention}[Chiffre d'affaires vs Bénéfice]
$CA = \text{Prix de vente} \times \text{Quantité vendue}$. $\text{Bénéfice} = CA - \text{Charges totales}$. Ne confonds pas marge unitaire (prix vente - prix achat) et bénéfice net (après frais de transformation, énergie, main-d'œuvre, transport...).
\end{attention}

\begin{corrige}[Exercice 9 — Type Bac : le brunissement des tranches de mangue]
\emph{Barème indicatif : question 1 : 2 pts ; question 2 : 1 pt ; question 3 : 3 pts ; question 4 : 2 pts ; question 5 : 2 pts. Total : 10 pts.}
\begin{enumerate}
\item \textbf{Variable expérimentale (facteur modifié) :} le traitement appliqué aux tranches (citronnage, blanchiment, réfrigération, absence de traitement). \textbf{Variable mesurée :} l'intensité du brunissement (échelle de 0 à 4). Le lot A est le \cle{lot témoin} : il sert de référence pour comparer l'effet de chaque traitement.
\item \textbf{Hypothèse :} le brunissement de la mangue est une réaction enzymatique ; tout traitement qui inactive les enzymes ou ralentit leur activité (chaleur, acidification, froid) réduit le brunissement.
\item \textbf{Interprétation :}
\begin{itemize}
\item \textbf{Lot C (blanchiment, intensité 0) :} l'eau bouillante dénature les polyphénoloxydases (protéines thermosensibles) : la réaction d'oxydation est totalement bloquée. C'est le traitement le plus efficace.
\item \textbf{Lot B (citron, intensité 1) :} l'acide citrique abaisse le pH en dehors de l'optimum des PPO et la vitamine C réduit les quinones formées : l'activité enzymatique est fortement ralentie mais pas supprimée.
\item \textbf{Lot D (réfrigération, intensité 2) :} le froid ralentit la vitesse des réactions enzymatiques sans dénaturer les enzymes : le brunissement est seulement retardé.
\end{itemize}
Comparaison au témoin A (intensité 4) : les trois traitements réduisent le brunissement, ce qui confirme sa nature enzymatique.
\item Trois principes de conservation :
\begin{itemize}
\item \textbf{traitement thermique (blanchiment)} $\Rightarrow$ neutralise les \emph{enzymes} ;
\item \textbf{acidification (citronnage)} $\Rightarrow$ neutralise les \emph{enzymes} (et inhibe certains micro-organismes) ;
\item \textbf{réfrigération} $\Rightarrow$ ralentit les \emph{enzymes} et le développement \emph{microbien}.
\end{itemize}
\item \textbf{Expérience complémentaire :} on prépare deux lots de tranches de mangue identiques ; le lot 1 est exposé à l'air libre, le lot 2 est placé sous vide (ou immergé dans de l'eau bouillie refroidie / sous film hermétique) dans les mêmes conditions de température et de durée. \textbf{Résultat attendu :} le lot 1 brunit (intensité élevée), le lot 2 ne brunit pas ou très peu. Conclusion : l'oxygène de l'air est bien un réactif nécessaire au brunissement enzymatique.
\end{enumerate}
\end{corrige}

\begin{attention}[Points de vigilance Bac]
Cite toujours le lot témoin et son rôle ; relie chaque résultat à l'activité des polyphénoloxydases ; dans la question 5, précise que toutes les autres conditions (température, durée, variété) doivent être identiques (un seul facteur modifié).
\end{attention}

\begin{corrige}[Exercice 10 — Type Bac : choix d'un barème de pasteurisation]
\emph{Barème indicatif : question 1 : 2 pts ; question 2 : 1 pt ; question 3 : 2 pts ; question 4 : 1 pt ; question 5 : 2 pts ; question 6 : 2 pts. Total : 10 pts.}
\begin{enumerate}
\item On porte la température en abscisse (60 à \SI{90}{\celsius}) et le temps $6D$ en ordonnée logarithmique (60 ; 6 ; \num{0,6} ; \num{0,06} min). Les points s'alignent sur une droite décroissante : chaque fois que la température augmente de \SI{10}{\celsius}, le temps $6D$ est divisé par 10.
\begin{popfigure}
\begin{tikzpicture}
\begin{semilogyaxis}[
width=0.85\linewidth, height=6cm,
xlabel={Température (\si{\celsius})}, ylabel={Temps $6D$ (min, échelle log)},
xmin=55, xmax=95, ymin=0.01, ymax=200,
xtick={60,70,80,90}, grid=both, grid style={popline!40},
]
\addplot[popcurve, mark=*, mark options={fill=poporange}] coordinates {(60,60) (70,6) (80,0.6) (90,0.06)};
\addplot[dashed, popgreen] coordinates {(85,0.01) (85,0.18)};
\addplot[dashed, popgreen] coordinates {(55,0.18) (85,0.18)};
\node[popgreen, anchor=south west] at (axis cs:56,0.22) {\scriptsize $6D \approx \num{0,18}$ min à \SI{85}{\celsius}};
\end{semilogyaxis}
\end{tikzpicture}
\legende{Courbe du temps de réduction $6D$ en fonction de la température (échelle semi-logarithmique). Lecture graphique à \SI{85}{\celsius} : $6D \approx \num{0,18}$ min.}
\end{popfigure}
\item Par interpolation sur la droite (à mi-chemin entre 80 et \SI{90}{\celsius} en échelle log) : $6D(\SI{85}{\celsius}) \approx \num{0,18}$ min, soit environ \SI{11}{s}.
\item Le barème choisi est \SI{85}{\celsius} pendant \SI{5}{min}, alors que \SI{0,18}{min} suffisent pour une réduction de $10^{6}$. La marge de sécurité est :
\[ \frac{5}{\num{0,18}} \approx 28. \]
Le traitement réalisé est environ 28 fois supérieur au minimum requis, soit une réduction théorique de plus de $10^{160}$ : la marge est très confortable et couvre les hétérogénéités de chauffage.
\item À pH $\approx 4$, les micro-organismes sont plus thermosensibles (la chaleur les détruit plus vite en milieu acide) et surtout les spores de bactéries pathogènes comme \emph{Clostridium botulinum} ne peuvent pas germer en dessous de pH $\num{4,5}$ : le traitement thermique modéré devient suffisant et sûr.
\item \textbf{Pasteurisation :} chauffage modéré (ici \SI{85}{\celsius}) détruisant les formes végétatives pathogènes et altérantes ; qualité gustative et nutritionnelle bien préservée. \textbf{Stérilisation :} chauffage à $\geq \SI{115}{\celsius}$ détruisant tous les micro-organismes et leurs spores ; produit stable à l'ambiante mais vitamines et arômes dégradés. Pour ce jus acide, la pasteurisation suffit (sécurité assurée par le pH) et préserve la vitamine C et les arômes de la mangue : c'est le meilleur compromis sécurité/qualité.
\item Deux précautions après traitement :
\begin{itemize}
\item remplissage et bouchage \emph{à chaud} (ou traitement en bouteille fermée) puis refroidissement rapide pour éviter toute recontamination et créer une dépression assurant l'étanchéité ;
\item contrôle de l'étanchéité des capsules et stockage à l'abri de la lumière et de la chaleur (idéalement au frais), avec vérification de l'absence de bombage ou de fermentation avant consommation.
\end{itemize}
\end{enumerate}
\end{corrige}

\begin{attention}[Points de vigilance]
En échelle semi-logarithmique, l'interpolation à \SI{85}{\celsius} se fait au \emph{milieu géométrique} ($\sqrt{\num{0,6}\times\num{0,06}} = \sqrt{\num{0,036}} \approx \num{0,19}$), pas à la moyenne arithmétique. Justifie toujours le choix du barème par le pH $< \num{4,5}$.
\end{attention}

\begin{corrige}[Exercice 11 — Type Bac : filière tomate et pertes post-récolte]
\emph{Barème indicatif : question 1 : 1 pt ; question 2 : 3 pts ; question 3 : 3 pts ; question 4 : 2 pts ; question 5 : 3 pts. Total : 12 pts.}
\begin{enumerate}
\item Masse de tomates perdue par campagne :
\[ m_{\text{perdue}} = 4000 \times \frac{35}{100} = \SI{1400}{tonnes}. \]
\item Trois causes de pertes post-récolte chez la tomate :
\begin{itemize}
\item \textbf{Respiration intense et production d'éthylène} (fruit climactérique) : la maturation se poursuit et s'accélère après récolte, conduisant au ramollissement puis à la sénescence ;
\item \textbf{forte teneur en eau} (environ \SI{90}{\percent}) et peau fine : le fruit est très sensible aux chocs et blessures pendant la manutention et le transport (pistes dégradées), portes d'entrée des germes ;
\item \textbf{développement microbien} (moisissures, bactéries de pourriture molle) favorisé par la chaleur ambiante et l'humidité, en l'absence de chaîne du froid.
\end{itemize}
\item Organigramme de fabrication du concentré de tomate :
\begin{center}
réception $\rightarrow$ tri et lavage $\rightarrow$ ébouillantage/épluchage $\rightarrow$ broyage et raffinage $\rightarrow$ concentration par évaporation $\rightarrow$ mise en boîte hermétique $\rightarrow$ stérilisation (appertisation) $\rightarrow$ refroidissement $\rightarrow$ stockage.
\end{center}
Principes (deux exemples) :
\begin{itemize}
\item \textbf{concentration par évaporation} : élimination d'eau $\Rightarrow$ abaissement de l'activité de l'eau ($a_w$), ce qui freine le développement microbien ;
\item \textbf{stérilisation en boîte fermée} : destruction des micro-organismes et de leurs spores par la chaleur (\SI{121}{\celsius}) et protection contre l'oxygène et les recontaminations.
\end{itemize}
\item Calculs :
\begin{align*}
\text{Masse de concentré} &= \frac{300}{6} = \SI{50}{tonnes} = \SI{50000}{kg}. \\
\text{Chiffre d'affaires} &= 50\,000 \times 1800 = \SI{90000000}{FCFA}. \\
\text{Charges} &= 0{,}40 \times 90\,000\,000 = \SI{36000000}{FCFA}. \\
\text{Bénéfice} &= 90\,000\,000 - 36\,000\,000 = \SI{54000000}{FCFA}.
\end{align*}
\item Vente fraîche des \SI{300}{tonnes} : $300\,000 \times 200 = \SI{60000000}{FCFA}$ de valeur brute, contre \SI{90000000}{FCFA} pour le concentré, soit \SI{30000000}{FCFA} de valeur ajoutée brute (et un bénéfice de 54 millions après charges).

\textbf{Conclusion (paragraphe attendu) :} la transformation des fruits de saison, comme la tomate de la plaine des Mbo ou la mangue de l'Extrême-Nord, répond à un double enjeu. D'une part, elle réduit les pertes post-récolte considérables (ici \SI{1400}{tonnes}, soit \SI{35}{\percent} de la récolte) causées par la respiration du fruit, sa fragilité et les pourritures microbiennes en période d'abondance. D'autre part, elle crée de la valeur ajoutée locale : le concentré rapporte ici 30 millions de FCFA de plus que la vente fraîche et dégage un bénéfice de 54 millions de FCFA par campagne, source de revenus pour les producteurs et d'emplois ruraux, notamment pour les femmes et les jeunes. En étalant la disponibilité des produits sur toute l'année, la conservation (séchage, appertisation, pasteurisation) renforce la sécurité alimentaire et nutritionnelle des populations. Elle stabilise aussi les prix payés aux producteurs en pleine saison et peut ouvrir des débouchés à l'exportation. La transformation locale apparaît ainsi comme un levier essentiel de l'économie rurale camerounaise.
\end{enumerate}
\end{corrige}

\begin{attention}[Points de vigilance]
Distingue bien \emph{chiffre d'affaires} et \emph{bénéfice} (soustraire les charges) ; convertis les tonnes en kilogrammes avant de multiplier par le prix au kilogramme ; dans la question 2, exige-toi des causes \emph{biologiques ou physiques} précises (climactère, $a_w$, blessures), pas seulement « mauvaise conservation ».
\end{attention}

\newpage

\coursec{Fiche bilan}

\courssub{Carte mentale de la leçon}

\begin{popfigure}
\begin{tikzpicture}[mindmap, grow cyclic,
  every node/.style={concept, font=\small, text=white},
  root concept/.append style={concept color=popdark, font=\small\bfseries, minimum size=3.2cm},
  level 1/.append style={level distance=3.6cm, sibling angle=60, minimum size=2.4cm},
  level 2/.append style={level distance=2.2cm, minimum size=1.7cm, font=\scriptsize}]
\node[root concept]{Transformation et conservation des fruits de saison}
  child[concept color=popblue]{ node {Pertes post-récolte}
    child[concept color=popblue!70]{ node {30--50\,\% des fruits perdus} }
    child[concept color=popblue!70]{ node {Enjeu : sécurité alimentaire} }
  }
  child[concept color=poporange]{ node {Causes d'altération}
    child[concept color=poporange!70]{ node {Microbiennes : moisissures, levures, bactéries} }
    child[concept color=poporange!70]{ node {Enzymatiques : brunissement} }
    child[concept color=poporange!70]{ node {Physico-chimiques : eau, O$_2$, chaleur} }
  }
  child[concept color=popgreen]{ node {Principes de conservation}
    child[concept color=popgreen!70]{ node {Stopper : froid, chaleur} }
    child[concept color=popgreen!70]{ node {Réduire l'eau : séchage, sucre, sel} }
    child[concept color=popgreen!70]{ node {Exclure O$_2$ : vide, stérilisation} }
  }
  child[concept color=poppurple]{ node {Cas de la mangue}
    child[concept color=poppurple!70]{ node {Jus, nectar, confiture} }
    child[concept color=poppurple!70]{ node {Séchage (tranches)} }
  }
  child[concept color=poppink]{ node {Cas de la tomate}
    child[concept color=poppink!70]{ node {Concentré, purée, sauce} }
    child[concept color=poppink!70]{ node {Conserve stérilisée} }
  }
  child[concept color=popgold]{ node {Qualité et enjeux}
    child[concept color=popgold!70]{ node {Hygiène, contrôle qualité} }
    child[concept color=popgold!70]{ node {Revenus, emplois locaux} }
  };
\end{tikzpicture}
\legende{Carte mentale de la leçon : les six branches essentielles à maîtriser --- (1) le problème des pertes post-récolte, (2) les causes d'altération des fruits, (3) les principes généraux de conservation, (4) la transformation de la mangue, (5) la transformation de la tomate, (6) la qualité et les enjeux socio-économiques.}
\end{popfigure}

\courssub{L'essentiel à retenir}

\begin{aretenir}[Fiche bilan --- Transformation et conservation des fruits de saison]
\textbf{1. Définitions clés}
\begin{itemize}
  \item \cle{Transformation} : ensemble des opérations (lavage, tri, épluchage, cuisson, broyage, conditionnement) qui modifient un fruit frais pour obtenir un produit stable (jus, confiture, concentré).
  \item \cle{Conservation} : ensemble des techniques visant à prolonger la durée de vie d'un fruit ou d'un produit dérivé en bloquant les altérations.
  \item \cle{Pertes post-récolte} : quantités de fruits détruites ou dégradées entre la récolte et la consommation (au Cameroun, 30 à 50\,\% des mangues en pleine saison).
  \item \cle{Pasteurisation} : chauffage modéré (environ 70--90\,$^\circ$C pendant quelques secondes à quelques minutes) détruisant les micro-organismes pathogènes et inactivant la plupart des enzymes, sans éliminer toutes les spores : le produit doit donc être conservé au frais ou en milieu acide/sucré.
  \item \cle{Stérilisation} : chauffage à plus de 100\,$^\circ$C (autoclave, 115--121\,$^\circ$C) détruisant tous les micro-organismes et leurs spores.
  \item \cle{Appertisation} : conservation par stérilisation d'un aliment placé dans un récipient hermétique (boîte, bocal), du nom d'Appert.
\end{itemize}

\textbf{2. Causes d'altération et moyens de lutte}
\begin{center}
\begin{tabularx}{\linewidth}{|l|X|X|}
\hline
\rowcolor{popblueL} \textbf{Cause} & \textbf{Mécanisme} & \textbf{Technique de lutte} \\
\hline
Micro-organismes (moisissures, levures, bactéries) & Fermentations, pourriture, toxines & Chauffage (pasteurisation, stérilisation), froid, acidification, hygiène \\
\hline
Enzymes du fruit (polyphénoloxydases, pectinases) & Brunissement, ramollissement, perte d'arômes & Blanchiment (choc thermique), jus de citron (acide citrique), froid \\
\hline
Eau libre ($a_w$ élevée) & Milieu favorable aux microbes & Séchage, ajout de sucre (confiture : 60--65\,\% d'extraits secs), sel \\
\hline
Oxygène de l'air & Oxydations, développement des aérobies & Conditionnement sous vide, remplissage à chaud, fermeture hermétique \\
\hline
Lumière et chaleur & Destruction des vitamines (vitamine C), rancissement & Stockage à l'ombre, au frais, emballages opaques \\
\hline
\end{tabularx}
\end{center}

\textbf{3. Filières de transformation à connaître par c\oe ur}
\begin{itemize}
  \item \textbf{Mangue} : réception $\to$ tri et lavage $\to$ épluchage/dénoyautage $\to$ pulpage $\to$ \emph{jus/nectar} (pulpe + eau + sucre, pasteurisation, embouteillage à chaud) ; \emph{confiture} (pulpe + sucre en proportion voisine de 1:1, cuisson jusqu'à 60--65\,\% d'extraits secs, mise en pot à chaud) ; \emph{mangue séchée} (tranches fines, séchage solaire ou en étuve à 50--60\,$^\circ$C).
  \item \textbf{Tomate} : tri et lavage $\to$ ébouillantage (facilite l'épluchage et inactive les enzymes) $\to$ broyage $\to$ passage au tamis (élimination des peaux et des graines) $\to$ concentration par évaporation (concentré à 28--30\,\% de matière sèche) $\to$ remplissage à chaud en boîtes/bocaux $\to$ stérilisation (115--121\,$^\circ$C) $\to$ refroidissement rapide et stockage.
\end{itemize}

\textbf{4. Méthodes express}
\begin{itemize}
  \item Règle d'or : \textbf{toute conservation repose sur au moins un obstacle} (chaleur, froid, sucre, acide, vide) ; combiner les obstacles augmente l'efficacité (technologie des barrières).
  \item Remplissage \textbf{à chaud} ($\geq 85\,^\circ$C) + fermeture immédiate $\Rightarrow$ création d'un vide partiel au refroidissement (sécurité microbiologique).
  \item Contrôle qualité minimal : aspect, odeur, goût, absence de bombage des boîtes, pH acide ($< 4{,}5$ pour les fruits), taux de sucre au réfractomètre.
  \item Enjeux : réduction des pertes, disponibilité hors saison, création de revenus (coopératives, PME agroalimentaires), valorisation des productions locales (mangues de l'Adamaoua, tomates de l'Ouest).
\end{itemize}
\end{aretenir}

\begin{attention}[Pièges fréquents au Bac]
\begin{itemize}
  \item Ne pas confondre \textbf{pasteurisation} et \textbf{stérilisation} : la pasteurisation ne détruit pas les spores ; seule la stérilisation (115--121\,$^\circ$C) rend le produit stable à température ambiante.
  \item Le sucre ne « tue » pas directement les microbes : il \textbf{réduit l'eau libre} ($a_w$), ce qui bloque leur développement. Rédige ce mécanisme, ne te contente pas de « le sucre conserve ».
  \item Le bombage d'une boîte de conserve signale une fermentation (production de gaz) : le produit est \textbf{impropre à la consommation}, même si l'odeur semble normale.
  \item Dans un ordinogramme, respecte l'\textbf{ordre chronologique} des opérations : le conditionnement précède toujours la stérilisation dans l'appertisation.
\end{itemize}
\end{attention}

\begin{aretenir}[Checklist avant le Bac]
\begin{itemize}
  \item $\square$ Je sais définir transformation, conservation, pasteurisation, stérilisation et appertisation.
  \item $\square$ Je cite les trois grandes causes d'altération des fruits (microbienne, enzymatique, physico-chimique).
  \item $\square$ J'explique le rôle de l'eau libre et comment le sucre ou le séchage la réduisent.
  \item $\square$ Je distingue pasteurisation (70--90\,$^\circ$C) et stérilisation (115--121\,$^\circ$C) avec leurs objectifs.
  \item $\square$ Je décris, dans l'ordre, les étapes de fabrication du jus de mangue et de la confiture de mangue.
  \item $\square$ Je décris, dans l'ordre, les étapes de fabrication du concentré de tomate.
  \item $\square$ Je justifie le remplissage à chaud et la fermeture hermétique (vide partiel, exclusion de l'air).
  \item $\square$ Je sais pourquoi le blanchiment/ébouillantage inactive les enzymes responsables du brunissement.
  \item $\square$ Je cite au moins deux paramètres de contrôle qualité (pH, taux de sucre, absence de bombage).
  \item $\square$ Je relie la conservation des fruits aux enjeux camerounais : pertes post-récolte, sécurité alimentaire, revenus ruraux.
  \item $\square$ Je sais schématiser une chaîne de transformation sous forme d'ordinogramme fléché.
  \item $\square$ Je rédige mes réponses avec le vocabulaire exact (pulpage, concentration, conditionnement, barrières).
\end{itemize}
\end{aretenir}

\findecours

\end{document}
